
%%FILE blueprint/src/chapter/CompactOperators.tex
\chapter{Compact Operators and Fredholm Theory}
\label{chap:compact}

Serre's \emph{Endomorphismes compl\`etement continus des espaces de Banach $p$-adiques}
\cite{Serre} attaches to a completely continuous operator $u$ on a $p$-adic Banach space a
characteristic power series $\det(1 - Tu)$, entire in $T$, whose reciprocal roots are exactly the
eigenvalues of $u$, and which controls a Riesz decomposition of the space.  That is the machine
we need: a Hecke operator will be completely continuous, its determinant will be computable from
a matrix, and the Newton polygon of the determinant will read off the eigenvalue valuations.

Serre works over a complete nonarchimedean field.  Buzzard \cite{Buzzard} works over a Banach
algebra over such a field, and Johansson--Newton \cite{JN} over a Banach--Tate ring with no
ground field at all.  We formalise the last of these, since the applications need the ground
ring to be an affinoid rather than a field, and since nothing in the arguments actually uses a
field.  The specialisations are then corollaries.

\section{The setting}

\begin{definition}
  \label{def:istate}
  \lean{TateFredholm.IsTate}
  \leanok
  A \emph{Banach--Tate ring} is a complete normed ring $R$ possessing a
  \emph{pseudo-uniformizer}: a unit $\varpi$ with $\lVert \varpi \rVert < 1$ whose norm is
  multiplicative.  We write $\mathrm{IsTate}\, R$.
\end{definition}

Note that the norm on $R$ is only submultiplicative; the pseudo-uniformizer is the one element
whose norm behaves multiplicatively, and that turns out to be all the analysis needs.

\begin{definition}
  \label{def:cspace}
  \lean{TateFredholm.cSpace}
  \leanok
  For an index type $I$ the \emph{model space} $c(I, R)$ is the space of families $(x_i)_{i \in
  I}$ in $R$ tending to $0$ along the cofinite filter, with the supremum norm.
\end{definition}

Every module we care about is either $c(I,R)$ or a direct summand of one:

\begin{definition}
  \label{def:onable}
  \uses{def:cspace}
  \lean{TateFredholm.IsONable, TateFredholm.HasPr}
  \leanok
  A Banach $R$-module $M$ is \emph{orthonormalisable} if it is isometrically isomorphic to some
  $c(I,R)$, and has \emph{property (Pr)} if it is a direct summand of an orthonormalisable
  module.
\end{definition}

Property (Pr) is Buzzard's hypothesis, and it is what a space of modular forms at a non-neat
level satisfies when orthonormalisability itself fails.

\begin{proposition}
  \label{prop:tate-algebra-onable}
  \uses{def:onable, def:tate-algebra}
  \lean{TateFredholm.restrictedEquivCSpace, TateFredholm.isONable_restricted}
  \leanok
  The Tate algebra $K\langle X \rangle$ is orthonormalisable, with the monomials $X^n$ as
  orthonormal basis; the isometry $K\langle X\rangle \cong c(\N, K)$ sends a series to its
  coefficient family.
\end{proposition}

This proposition is the reason the weight-$\kappa$ action of Chapter~\ref{chap:qmf}, which acts
on a Tate algebra, is an operator on a model space and so has matrix coefficients.

\section{Compactness}

\begin{definition}
  \label{def:matrixcoeff}
  \uses{def:cspace}
  \lean{TateFredholm.matrixCoeff, TateFredholm.rowNorm}
  \leanok
  For $u : c(I,R) \to c(J,R)$ continuous linear, the \emph{matrix coefficient} $u_{j i}$ is the
  $j$-th coordinate of $u(e_i)$, where $e_i$ is the $i$-th basis vector.  The $j$-th
  \emph{row norm} is $\sup_i \lVert u_{ji}\rVert$.
\end{definition}

\begin{definition}
  \label{def:iscompactoid}
  \uses{def:matrixcoeff}
  \lean{TateFredholm.IsCompletelyContinuous, TateFredholm.IsCompactoid}
  \leanok
  $u$ is \emph{completely continuous} if it is an operator-norm limit of finite-rank operators,
  and \emph{compactoid} if its row norms tend to $0$ along the cofinite filter.
\end{definition}

\begin{proposition}
  \label{prop:compactoid-criterion}
  \uses{def:iscompactoid}
  \lean{TateFredholm.isCompactoid_of_row_decay'}
  \leanok
  On $c(I,R)$ the two notions agree, and a row-decay estimate $\lVert u_{ji} \rVert \leq C
  \sigma^{j}$ with $\sigma < 1$ suffices for both.
\end{proposition}

The row-decay criterion is the only form of compactness we ever verify in practice: for a Hecke
operator the estimate falls out of the weight action, see Theorem~\ref{thm:qmf-compact}.

\section{The Fredholm determinant}

\begin{definition}
  \label{def:minor}
  \uses{def:matrixcoeff}
  \lean{TateFredholm.minor, TateFredholm.charCoeff, TateFredholm.charPowerSeries}
  \leanok
  For a finite $S \subseteq I$ let $\mathrm{minor}(u, S)$ be the determinant of the $S \times S$
  submatrix of $u$.  For compactoid $u$ the family of $n$-element minors is summable, and
  \[
    c_n(u) = (-1)^n \sum_{\lvert S \rvert = n} \mathrm{minor}(u, S), \qquad
    \det(1 - Tu) = \sum_{n \geq 0} c_n(u) T^n .
  \]
\end{definition}

\begin{lemma}[Ultrametric Hadamard bound]
  \label{lem:hadamard}
  \uses{def:minor}
  \lean{TateFredholm.norm_det_le_pow_of_row_bound}
  \leanok
  If every entry of a square matrix satisfies $\lVert A_{ji}\rVert \leq \sigma^{w_j}$ for a row
  weight $w$, then $\lVert \det A \rVert \leq \sigma^{\sum_j w_j}$.
\end{lemma}

\begin{proof}
  \leanok
  Each Leibniz monomial picks one entry from every row, so has norm at most $\prod_j
  \sigma^{w_j}$; the ultrametric inequality bounds the sum of the monomials by their maximum.
\end{proof}

Summability of the minors comes from the same bound: a minor of norm $\geq \varepsilon$ forces
its index set inside the finite set of ``large'' rows, and there are finitely many $n$-subsets of
a finite set.

\begin{theorem}
  \label{thm:charpowerseries-entire}
  \uses{def:minor}
  \lean{TateFredholm.charPowerSeries_isEntire, TateFredholm.charCoeff_zero}
  \leanok
  For compactoid $u$ the series $\det(1 - Tu)$ is entire, with constant term $1$.
\end{theorem}

\begin{theorem}[Invariances]
  \label{thm:charpowerseries-inv}
  \uses{def:minor}
  \lean{TateFredholm.charPowerSeries_comm, TateFredholm.charPowerSeries_conj}
  \leanok
  $\det(1 - Tuv) = \det(1 - Tvu)$, and the determinant is invariant under conjugation by an
  isomorphism.  Consequently it may be computed in any model of the space.
\end{theorem}

\begin{proposition}[Block operators]
  \label{prop:blockop}
  \uses{def:minor}
  \lean{TateFredholm.blockOp, TateFredholm.isCompactoid_blockOp}
  \leanok
  An operator on $c(\sigma \times I, R)$ with $\sigma$ finite is assembled from its $\sigma
  \times \sigma$ blocks; it is compactoid iff each block is, and its determinant is computed
  blockwise.
\end{proposition}

This is the shape a Hecke operator takes: the class set indexes $\sigma$, and each block is the
weight action of a single certificate element.

\section{Serre's theorems}

\begin{theorem}[{\cite[Prop.\ 11]{Serre}}]
  \label{thm:serre11}
  \uses{thm:charpowerseries-entire}
  \lean{TateFredholm.isUnit_one_sub_smul_iff_isUnit_evalT}
  \leanok
  For compactoid $u$ and $a \in R$, the operator $1 - a u$ is invertible if and only if
  $\det(1-Tu)$ evaluated at $a$ is a unit.
\end{theorem}

\begin{theorem}[{\cite[Prop.\ 12]{Serre}}]
  \label{thm:serre12}
  \uses{thm:serre11}
  \lean{TateFredholm.exists_eigenvector_of_evalT_charPowerSeries_eq_zero}
  \leanok
  Over a field, if $a \neq 0$ is a zero of $\det(1 - Tu)$ then $a^{-1}$ is an eigenvalue of $u$;
  and conversely.
\end{theorem}

Theorem~\ref{thm:serre12} is the statement that makes the Newton polygon of the determinant say
something about the operator, and hence the hinge of the whole thesis.

\begin{theorem}[Riesz decomposition]
  \label{thm:riesz}
  \uses{thm:serre12}
  \lean{TateFredholm.exists_riesz_decomposition}
  \leanok
  Let $u$ be compactoid on $c(I,K)$ over a discretely valued complete field, and let $a$ be a
  zero of $\det(1 - Tu)$ of multiplicity $h$.  Then $c(I,K) = N \oplus F$ with both summands
  $u$-stable, $(1 - au)^h$ vanishing on $N$, $(1-au)$ invertible on $F$, and
  \[
    \dim_K N = h .
  \]
\end{theorem}

The subspace $N$ is the \emph{finite-slope subspace} attached to $a$; it is finite dimensional,
which is what makes eigenvalue statements about it meaningful, and everything commuting with $u$
preserves it, which is what makes a Hecke algebra act on it (Proposition~\ref{prop:qmf-eigensystem}).

%% The projector p = e^h and Serre's Lemme 2 (block-triangular determinants) are the internal
%% steps: TateFredholm.exists_rieszProjection, charPowerSeries_blockTriangular.


%%FILE blueprint/src/chapter/Jacobs.tex
\chapter{The Slopes of $U_3$}
\label{chap:jacobs}

We now instantiate the general theory at the case worked out by Jacobs \cite{Jacobs}, and prove
the theorem the thesis is aimed at: for the definite quaternion algebra over $\Q$ ramified at
$2$ and $\infty$, at level $U_1(9)$ and at a suitable weight, the eigenvalues of $U_3$ have
$3$-adic valuations
\[
  \tfrac12, \quad \tfrac32, \quad \tfrac52, \quad \dots
\]
--- an arithmetic progression of half-integers.  The half-integrality is not an artefact: it is
what makes the example interesting, and it is possible only because the coefficient field is
ramified over $\Q_3$.

Everything general has been done.  What remains is to exhibit the data --- the algebra, the
place, the level, the weight, the coset representatives --- and to compute one matrix.

\section{The data}

\begin{definition}
  \label{def:jacobs-setting}
  \lean{JacobsSlash.D, JacobsSlash.v₃, JacobsSlash.K₃, JacobsSlash.theta}
  \leanok
  Let $D = \bbH_{\Q}$ be Hamilton's quaternions over $\Q$, a definite quaternion algebra of
  discriminant $2$; let $v_3$ be the place above $3$ and $K_3 = \Q_3$ its completion.  Since $3
  \nmid 2$, $D$ splits at $v_3$, and a splitting is fixed by choosing a square root $\nu$ of
  $-2$ in $K_3$; the resulting isomorphism $\theta_3 : D \otimes_{\Q} K_3 \cong M_2(K_3)$ is the
  rigidification of Definition~\ref{def:rigidification}.
\end{definition}

%% ν₃ is located explicitly by Hensel at 2695 mod 3^10; the location is what makes the matrix of
%% U₃ computable rather than merely existent.

\begin{definition}
  \label{def:hurwitz}
  \lean{JacobsSlash.hurwitzOrder}
  \leanok
  The \emph{Hurwitz order} is $\calO = \Z\langle i, j, \tfrac12(1+i+j+k)\rangle \subseteq \bbH_{\Q}$,
  the unique maximal order up to conjugacy.  We record membership by the parity condition that
  $2x$ has integer coordinates all of the same parity.
\end{definition}

\begin{proposition}
  \label{prop:hurwitz-units}
  \uses{def:hurwitz}
  \lean{JacobsSlash.isUnit_iff_hnorm_eq_one, JacobsSlash.card_units_hurwitzOrder}
  \leanok
  A Hurwitz quaternion is a unit exactly when its reduced norm is $1$, and $\calO^\times$ has
  exactly $24$ elements.
\end{proposition}

\begin{definition}
  \label{def:jacobs-levels}
  \uses{def:hurwitz, def:sigma0}
  \lean{JacobsSlash.U0, JacobsSlash.U1_9}
  \leanok
  $U_0(1) \subseteq D_f^\times$ consists of the adelic units integral at every finite place ---
  the unit group of $\calO \otimes \Zhat$ --- and $U_1(9)$ refines it by the congruence
  $\theta_3(g) \equiv \left(\begin{smallmatrix} * & * \\ 0 & 1\end{smallmatrix}\right)$ modulo
  $9$.
\end{definition}

\begin{proposition}
  \label{prop:jacobs-levels-compact-open}
  \uses{def:jacobs-levels, prop:compact-open-levels}
  \lean{JacobsSlash.isCompact_U0, JacobsSlash.isCompact_U1_9, JacobsSlash.isOpen_U1_9,
        JacobsSlash.finite_image_doubleCoset_U1_9}
  \leanok
  $U_0(1)$ and $U_1(9)$ are compact open, so Theorem~\ref{thm:hecke-finiteness} applies and the
  Hecke operators at this level are defined without exhibiting representatives.
\end{proposition}

\begin{proof}
  \leanok
  \uses{prop:compact-open-levels}
  The everywhere-integral elements of $D \otimes \A_f$ are exactly the $\Zhat$-span of the
  Hurwitz basis $1, i, j, \omega$: one inclusion is clear, and the other extracts each coordinate
  by a $\Q$-linear functional, extends it to the completion, and uses that a rational integral at
  every finite place is a rational integer.  That set is compact by
  Proposition~\ref{prop:compact-open-levels} and open because the integral adeles are open, so
  $U_0(1)$ is compact open; $U_1(9)$ is cut out of it by a congruence condition at $3$, which is
  clopen because closed balls in an ultrametric field are open.
\end{proof}

\section{Class number one}

\begin{theorem}[{[Jacobs, Theorem 2.1]}]
  \label{thm:jacobs-classset}
  \uses{def:jacobs-levels, thm:fujisaki}
  \lean{JacobsSlash.hClassNumberOne, JacobsSlash.classRep_bijective}
  \leanok
  $D_f^\times = D^\times \cdot U_0(1)$, and the class set at level $U_1(9)$ has exactly three
  elements, represented by explicit adelic $c_0, c_1, c_2$.
\end{theorem}

\begin{proof}
  \leanok
  Class number one for the Hurwitz order is classical; formalised here by the idelic dictionary,
  reducing a global integrality statement to local approximation.  Passing from $U_0(1)$ to
  $U_1(9)$ replaces the class set by the orbits of $\calO^\times$ on the primitive vectors of
  $(\Z/9)^2$, which are computed to be three.
\end{proof}

\begin{lemma}[{[Jacobs, Lemma 2.2]}]
  \label{lem:jacobs-neat}
  \uses{thm:jacobs-classset, thm:evalatreps}
  \lean{JacobsSlash.stabilizerAt_classRep}
  \leanok
  The stabilisers $\Gamma_i = c_i^{-1}D^\times c_i \cap U_1(9)$ are trivial.
\end{lemma}

So the level is neat, and Theorem~\ref{thm:evalatreps} gives $S_\kappa(U_1(9)) \cong c(\mathrm{Fin}\,3
\times \N, K_3)$: a form is exactly a triple of elements of the Tate algebra.

\section{The weight}

\begin{definition}
  \label{def:jacobs-weight}
  \uses{def:analyticweight}
  \lean{JacobsSlash.jacobsCharOf, JacobsSlash.jacobsWeight}
  \leanok
  For $t \in K_3$ with $\lVert t \rVert < 1$ let $\kappa_t(u) = u^t := \exp(t \log u)$ on the
  $1$-units.  The expansion of $\kappa_t(cz+d)$ is the binomial series, whose coefficients
  satisfy the decay required by Definition~\ref{def:weightseries}; the pair is an analytic weight
  in the sense of Definition~\ref{def:analyticweight}.
\end{definition}

%% The hypothesis ‖t‖ < 1 is Jacobs' t ∈ m₃ (p. 38), not an artefact of the formalisation.

\section{The operator and its matrix}

\begin{definition}
  \label{def:u3}
  \uses{def:jacobs-weight, def:heckeoperator, def:jacobs-levels}
  \lean{JacobsSlash.eta3, JacobsSlash.heckeU3}
  \leanok
  $U_3 = [U_1(9)\, \eta_3\, U_1(9)]$ where $\eta_3$ is the adelic element which is $1$ away from
  $3$ and $\left(\begin{smallmatrix}1&0\\0&3\end{smallmatrix}\right)$ at $3$.
\end{definition}

\begin{theorem}
  \label{thm:u3-matrix}
  \uses{def:u3, thm:transport, lem:jacobs-neat}
  \lean{JacobsSlash.heckeU3_apply_classRep, JacobsSlash.matrixCoeff_blockEntry}
  \leanok
  With respect to the three representatives, $(U_3\varphi)(c_i) = \sum_j \varepsilon_{ij}
  (\varphi(c_j))$, and the matrix coefficients of the blocks $\varepsilon_{ij}$ are the explicit
  coefficients of the generating functions of Proposition~\ref{prop:kappaslash-matrix} evaluated
  at the certificate elements.
\end{theorem}

The certificate data --- four right coset representatives $v_t$ of $U_1(9)\eta_3 U_1(9)$, the
permutation $\sigma$, and the factorisations $c_i v_t^{-1} = d\, c_{\sigma(i,t)}\, u$ --- is
computed by hand; it is the only part of the argument that is a computation rather than a
theorem, and it is what the general framework was built to consume.

\begin{definition}
  \label{def:u3-det}
  \uses{thm:u3-matrix, def:heckecharpowerseries}
  \lean{JacobsSlash.charPowerSeriesU3}
  \leanok
  $\det(1 - T\,U_3)$, computed in the block model through the three evaluations.
\end{definition}

\section{The eigenvalues}

Over $K_3$ the matrix of $U_3$ does not have unit top-left minors, and no arithmetic progression
of integer slopes can appear: the answer involves half-integers, so the polygon must be read over
a ramified extension.  Adjoining a primitive cube root of unity $\omega$ gives $K_3(\omega) =
\Q_3(\sqrt{-3})$, ramified of degree $2$, in which $v_3$ takes half-integer values; there the
matrix splits into eigenblocks.

\begin{theorem}
  \label{thm:u3-factorisation}
  \uses{def:u3-det, prop:qmf-basechange}
  \lean{JacobsSlash.charPowerSeriesU3_factorisation}
  \leanok
  Over an isometric extension containing $\omega$, $\det(1 - TU_3)$ factors according to the
  $\omega$-eigenspaces of the block structure, and on the relevant factor the rescaled matrix has
  every top-left minor of norm $1$.
\end{theorem}

\begin{theorem}[The crux]
  \label{thm:u3-eigenvalues}
  \uses{thm:u3-factorisation, thm:unit-minors, cor:slope-reading}
  \lean{JacobsSlash.exists_eigenvalue_U3_halfIntegral}
  \leanok
  For every $j \in \N$ there is an eigenvalue $a$ of $U_3$ with
  \[
    v_3(a) = j + \tfrac12 .
  \]
\end{theorem}

\begin{proof}
  \leanok
  \uses{thm:unit-minors, cor:slope-reading, thm:np-slope-reading}
  Theorem~\ref{thm:u3-factorisation} puts the relevant factor in the shape of
  Theorem~\ref{thm:unit-minors} with $\varpi = \sqrt{-3}$, whose valuation is $\tfrac12$ in the
  normalisation $v_3(3) = 1$.  The exact coefficient valuations $\binom m2$ in that normalisation
  become $\tfrac12\binom m2$ here, whose successive differences are $\tfrac12, \tfrac32, \dots$;
  Corollary~\ref{cor:slope-reading} reads these off the polygon and
  Theorem~\ref{thm:serre12} converts slopes into eigenvalues.
\end{proof}

\begin{theorem}[The crux, for forms]
  \label{thm:u3-eigenforms}
  \uses{thm:u3-eigenvalues, thm:qmf-eigenform-criterion, prop:qmf-basechange}
  \lean{JacobsSlash.exists_eigenform_U3_halfIntegral}
  \leanok
  For every $j \in \N$ there exist $a \in \bbC_3$ and a nonzero overconvergent quaternionic
  modular form $\varphi$ of weight $\kappa_t$ and level $U_1(9)$ over $\bbC_3$ with
  \[
    U_3 \varphi = a\varphi, \qquad v_3(a) = j + \tfrac12 .
  \]
\end{theorem}

\begin{proof}
  \leanok
  \uses{thm:qmf-eigenform-criterion, thm:evalatreps}
  Theorem~\ref{thm:u3-eigenvalues} produces an eigenvector in the block model; the level is neat
  by Lemma~\ref{lem:jacobs-neat}, so Theorem~\ref{thm:evalatreps} transports it back to a form,
  and the transport intertwines the block operator with $U_3$ by Theorem~\ref{thm:transport}.
  Base change along $K_3(\omega) \hookrightarrow \bbC_3$ is
  Proposition~\ref{prop:qmf-basechange}.
\end{proof}

This is the statement announced in the introduction: the slopes of a compact Hecke operator, for
this weight and level, form an arithmetic progression --- and one that is visibly not integral,
so that the fibre of the ``quaternionic eigenvariety'' over this weight is not what a naive guess
would predict.


%%FILE blueprint/src/chapter/LWXHalo.tex
\chapter{The Halo Estimate of Liu--Wan--Xiao}
\label{chap:lwx-halo}

Chapters~\ref{chap:qmf}--\ref{chap:jacobs} computed slopes at one weight at a time.  Weight space
$\mathcal{W}$ is a union of open discs, and the picture above it changes character near the
boundary: Liu, Wan and Xiao \cite{LWX} proved, for the spectral curve of overconvergent
automorphic forms on a definite quaternion algebra, that over the boundary annulus it is a
disjoint union of ``pieces'' finite and flat over the annulus, on each of which the slope of
$a_p$ is a fixed multiple of $v(T)$ --- the multiple being independent of the weight.  This is
the \emph{spectral halo}, an analogue for $\mathrm{Spc}_D$ of the Coleman--Mazur--Buzzard--Kilford
conjecture; the precise statements, and exactly which parts of them are formalised here, are
laid out in Chapter~\ref{chap:lwx-slopes} and \S\ref{sec:what-15-says}.

Their argument has three layers, and this chapter is the first: a purely integral estimate.
Let $\Lhalo = \Zp\llbracket T, pT^{-1}\rrbracket$ be the ring of \cite[Lemma 3.15]{LWX}.  On the
integral model of overconvergent forms the operator $U_p$ has a matrix $P$ with entries in
$\Lhalo$, and \cite[Theorem 3.16]{LWX} says that its Fredholm determinant
$\Char(P) = \sum_n c_n X^n$ exists and satisfies
\[
  c_n \in T^{\lambda(n)}\,\Lhalo, \qquad
  \lambda(n) = \sum_{k<n}\Bigl(\Bigl\lfloor \tfrac{k}{t}\Bigr\rfloor -
  \Bigl\lfloor \tfrac{k}{pt}\Bigr\rfloor\Bigr),
\]
where $t$ is the size of the class set.  Specialising $T$ at a point $T_0$ of the boundary
annulus $p^{-1} < \lVert T_0\rVert < 1$ gives \cite[Corollary 3.18]{LWX}: the Newton polygon of
the specialised series lies on or above an explicit polygon whose slopes are the numbers
$(\lfloor k/t\rfloor - \lfloor k/pt\rfloor)\,v(T_0)$.  The remarkable feature, and the whole
point, is that this lower bound is a \emph{multiple of $v(T_0)$}: it degenerates to the trivial
bound as one moves into the interior, and it is sharp near the boundary.

Everything in this chapter is unconditional and, except where stated, is proved for an arbitrary
combinatorial datum of the shape that \cite[Proposition 3.1]{LWX} produces; the identification of
that datum with the genuine Hecke operator is Chapter~\ref{chap:lwx-model}, and the slope
consequences are Chapter~\ref{chap:lwx-slopes}.  Throughout $p$ is an odd prime, so that
$q = p$ in \cite{LWX}'s notation; $p = 2$ is out of scope.

\section{The integral halo ring}

\begin{definition}
  \label{def:haloint}
  \lean{LWX.HaloInt, LWX.HaloInt.norm_def, LWX.HaloInt.norm_le_one}
  \leanok
  $\Lhalo$ is realised coefficientwise, as \cite[proof of Corollary 3.18]{LWX} describes its
  elements: a two-sided stream $d : \Z \to \Zp$ with $v(d_j) \geq \max(0, -j)$, that is
  $\lVert d_j\rVert \leq p^{\min(0,j)}$, with pointwise addition, convolution product, and the
  gauge norm
  \[
    \lVert d\rVert = \sup_{j\in\Z} \lVert d_j \rVert\, p^{-j} \in [0,1] .
  \]
\end{definition}

Both tails of the convolution die --- the left-hand one by the halo bound, the right-hand one by
integrality --- so the product is a convergent sum over $\Z$ and the ring is complete.

\begin{proposition}
  \label{prop:haloint-norm}
  \uses{def:haloint}
  \lean{LWX.HaloInt.T, LWX.HaloInt.norm_T_mul, LWX.HaloInt.constRingHom,
        LWX.HaloInt.norm_const, LWX.HaloInt.exists_T_pow_mul_of_norm_le,
        LWX.HaloInt.norm_le_of_eq_T_pow_mul}
  \leanok
  The constant embedding $\Zp \to \Lhalo$ is isometric; the element $T = \delta_1$ acts as the
  coefficient shift and satisfies $\lVert Tx\rVert = p^{-1}\lVert x\rVert$ exactly.
  Consequently
  \[
    x \in T^{k}\,\Lhalo \iff \lVert x\rVert \leq p^{-k},
  \]
  which is \cite[Lemma 3.15]{LWX} in the only form the estimate consumes.
\end{proposition}

Divisibility by a power of $T$ has become a norm inequality; this is what makes the halo estimate
an instance of the general Hadamard bound of Chapter~\ref{chap:compact} rather than an argument
about ideals.

\begin{definition}[Specialisation at a halo point]
  \label{def:halo-specialize}
  \uses{def:haloint}
  \lean{LWX.HaloInt.specialize, LWX.HaloInt.specializeHom, LWX.HaloInt.continuous_specialize,
        LWX.HaloInt.norm_specialize_le}
  \leanok
  Let $K$ be a complete ultrametric field, $\psi : \Zp \to K$ an isometric ring homomorphism and
  $T_0 \in K$ with $p^{-1} < \lVert T_0\rVert < 1$.  Then $\sum_{j\in\Z}\psi(d_j)T_0^{\,j}$
  converges, and $d \mapsto d(T_0)$ is a continuous ring homomorphism $\Lhalo \to K$
  satisfying $\lVert d(T_0)\rVert \leq \lVert T_0\rVert^{k}$ whenever $\lVert d\rVert\leq p^{-k}$.
\end{definition}

The two hypotheses on $\lVert T_0\rVert$ are exactly what makes the two tails converge; this is
the annulus $\mathcal{W}^{>1/p}$ of \cite{LWX}.  Continuity is needed because $\Char(P)$ is a
limit of determinants of minors, and the specialisation has to be pushed through that limit
(\S\ref{sec:seam}).

\begin{proposition}[{The Banach--Tate ring $A = \Lhalo[1/T]$}]
  \label{prop:halotate}
  \uses{def:haloint, def:istate}
  \lean{LWX.HaloTate, LWX.HaloTate.ofIntRingHom, LWX.HaloTate.T,
        LWX.HaloTate.norm_T_mul, LWX.HaloTate.pseudoUniformizer, LWX.HaloTate.norm_le_one_iff,
        LWX.HaloTate.exists_T_zpow_mul_ofInt}
  \leanok
  Inverting $T$ gives the ring $A$ of streams satisfying a \emph{shifted} halo bound.  It is
  complete for the gauge norm, contains $\Lhalo$ isometrically as its unit ball, and $T$ is a
  unit with $\lVert Tx\rVert = p^{-1}\lVert x\rVert$: a multiplicative pseudo-uniformizer.  So
  $A$ is a Banach--Tate ring in the sense of Definition~\ref{def:istate}, and every element is
  $T^{-k}$ times an integral one.
\end{proposition}

No Noetherian hypothesis is claimed, and none is available: this is precisely why
Chapter~\ref{chap:riesz-coleman} was developed Noetherian-free.

\section{The universal character}

\begin{definition}
  \label{def:univchar}
  \uses{def:haloint}
  \lean{LWX.teichmuller, LWX.oneUnitPart, LWX.qlog, LWX.qlog_mul, LWX.oneAddTPow,
        LWX.univChar, LWX.univChar_mul, LWX.norm_univChar_le}
  \leanok
  Write $\Zp^{\times} = \Delta \times (1+p\Zp)$, $x = \omega_{\mathrm{T}}(x)\langle x\rangle$,
  with the Teichm\"uller lift $\omega_{\mathrm{T}}(x) = \lim_n x^{p^n}$; let
  $\ell\langle x\rangle = \tfrac1p\log\langle x\rangle \in \Zp$.  For a nebentypus
  $\omega : (\Z/p)^{\times}\to\Zp^{\times}$ the \emph{universal character} of
  \cite[Notation 2.1]{LWX} is
  \[
    [x] \;=\; \omega(\bar x)\,(1+T)^{\ell\langle x\rangle} \;\in\; (\Lhalo)^{\times},
    \qquad (1+T)^{s} = \sum_{r\geq 0}\binom{s}{r}T^{r} .
  \]
  It is multiplicative and takes values in the unit ball.
\end{definition}

Mathlib's \texttt{WittVector.teichmuller} is a different map, so the lift, the logarithm on
$1+p\Zp$ and its additivity are all built here; additivity of $\log$ is what allows
$\log((cz+d)/d_0)/p$ to be split into a constant plus a series, and that split is what makes the
binomial estimates of the next section applicable.

\section{Tilted degree}

\cite{LWX} call \S3.3--\S3.13 ``a mild $p$-adic analysis computation (which is the core of our
paper)''.  Its object is a notion of degree for functions $\Zp\to\Zp$ that is stable under the
operations the matrix entries are built from.

\begin{definition}
  \label{def:tilteddeg}
  \lean{LWX.TiltedDeg}
  \leanok
  A function $f : \Zp\to\Zp$ has \emph{tilted degree $\leq n$} if
  $\lVert \tilde\Delta^{m}f(z)\rVert \leq p^{\,n-m}$ for all $m$ and all $z$, where
  $\tilde\Delta$ is the forward difference $f(z+1)-f(z)$.
\end{definition}

This is clause (1) of \cite[Definition-Proposition 3.8]{LWX}.  Taking the bound
\emph{everywhere} rather than only at $z=0$ costs nothing and buys two things: no Mahler
expansion is needed to state it, and the class is closed under pointwise limits.  Clause (2),
the Mahler-coefficient form, is then never needed and is not stated.

\begin{lemma}[The calculus]
  \label{lem:tilted-calculus}
  \uses{def:tilteddeg}
  \lean{LWX.TiltedDeg.mul, LWX.TiltedDeg.of_tendsto, LWX.TiltedDeg.mono, LWX.TiltedDeg.sum,
        LWX.tiltedDeg_const, LWX.fwdDiff_iter_mul, LWX.fwdDiff_iter_choose}
  \leanok
  Tilted degree is monotone in $n$, additive under products, closed under convergent sums and
  pointwise limits, and constants have tilted degree $0$.  The engine is the Leibniz rule
  $\tilde\Delta^{m}(fg) = \sum_{i+j=m}\tilde\Delta^{i}f \cdot (\tilde\Delta^{j}g)(\cdot + i)$,
  which is \cite[(3.5.1)]{LWX}, together with
  $\tilde\Delta\binom{z}{n} = \binom{z}{n-1}$, which is \cite[(3.5.2)]{LWX}.
\end{lemma}

The Chu--Vandermonde input \cite[(3.5.3)]{LWX} is mathlib's binomial-ring
\texttt{add\_choose\_eq}, and the Mahler coefficient $a_m(f) = \tilde\Delta^{m}f(0)$ is mathlib's
\texttt{fwdDiff\_iter\_eq\_sum\_shift}.

\begin{lemma}[{\cite[Lemmas 3.12 and 3.13]{LWX}}]
  \label{lem:tilted-binomial}
  \uses{def:tilteddeg, lem:tilted-calculus}
  \lean{LWX.tiltedDeg_choose_of_pMul, LWX.tiltedDeg_choose_monomial_p,
        LWX.tiltedDeg_choose_of_logShape}
  \leanok
  If $f : \Zp\to p\Zp$ then $z\mapsto\binom{f(z)}{n}$ has tilted degree $\leq \lfloor n/p\rfloor$;
  and if $g(z) = \sum_k p^{k-1}(a_k/k)z^{k}$ has the shape of a normalised logarithm then
  $z\mapsto\binom{g(z)}{r}$ has tilted degree $\leq 0$.
\end{lemma}

These two estimates are where the odd prime is used: the proof rests on the Legendre margin
$2\,v_p(n!)\leq n-1$ and on $4v\leq p^{v}+3$, both of which fail at $p=2$.

\section{The matrix of $U_p$}

\begin{definition}
  \label{def:localmat}
  \lean{LWX.LocalMat, LWX.LocalMat.IsUpShape}
  \leanok
  A \emph{local matrix} is $\delta = \begin{pmatrix} a & b\\ c& d\end{pmatrix}\in M_2(\Zp)$ with
  $p \mid c$ and $d \in \Zp^{\times}$: the monoid $M_1$ of \cite[(2.3.3)]{LWX}.  It has
  \emph{$U_p$-shape} if in addition $a \in p\Zp$, which is clause (3) of
  \cite[Proposition 3.1]{LWX}.
\end{definition}

\begin{definition}[The entry stream, {\cite[Propositions 3.4 and 3.14]{LWX}}]
  \label{def:entry}
  \uses{def:localmat, def:univchar, def:haloint}
  \lean{LWX.entry, LWX.entryCoeff, LWX.coeff_entry}
  \leanok
  \cite[Proposition 3.4]{LWX} computes the matrix of a single $\Vert_\delta$ in the Mahler basis:
  \[
    P_{m,n}(\delta) \;=\; \tilde\Delta^{m}\Bigl(\;\binom{\tfrac{az+b}{cz+d}}{n}\,[\,cz+d\,]\;
    \Bigr)\Big|_{z=0} .
  \]
  Expanding $[cz+d] = [d_0](1+T)^{g(z)}$ with $g(z) = \tfrac1p\log\bigl((cz+d)/d_0\bigr)$ gives
  the $T$-coefficient stream
  $r \mapsto [d_0]\,\tilde\Delta^{m}\bigl(\binom{f(z)}{n}\binom{g(z)}{r}\bigr)\big|_{z=0}$,
  and we take that stream as the definition of the entry $P_{m,n}(\delta)\in\Lhalo$.
\end{definition}

That the stream computes the action on the integral model is a separate statement, proved in
Chapter~\ref{chap:lwx-model} (Proposition~\ref{prop:prop34}); here it is simply the matrix we
estimate.

\begin{proposition}[{\cite[Proposition 3.14(1)]{LWX}}]
  \label{prop:entry-bound}
  \uses{def:entry, lem:tilted-binomial, prop:haloint-norm}
  \lean{LWX.norm_entry_le, LWX.norm_entryCoeff_le, LWX.norm_entry_le_M1}
  \leanok
  For a local matrix of $U_p$-shape,
  \[
    \bigl\lVert P_{m,n}(\delta)\bigr\rVert \;\leq\; p^{-(m - \lfloor n/p\rfloor)},
    \qquad\text{i.e.}\qquad
    P_{m,n}(\delta) \in \mathfrak{m}^{\max(m-\lfloor n/p\rfloor,\,0)} .
  \]
  Without the shape hypothesis, that is for a general element of $M_1$, one still has
  $\lVert P_{m,n}(\delta)\rVert \leq p^{-(m-n)}$.
\end{proposition}

\begin{proof}
  \leanok
  \uses{lem:tilted-binomial, lem:tilted-calculus, def:univchar}
  The $U_p$-shape makes $f(z) = (az+b)/(cz+d)$ take values in $p\Zp$, so
  $\binom{f(z)}{n}$ has tilted degree $\leq\lfloor n/p\rfloor$ by the first half of
  Lemma~\ref{lem:tilted-binomial}; $g$ is a normalised logarithm, so $\binom{g(z)}{r}$ has tilted
  degree $\leq 0$ by the second half; the product has tilted degree
  $\leq \lfloor n/p\rfloor$ by Lemma~\ref{lem:tilted-calculus}, and evaluating the definition of
  tilted degree at the $m$-th difference and $z=0$ is the claim.  The constant $[d_0]$ is a unit.
\end{proof}

\begin{definition}
  \label{def:updatum}
  \uses{def:localmat, def:entry}
  \lean{LWX.UpDatum, LWX.UpDatum.matrix, LWX.UpDatum.op, LWX.UpDatum.matrixCoeff_op}
  \leanok
  A \emph{$U_p$-datum} on a class set $\iota$ consists of, for each $i\in\iota$, exactly $p$
  summands, each given by a target index $\mathrm{tgt}(i,j)\in\iota$ and a local matrix
  $\mathrm{mat}(i,j)$ of $U_p$-shape.  This is the content of \cite[Proposition 3.1]{LWX} that
  the estimate consumes; the column count of clause (2) is not needed.  Assembling the entry
  streams block by block gives a matrix on $\iota\times\N$ and, by cofinite column decay, a
  continuous operator $D.\mathrm{op}$ on $c(\iota\times\N, \Lhalo)$.
\end{definition}

\section{Theorem 3.16 and Corollary 3.18}

\begin{definition}
  \label{def:lwxlambda}
  \lean{LWX.lwxLambda, LWX.lwxLambda_succ, LWX.lwxLambda_eq_sum_comp}
  \leanok
  With $t = \lvert\iota\rvert$, set
  $\lambda(n) = \sum_{k<n}\bigl(\lfloor k/t\rfloor - \lfloor k/pt\rfloor\bigr)$.
\end{definition}

\begin{theorem}[{\cite[Theorem 3.16]{LWX}}]
  \label{thm:lwx316}
  \uses{def:updatum, def:lwxlambda, lem:two-sided, prop:entry-bound, def:iscompactoid}
  \lean{LWX.norm_charCoeff_upOp_le, LWX.exists_charCoeff_upOp_eq_T_pow_mul,
        LWX.summable_minor_upOp}
  \leanok
  For every $U_p$-datum $D$ on a class set of size $t$ and every nebentypus $\omega$, the
  Fredholm determinant $\Char(P) = \sum_n c_n X^n$ of $D.\mathrm{op}$ is defined over $\Lhalo$
  and
  \[
    \lVert c_n\rVert \;\leq\; p^{-\lambda(n)},
    \qquad\text{equivalently}\qquad
    c_n \in T^{\lambda(n)}\,\Lhalo .
  \]
\end{theorem}

\begin{proof}
  \leanok
  \uses{lem:two-sided, prop:entry-bound, def:lwxlambda, prop:haloint-norm}
  Proposition~\ref{prop:entry-bound} says the matrix entry at $\bigl((i,m),(j,n)\bigr)$ has norm
  at most $p^{-\max(m-\lfloor n/p\rfloor,\,0)}$: a two-sided weight, with row weight
  $w(i,m) = m$ and column weight $w'(j,n) = \lfloor n/p\rfloor$.  Lemma~\ref{lem:two-sided}
  therefore bounds $\lVert c_n\rVert$ by $p^{-f(n)}$ for any $f$ with
  $f(n)\leq \sum_{a\in S}w(a) - \sum_{a\in S}w'(a)$ over $n$-element $S$, and the combinatorial
  content is that $\lambda$ is such an $f$: on $\iota\times\N$ with $\lvert\iota\rvert = t$, the
  smallest possible value of $\sum_{(i,m)\in S}\bigl(m - \lfloor m/p\rfloor\bigr)$ over
  $\lvert S\rvert = n$ is attained by the first $n$ indices in the block order, which gives
  $\sum_{k<n}\bigl(\lfloor k/t\rfloor - \lfloor \lfloor k/t\rfloor/p\rfloor\bigr) = \lambda(n)$.
  Summability of the minors is the same estimate.  Finally $T$-divisibility is
  Proposition~\ref{prop:haloint-norm}.
\end{proof}

\cite{LWX} obtain the estimate by conjugating $P$ by $\mathrm{diag}(1,\dots,T,\dots,T^2,\dots)$
and reading off column divisibility.  Lemma~\ref{lem:two-sided} performs the same computation at
the level of minors, which are invariant under diagonal conjugation, so no $T^{-1}$ ever appears
and the whole argument stays inside the integral ring.

\begin{definition}
  \label{def:speccharseries}
  \uses{thm:lwx316, def:halo-specialize}
  \lean{LWX.specCharSeries, LWX.specCharSeries_coeff_zero, LWX.lwxSlopes,
        LWX.monotone_lwxSlopes}
  \leanok
  At a halo point $T_0$ with $p^{-1}<\lVert T_0\rVert<1$ set
  $\Char(P)(T_0) = \sum_n c_n(T_0)X^n \in K\llbracket X\rrbracket$.  Its constant coefficient is
  $1$.  The comparison slope sequence is
  $s_k = \bigl(\lfloor k/t\rfloor - \lfloor k/pt\rfloor\bigr)\cdot v(T_0)$, which is monotone.
\end{definition}

\begin{theorem}[{\cite[Corollary 3.18]{LWX}}]
  \label{thm:lwx318}
  \uses{def:speccharseries, thm:lwx316, def:ofslopes, def:isnewtonpolygonof, thm:ofslopes}
  \lean{LWX.norm_specCharSeries_coeff_le, LWX.isBelow_newtonPolygon_specCharSeries}
  \leanok
  For every $T_0$ in the boundary annulus, $v\bigl(c_n(T_0)\bigr) \geq \lambda(n)\,v(T_0)$, and
  the Newton polygon of $\Char(P)(T_0)$ lies on or above the polygon anchored at the origin
  whose $k$-th unit slope is $\bigl(\lfloor k/t\rfloor - \lfloor k/pt\rfloor\bigr)v(T_0)$.
\end{theorem}

\begin{proof}
  \leanok
  \uses{thm:lwx316, def:halo-specialize, thm:ofslopes, def:isnewtonpolygonof}
  The coefficient bound is Theorem~\ref{thm:lwx316} pushed through the specialisation estimate of
  Definition~\ref{def:halo-specialize}.  For the polygon statement, the constant coefficient is
  $1$ so the polygon is anchored at the origin, all coefficients have norm at most $1$ so the
  valuation sequence is admissible, and $\sum_{i<k}s_i = \lambda(k)\,v(T_0) \leq
  v(c_k(T_0))$ is the coefficient bound after taking logarithms; now apply the maximality clause
  of Definition~\ref{def:isnewtonpolygonof} to the comparison polygon of
  Definition~\ref{def:ofslopes}.
\end{proof}

\section{Sharpness}

The lower bound of Theorem~\ref{thm:lwx318} is attained exactly when a single integral
coefficient is a unit.  Write $c_n = \sum_m b_{n,m}T^{m}$ for the $T$-expansion of the integral
coefficient.

\begin{theorem}[The equality clause of {\cite[Corollary 3.18]{LWX}}]
  \label{thm:lwx-sharp}
  \uses{thm:lwx318, def:speccharseries}
  \lean{LWX.norm_specCharSeries_coeff_eq_iff, LWX.norm_specCharSeries_coeff_eq_of_isUnit,
        LWX.norm_specCharSeries_coeff_le_of_not_isUnit}
  \leanok
  $v\bigl(c_n(T_0)\bigr) = \lambda(n)v(T_0)$ if and only if $b_{n,\lambda(n)} \in \Zp^{\times}$.
  Otherwise there is a definite margin:
  \[
    v\bigl(c_n(T_0)\bigr) \;\geq\; \lambda(n)\,v(T_0) + \min\bigl(v(T_0),\, 1-v(T_0)\bigr).
  \]
\end{theorem}

\begin{proof}
  \leanok
  \uses{thm:lwx318}
  The term-by-term estimate \cite[(3.18.1)]{LWX} shows that in
  $c_n(T_0) = \sum_m \psi(b_{n,m})T_0^{\,m}$ the term $m = \lambda(n)$ strictly dominates all
  others when $b_{n,\lambda(n)}$ is a unit: terms with $m < \lambda(n)$ are smaller because
  integrality forces $v(b_{n,m}) \geq \lambda(n)-m$, and terms with $m > \lambda(n)$ are smaller
  because $\lVert T_0\rVert<1$.  An ultrametric sum with a unique dominant term has that term's
  norm.  If $b_{n,\lambda(n)}$ is not a unit then every term loses either a factor
  $\lVert T_0\rVert$ or a factor $p^{-1}\lVert T_0\rVert^{-1}$, which is the stated margin.
\end{proof}

The if-and-only-if is the hinge of Chapter~\ref{chap:lwx-slopes}: it converts a geometric
statement about the Newton polygon at one halo point into the $T$-free arithmetic statement that
a particular integral coefficient is a unit --- and a $T$-free statement can then be transported
to \emph{every} halo point.

\section{Riesz theory at a halo vertex}
\label{sec:tate-riesz}

Over the Banach--Tate ring $A$ of Proposition~\ref{prop:halotate} the halo $U_p$ is an honest
compact operator, and Chapter~\ref{chap:riesz-coleman} applies to it.

\begin{definition}
  \label{def:tateop}
  \uses{prop:halotate, def:updatum, def:iscompactoid}
  \lean{LWX.UpDatum.tateMatrix, LWX.UpDatum.tateOp, LWX.UpDatum.isCompactoid_tateOp}
  \leanok
  Conjugate the matrix $P$ by $\mathrm{diag}(T^{\lfloor n/t\rfloor - \lfloor n/pt\rfloor})$, as
  \cite[proof of Theorem 3.16]{LWX} does, and transpose.  The result is a compactoid operator on
  $c(\iota\times\N, A)$.
\end{definition}

The transpose is not cosmetic: after rescaling it is the \emph{columns} of $P'$ that decay
uniformly, so the honest compact operator in the sense of Definition~\ref{def:iscompactoid} ---
which asks for row decay --- is $(P')^{\mathsf{T}}$.

\begin{proposition}
  \label{prop:tateop-det}
  \uses{def:tateop, lem:diag-intertwine}
  \lean{LWX.UpDatum.minor_tateOp, LWX.UpDatum.charPowerSeries_tateOp,
        LWX.UpDatum.charCoeff_tateOp}
  \leanok
  Principal minors are invariant under transposition and diagonal conjugation, so the Fredholm
  determinant of the operator of Definition~\ref{def:tateop} is $\Char(P)$ read in
  $A\llbracket X\rrbracket$.
\end{proposition}

\begin{definition}
  \label{def:halovertex}
  \uses{def:tateop, thm:lwx-sharp}
  \lean{LWX.UpDatum.IsHaloVertex, LWX.HaloInt.isUnit_of_isUnit_coeff_zero,
        LWX.UpDatum.isHaloVertex_of_isUnit_coeff}
  \leanok
  An index $n$ is a \emph{halo vertex} if $c_n = T^{\lambda(n)}\cdot(\text{unit})$ in $\Lhalo$.
  This is the situation of \cite[Remark 3.25]{LWX}; by Theorem~\ref{thm:lwx-sharp} it is implied
  by $b_{n,\lambda(n)}$ being a unit, that is, by the Newton polygon touching the lower bound at
  $n$.
\end{definition}

\begin{theorem}[Riesz decomposition at a halo vertex]
  \label{thm:halo-riesz}
  \uses{def:halovertex, thm:slope-factorisation, thm:riesz-coleman, prop:tateop-det}
  \lean{LWX.UpDatum.isDominantIndex_charPowerSeries_tateOp,
        LWX.UpDatum.isMultiplicative_charCoeff_tateOp,
        LWX.UpDatum.exists_isDominantFactorization_tateOp,
        LWX.UpDatum.exists_rieszColemanProjection_tateOp}
  \leanok
  At a halo vertex $n$, the index $n$ is dominant for $\Char(P)$ at radius
  $p^{\lambda(n)-\lambda(n-1)}$ and the dominant coefficient is a multiplicative unit.  Hence
  $\Char(P) = P_n\,G$ with $\deg P_n = n$ (Theorem~\ref{thm:slope-factorisation}), and
  Theorem~\ref{thm:riesz-coleman} gives the corresponding decomposition of
  $c(\iota\times\N, A)$: a $U_p$-stable summand of rank $n$ carrying exactly the first $n$
  slopes, over the whole boundary annulus at once.
\end{theorem}

\begin{proof}
  \leanok
  \uses{def:halovertex, thm:lwx316, thm:slope-factorisation, thm:riesz-coleman}
  Dominance at the stated radius is convexity of $\lambda$: the increments
  $\lfloor k/t\rfloor - \lfloor k/pt\rfloor$ are non-decreasing, so $\lambda(k) - k\nu$ is
  minimised at $k = n$ for $\nu$ the $n$-th increment, and the inequality is strict beyond $n$
  because the increments grow.  The coefficient bound of Theorem~\ref{thm:lwx316} gives the
  inequality for $k \neq n$ and the vertex hypothesis gives equality at $n$ together with
  multiplicativity of the norm on $c_n$, since $\lVert T\rVert$ is multiplicative and units of
  $\Lhalo$ are detected by their leading coefficient.  Now apply
  Theorem~\ref{thm:slope-factorisation} and then Theorem~\ref{thm:riesz-coleman}.
\end{proof}

\cite[Remark 3.25]{LWX} appeals to \cite[Proposition 3.2.2]{Kedlaya} at exactly this point; the
hypothesis that $n$ is a halo vertex is derived in \cite{LWX} from the touching argument of
Theorem 1.3, which is \S\ref{sec:stepone} below, so we carry it as an explicit hypothesis.


%%FILE blueprint/src/chapter/LWXModel.tex
\chapter{The Integral Model and the Model Seam}
\label{chap:lwx-model}

Theorem~\ref{thm:lwx316} is an estimate about a matrix.  For it to say anything about modular
forms one has to know two things: that the matrix really is the matrix of $U_p$ on a space of
overconvergent automorphic forms, and that its Fredholm determinant really is the Fredholm
determinant of that operator.  These are \cite[Proposition 3.1]{LWX} and
\cite[Proposition 2.17]{LWX}, and they are what this chapter proves.

The difficulty is that the two sides use different coordinates.  The integral side of
Chapter~\ref{chap:lwx-halo} works in the Mahler basis $\binom{z}{n}$ of $C(\Zp, \Lhalo)$, in
which the entries are integral but the operator is not visibly compact; the analytic side of
Chapter~\ref{chap:qmf} works in the monomial basis of the Tate algebra, in which the operator is
visibly compact but the entries are not integral.  \cite[proof of Proposition 2.17]{LWX} bridges
them in one sentence: ``$P$ and $P'$ are conjugated by an infinite diagonal matrix with entries
$\lfloor n/(q^{-1}p^{m})\rfloor!$''.  The intermediate basis is Colmez's, and the statement that
it is orthonormal is Amice's theorem.  So the chapter is: analysis
(\S\ref{sec:explog}--\S\ref{sec:haloweight}), the integral model (\S\ref{sec:intmodel}), Amice's
theorem (\S\ref{sec:amice}), and the seam (\S\ref{sec:seam}).

\section{$p$-adic analysis over an abstract field}
\label{sec:explog}

\begin{definition}
  \label{def:padicexplog}
  \lean{LWX.PadicExpLog.padicLog, LWX.PadicExpLog.padicExp, LWX.PadicExpLog.padicLog_mul,
        LWX.PadicExpLog.padicExp_add, LWX.PadicExpLog.padicExp_padicLog,
        LWX.PadicExpLog.padicLog_padicExp}
  \leanok
  Over a complete ultrametric field $K$ with $\lVert p\rVert<1$ and $p$ odd,
  $\log u = \sum_{n\ge0}(-1)^{n}\tfrac{(u-1)^{n+1}}{n+1}$ for $\lVert u-1\rVert<1$ and
  $\exp w = \sum_n w^{n}/n!$ for $\lVert w\rVert^{2}<\lVert p\rVert$.  On the joint disc they are
  mutually inverse, and $\log$ is additive.
\end{definition}

The convergence discs are \cite[Chapter IV \S2]{Koblitz}.  Additivity of $\log$ is obtained from
$\exp\circ\log = \mathrm{id}$ and injectivity rather than by rearranging series.  Odd $p$ enters
only through two elementary inequalities on $\N$, $4v\leq p^{v}+3$ and the Legendre margin
$2v_p(n!)\leq n-1$; both fail at $p=2$.

\begin{theorem}[The $p$-adic binomial theorem at an arbitrary exponent]
  \label{thm:padic-binomial}
  \uses{def:padicexplog}
  \lean{LWX.PadicExpLog.hasSum_choose_mul_pow, LWX.PadicExpLog.chooseCoeff,
        LWX.PadicExpLog.norm_chooseCoeff_le,
        LWX.PadicExpLog.eq_zero_of_forall_tsum_natCast_eq_zero}
  \leanok
  For $e, x \in K$ with $\lVert x\rVert\leq c$, $\lVert e\rVert\lVert x\rVert\leq c$ and
  $c^{2}<\lVert p\rVert$,
  \[
    \sum_{r\geq 0}\binom{e}{r}x^{r} \;=\; \exp\bigl(e\log(1+x)\bigr).
  \]
\end{theorem}

\begin{proof}
  \leanok
  \uses{def:padicexplog}
  Fix $x$ and regard both sides as power series in the exponent $e$, with geometrically decaying
  coefficients.  They agree at every $e = n \in \N$, by the finite binomial theorem on the left
  and $\exp(n\log u) = u^{n}$ on the right.  A power series vanishing at all natural numbers
  vanishes identically, by Strassman's theorem applied along $p^{k}\to 0$.
\end{proof}

The exponent has to be arbitrary: in the application it is either the normalised logarithm
$\ell\langle x\rangle$ of a $1$-unit, which need not lie in $\Zp$ when $K/\Qp$ is ramified, or
the halo exponent $s = \log(1+T_0)/p$, whose norm $p\lVert T_0\rVert$ exceeds $1$.  So the
density-of-$\N$ proof used in Chapter~\ref{chap:jacobs} does not apply, and the identity theorem
above replaces it.

\begin{lemma}[The binomial power identity]
  \label{lem:pow-sub-one}
  \uses{thm:padic-binomial}
  \lean{LWX.oneAddPow_pow_mul, LWX.norm_pow_prime_pow_sub_one_le,
        LWX.tendsto_pow_prime_pow_sub_one, LWX.exists_sq_norm_pow_prime_pow_sub_one_lt}
  \leanok
  $(1+T)^{p^{h}u} = \bigl((1+T)^{p^{h}}\bigr)^{u}$ for $u\in\Zp$; and
  $(1+T)^{p^{h}} - 1 \to 0$ as $h\to\infty$.  Hence every halo point reaches the joint
  $\exp/\log$ disc at some analyticity level.
\end{lemma}

\begin{proof}
  \leanok
  Both sides of the identity are continuous in $u$ and agree on $\N$, where they are the ordinary
  binomial theorem, and $\N$ is dense in $\Zp$.  For the decay, iterate the one-step Kummer bound
  $\lVert (1+T)^{p}-1\rVert \leq \max(\lVert p\rVert\lVert T\rVert, \lVert T\rVert^{p})$, whose
  contraction factor is $\max(\lVert p\rVert, \lVert T\rVert^{p-1})<1$.
\end{proof}

The second half is the quantitative form of \cite[\S2.7]{LWX}'s statement that the universal
character is $m$-locally analytic for $m$ large enough, and it is what lets
Theorem~\ref{thm:lwx-prop217-h} cover the whole annulus.

\section{The halo weight}
\label{sec:haloweight}

\begin{definition}
  \label{def:haloweight}
  \uses{def:univchar, thm:padic-binomial, def:analyticweight}
  \lean{LWX.haloUnits, LWX.haloChar, LWX.M1K, LWX.haloRho, LWX.haloCol, LWX.haloExpansion,
        LWX.haloWeight, LWX.evalAt_autFactor_haloWeight}
  \leanok
  The specialisation of the universal character at $T_0$ is
  $\kappa_{T_0}(x) = \omega(\bar x)\,(1+T_0)^{\ell\langle x\rangle}
  = \omega(\bar x)\exp\bigl(s\log\langle x\rangle\bigr)$, with halo exponent
  $s = \log(1+T_0)/p$.  On the sub-annulus $p^{-1}<\lVert T_0\rVert$, $\lVert T_0\rVert^{2}<p^{-1}$
  it is an analytic weight in the sense of Definition~\ref{def:analyticweight}: the group is
  $\psi(\Zp^{\times})(1+p\calO_K)$, the level is the image of $M_1$, and the expansion of
  $\kappa_{T_0}(cz+d)$ on the closed unit ball is
  \[
    (cz+d)^{2}\,\kappa_{T_0}(d)\sum_{m\geq 0}\binom{s}{m}\Bigl(\frac{c}{d}\Bigr)^{m}z^{m},
    \qquad \text{decay } \rho = \lVert T_0\rVert\sqrt{p} .
  \]
\end{definition}

The factor $(cz+d)^{2}$ is the normalisation of Definition~\ref{def:jacobs-weight}, chosen so
that the automorphy factor $\kappa(cz+d)/(cz+d)^{2}$ of Definition~\ref{def:kappaslash} is
literally \cite[(2.3.2)]{LWX}'s $\chi(cz+d)$.  The expansion is
Theorem~\ref{thm:padic-binomial}, and the constraint $\lVert T_0\rVert^{2}<p^{-1}$, that is
$v(T_0)>\tfrac12$, is exactly the condition for its convergence.  \cite[\S2.7]{LWX} record only
the cruder radius and remark that ``the radius here is not optimal''.

\begin{definition}[The halo weight at analyticity level $h$]
  \label{def:haloweighth}
  \uses{def:haloweight, lem:pow-sub-one}
  \lean{LWX.Mh, LWX.M1Kh, LWX.TH, LWX.haloExponentH, LWX.haloCharH, LWX.haloRhoH,
        LWX.haloWeightH, LWX.specialize_univChar_eq_padicExp,
        LWX.evalAt_autFactor_haloWeightH}
  \leanok
  With $h = m-1$, $T'_h = (1+T_0)^{p^{h}}-1$ and $s_h = \log(1+T'_h)/p^{h+1}$, the character on
  the disc of radius $p^{-(h+1)}$ about a $\Zp$-unit $a$ is
  $x\mapsto [a](T_0)\exp\bigl(s_h\log(x/a)\bigr)$, an analytic weight at the level
  $M_1^{(h)}$ of matrices with $p^{h+1}\mid c$, of radius
  $\max(p^{-(h+1)}, \lVert T'_h\rVert)\sqrt p$.
\end{definition}

That the two descriptions agree, that is
$(1+T_0)^{\log u/p} = (1+T'_h)^{\log u/p^{h+1}} = \exp(s_h\log u)$ on $1+p^{h+1}\Zp$, is
Lemma~\ref{lem:pow-sub-one}.  This is \cite[\S2.1]{LWX}: ``$\chi(a\cdot x) =
\chi(a)\cdot\chi(\exp(p^{m}))^{(\log x)/p^{m}}$''.

\section{The integral model $S^{D}_{\mathrm{int}}$}
\label{sec:intmodel}

\cite[\S2.7]{LWX} define the space of integral $p$-adic automorphic forms as the
$\Iw_q$-equivariant functions valued in $\Ind^{\Iw_q}_{B(\Zp)}[-]$, identified by
\cite[(2.3.1)]{LWX} with $C(\Zp;\Lambda)$ carrying the action
\[
  (h\Vert_\delta)(z) \;=\; [\,cz+d\,]\cdot h\Bigl(\frac{az+b}{cz+d}\Bigr)
  \tag{2.3.2}
\]
of the monoid $M_1$.  The repository's Hecke layer (Chapter~\ref{chap:qmf}) is stated over an
arbitrary semiring, so it can be instantiated at $R = \Lhalo$ and reused rather than re-proved.

\begin{definition}
  \label{def:intforms}
  \uses{def:univchar, def:localmat, def:forms}
  \lean{LWX.CFun, LWX.cfunSlash, LWX.cfunSlashAction, LWX.IntForms, LWX.intAutSlashAction}
  \leanok
  $C(\Zp,\Lhalo)$ with the action (2.3.2) of $M_1$, and $S^{D}_{\mathrm{int}}$ the corresponding
  space of automorphic functions at a tame level $U$ --- the same construction as
  Definition~\ref{def:forms}, at the ring $\Lhalo$ instead of a field.
\end{definition}

\begin{proposition}[{\cite[Proposition 3.4]{LWX}}]
  \label{prop:prop34}
  \uses{def:intforms, def:entry}
  \lean{LWX.mahlerON, LWX.mahlerCoeffs, LWX.cfunSlash_mahlerON, LWX.entry_eq_fwdDiff,
        LWX.mahlerEmbed}
  \leanok
  The Mahler basis $\binom{z}{n}$ realises $C(\Zp,\Lhalo)$ as $c(\N,\Lhalo)$, and in these
  coordinates the matrix of $\Vert_\delta$ is exactly the entry stream $P_{m,n}(\delta)$ of
  Definition~\ref{def:entry}.  The $T$-rescaled submodule
  $\widehat{\bigoplus}_n T^{n}\Lhalo\binom{z}{n}$ of \cite[(5.4.1)]{LWX} is realised by an
  explicit embedding.
\end{proposition}

\begin{proof}
  \leanok
  \uses{def:entry, def:univchar}
  Mahler's theorem applies because $\Lhalo$ is an ultrametric $\Zp$-Banach algebra.  The Mahler
  coefficients of a function are its iterated forward differences at $0$, so the coefficient of
  $\binom{z}{n}$ in $\binom{z}{n'}\Vert_\delta$ is the $m$-th difference at $0$ of
  $\binom{f(z)}{n'}[cz+d]$, which is Definition~\ref{def:entry} verbatim; the $T$-expansion of
  $[cz+d]$ is the one performed there.
\end{proof}

So the definition of Chapter~\ref{chap:lwx-halo}, which was taken as a definition there, is a
theorem here.  Note that the rescaling of \cite[(5.4.1)]{LWX} lives in the estimates and not in
the model coordinates: the action (2.3.2) does not preserve $c_0$-coordinates in the rescaled
basis.

\begin{theorem}[{\cite[(2.11.1)]{LWX} and \cite[Proposition 3.1]{LWX}}]
  \label{thm:prop31}
  \uses{def:intforms, def:updatum, thm:evalatreps, def:heckeoperator}
  \lean{LWX.intEvalAtReps, LWX.intEvalAtReps_comm,
        LWX.bijective_intEvalAtReps_of_stabilizer_eq_bot, LWX.certM1, LWX.UpDatum.ofCerts,
        LWX.isUpShape_certM1, LWX.intHeckeOperator, LWX.intEvalAtReps_intHeckeOperator}
  \leanok
  At a family of class representatives with trivial stabilisers, evaluation at the
  representatives is a bijection $S^{D}_{\mathrm{int}} \to c(\iota\times\N,\Lhalo)$; and for the
  genuine integral Hecke operator $[U\eta U]$ the diagram
  \[
    \mathrm{ev}\circ[U\eta U] \;=\; D.\mathrm{op}\circ\mathrm{ev}
  \]
  commutes, where $D$ is the $U_p$-datum assembled from right-coset certificates of $U\eta U$.
  Its local matrices have $U_p$-shape because they lie in
  $\Iw_q\begin{pmatrix} p&0\\0&1\end{pmatrix}\Iw_q \subseteq
  \bigl(\begin{smallmatrix} p\Zp&\Zp\\ q\Zp&\Zp^{\times}\end{smallmatrix}\bigr)$.
\end{theorem}

\begin{proof}
  \leanok
  \uses{thm:evalatreps, def:heckeoperator, def:updatum}
  \cite[Proposition 3.1]{LWX} writes $\gamma_i v_j^{-1} = \delta_{i,j}^{-1}\gamma_{\lambda_{i,j}}
  u_{i,j}$ and finds $(U_p\varphi)(\gamma_i) = \sum_j \varphi(\gamma_{\lambda_{i,j}})
  \Vert_{u_{i,j,p}v_j}$.  That display is already proved in the ring-generic Hecke layer
  (Theorem~\ref{thm:evalatreps} and its companion for $[U\eta U]$), so what is left is to check
  that any operator satisfying it intertwines evaluation with $D.\mathrm{op}$, which is a
  computation with the block decomposition, and to verify the shape of
  $\delta_{i,j,p} = u_{i,j,p}v_j$ from the Iwahori double coset.
\end{proof}

\section{Amice's theorem and the Colmez basis}
\label{sec:amice}

\cite[\S2.16]{LWX} quote \cite[Th\'eor\`eme 1.29]{Colmez}: the functions
$\lfloor n/(q^{-1}p^{m})\rfloor!\binom{z}{n}$ form an orthonormal basis of the space
$\mathcal{O}B_{qp^{-m}}$ of functions analytic on the discs $a+q^{-1}p^{m}\Zp$.  This is Amice's
theorem \cite{Amice}.  We prove it at every level, in the form the seam needs.

\begin{proposition}[Level $m=1$]
  \label{prop:colmez-m1}
  \uses{def:tate-algebra, def:cspace}
  \lean{LWX.colmezToMonomial, LWX.norm_colmezToMonomial, LWX.colmezEquiv,
        LWX.monomialToMahler, LWX.diagFactorial, LWX.monomialToMahler_eq,
        LWX.mahlerCoeffPow, LWX.pow_eq_sum_choose_mul_mahlerCoeffPow,
        LWX.fwdDiff_iter_evalAt_natCast}
  \leanok
  At $m=1$ (so $p$ odd, $q=p$) the disc is the closed unit disc and the space is the Tate algebra
  $K\langle z\rangle = c(\N,K)$; the Colmez basis is
  $n!\binom{z}{n} = z(z-1)\cdots(z-n+1)$, a monic integral polynomial.  Change of coordinates
  from the Colmez basis to the monomial basis is an isometry with dense range, hence a continuous
  linear equivalence, and
  \[
    \mathrm{monomialToMahler} \;=\; \mathrm{diag}(m!)\circ\mathrm{colmezEquiv}^{-1} .
  \]
\end{proposition}

The last display is the ``conjugated by an infinite diagonal matrix'' of
\cite[proof of Proposition 2.17]{LWX} at $m=1$.  The isometry is proved by a largest-index
argument; the same argument at every level is Theorem~\ref{thm:unitriangular}.

\begin{lemma}[{\cite[Lemme 1.4.9]{Colmez}}]
  \label{lem:colmez-valuation}
  \lean{LWX.colmezPoly, LWX.discPoly, LWX.discFactor, LWX.discConst, LWX.discPoly_eq,
        LWX.norm_discConst, LWX.norm_coeff_discPoly_le, LWX.norm_coeff_discPoly_le_inv_of_lt,
        LWX.norm_coeff_discPoly_diag}
  \leanok
  Let $g_n = \lfloor n/p^{h}\rfloor!\binom{z}{n}$ and let $a < p^{h}$ be a disc centre.  The
  polynomial $g_{n,a}(w) = g_n(a+p^{h}w)$ has $\Zp$-coefficients, and factors as
  $g_{n,a} = c_{n,a}f_{n,a}$ where $f_{n,a}$ reduces mod $p$ to a monic polynomial of degree
  $\lvert K_{n,a}\rvert$, $K_{n,a} = \{k<n : k\equiv a \pmod{p^{h}}\}$, and
  $v(c_{n,a}) = \sum_{\ell=1}^{h}[\,a \bmod p^{\ell} < n\bmod p^{\ell}\,] \geq 0$.
\end{lemma}

\begin{proof}
  \leanok
  Each linear factor $a-k+p^{h}w$ is $p^{h}(w-(k-a)/p^{h})$ when $p^{h}\mid k-a$, and otherwise
  $(a-k)\bigl(1+p^{h}w/(a-k)\bigr)$; collecting the two kinds of factor gives the factorisation,
  and the valuation of the scalar is the count of carries, which is the displayed sum.
\end{proof}

\cite{Colmez} centres the discs at $-j$ with $1\leq j\leq p^{h}$; we centre at the natural
representative $a = n\bmod p^{h}$, which is a translation of the unit disc and changes neither
integrality, nor vanishing mod $p$, nor degrees.

\begin{theorem}[Amice's theorem at level $h$]
  \label{thm:amice}
  \uses{lem:colmez-valuation, thm:unitriangular, def:cspace}
  \lean{LWX.discIdx, LWX.revIdx, LWX.colmezMatrix,
        LWX.isUnitriangularPerturbation_colmezMatrix, LWX.colmezToDisc,
        LWX.colmezDiscEquiv, LWX.diagFactorialH, LWX.discToMahler,
        LWX.discToMahler_apply_eq_fwdDiff}
  \leanok
  Model $\mathcal{O}B_{p^{-h}}$, the functions analytic on every disc $a+p^{h}\Zp$ with the
  maximum of the $p^{h}$ Gauss norms, by $c(\Z/p^{h}\times\N, K)$ --- the Taylor coefficients on
  each disc in the coordinate $z = a+p^{h}w$.  Then the map sending $e_n$ to the disc model of
  $g_n$ is an isometric equivalence $c(\N,K) \to c(\Z/p^{h}\times\N,K)$, and composing with
  $\mathrm{diag}(\lfloor n/p^{h}\rfloor!)$ turns Colmez coordinates into Mahler coefficients.
\end{theorem}

\begin{proof}
  \leanok
  \uses{lem:colmez-valuation, thm:unitriangular}
  Order the basis by $n = (m(n)+1)p^{h} - i(n)$ with $1\leq i(n)\leq p^{h}$, and the discs by
  their centres, as \cite{Colmez} does; in that order Lemma~\ref{lem:colmez-valuation} says the
  reduction of the basis-change matrix is block upper triangular with each diagonal block lower
  triangular with invertible diagonal entries.  Reindexing the columns by the involution
  reversing each block of $p^{h}$ consecutive integers turns the matrix into a perturbed
  unitriangular one, so Theorem~\ref{thm:unitriangular} applies.  The two reindexings cancel, so
  the resulting map sends $e_n$ to $g_n$ on the nose.  The final clause is the level-$h$ form of
  the statement that Mahler coefficients are forward differences of values at $\N$-points.
\end{proof}

\begin{definition}[The disc model of $S^{D,\dagger,m}$]
  \label{def:discforms}
  \uses{thm:amice, def:haloweighth, def:forms, thm:qmf-compact}
  \lean{LWX.discImage, LWX.discConj, LWX.discSlash, LWX.discSlashAction, LWX.DiscForms,
        LWX.discEvalAtReps, LWX.bijective_discEvalAtReps_of_stabilizer_eq_bot,
        LWX.discHeckeOperator, LWX.discHeckeBlockOp, LWX.isCompactoid_discHeckeBlockOp,
        LWX.discHeckeCharPowerSeries, LWX.evalT_discHeckeCharPowerSeries_eq_zero_iff}
  \leanok
  The action (2.3.2) does not preserve a single disc: it \emph{permutes} the discs.  Writing
  $t_a(w) = a+p^{h}w$ and $a' = \delta(a)\bmod p^{h}$, the matrix identity
  $\delta\, t_a = t_{a'}\,\delta'$ with $\delta' = t_{a'}^{-1}\delta t_a$ of level $h$ says that
  the action carries the Taylor expansion on $a'$ to the one on $a$ through the single-disc
  weight action of Definition~\ref{def:haloweighth}.  So the disc action is a block operator of
  $\kappa$-slashes permuted by $a\mapsto a'$, and it is a right action because the two data
  satisfy the cocycle rules.  Everything in Chapter~\ref{chap:qmf} then transcribes: the space
  $S^{D,\dagger,m}$, evaluation at representatives, the Hecke operator, its compactness, its
  Fredholm determinant, and the eigenvalue reading of the zeros.
\end{definition}

\section{Proposition 2.17: the seam}
\label{sec:seam}

\begin{theorem}[{\cite[Proposition 2.17]{LWX}} at $m=1$]
  \label{thm:lwx-prop217}
  \uses{prop:colmez-m1, def:haloweight, thm:prop31, lem:diag-intertwine,
        thm:charpowerseries-inv, def:heckecharpowerseries}
  \lean{LWX.specOp, LWX.charPowerSeries_specOp, LWX.monomialToMahler_comp_kappaSlash,
        LWX.monomialToMahlerBlock_comp_heckeBlockOp,
        LWX.specCharSeries_ofCerts_eq_heckeCharPowerSeries}
  \leanok
  On the sub-annulus $p^{-1}<\lVert T_0\rVert$, $\lVert T_0\rVert^{2}<p^{-1}$, the integral
  characteristic series evaluated at $T_0$ is the Fredholm determinant of $U_p$ on the
  overconvergent forms of the halo weight:
  \[
    \Char(P)(T_0) \;=\; \det\bigl(1 - X\,[U\eta U]\bigr)
    \quad\text{on } S^{D,\dagger,1}_{\kappa_{T_0}}(U).
  \]
\end{theorem}

\begin{proof}
  \leanok
  \uses{prop:colmez-m1, lem:diag-intertwine, thm:charpowerseries-inv, prop:prop34}
  Four steps.  (i) The specialised operator $P(T_0)$ on $c(\iota\times\N,K)$ has
  $c_n\bigl(P(T_0)\bigr) = c_n(P)(T_0)$, because specialisation is a continuous ring homomorphism
  (Definition~\ref{def:halo-specialize}) and the coefficients are limits of determinants of
  minors.  (ii) The \emph{seam identity}: monomial-to-Mahler coordinates intertwine the analytic
  action of a certificate matrix with $P(T_0)$ --- both are $f\mapsto \chi(cz+d)f(\text{M\"ob}\,z)$
  read at $\N$-points --- and hence blockwise for $[U\eta U]$.  (iii) By
  Proposition~\ref{prop:colmez-m1}, $\mathrm{monomialToMahler} = \mathrm{diag}(m!)\circ
  \mathrm{colmezEquiv}^{-1}$, so the Colmez-basis matrix $P'$ satisfies
  $\mathrm{diag}(m!)P' = P(T_0)\mathrm{diag}(m!)$, and Lemma~\ref{lem:diag-intertwine} gives
  $\det(1-XP') = \det(1-XP(T_0))$ --- with no inversion of the unbounded diagonal.  (iv) Basis
  independence (Theorem~\ref{thm:charpowerseries-inv}) gives
  $\det(1-XP') = \det(1-X[U\eta U])$.
\end{proof}

\begin{theorem}[{\cite[Proposition 2.17]{LWX}} at every level]
  \label{thm:lwx-prop217-h}
  \uses{thm:lwx-prop217, thm:amice, def:discforms, lem:pow-sub-one}
  \lean{LWX.discToMahler_comp_discSlash, LWX.discToMahlerBlock_comp_discHeckeBlockOp,
        LWX.specCharSeries_ofCerts_eq_discHeckeCharPowerSeries,
        LWX.exists_level_specCharSeries_eq_heckeCharPowerSeries}
  \leanok
  The same identity holds at every analyticity level $m = h+1$, with $S^{D,\dagger,m}$ in its
  disc model.  Consequently every point of the full boundary annulus
  $p^{-1}<\lVert T_0\rVert<1$ is covered, at a level supplied by Lemma~\ref{lem:pow-sub-one}.
\end{theorem}

\begin{proof}
  \leanok
  \uses{thm:amice, def:discforms, lem:diag-intertwine, thm:charpowerseries-inv}
  The proof of Theorem~\ref{thm:lwx-prop217} with the disc model in place of the Tate algebra:
  Theorem~\ref{thm:amice} supplies the Colmez equivalence and the diagonal
  $\mathrm{diag}(\lfloor n/p^{h}\rfloor!)$, the disc action is (2.3.2) pointwise so on Mahler
  coefficients it is the specialised entry stream, and the two determinant identities are as
  before.
\end{proof}

\begin{corollary}[{\cite[Lemma 4]{Bu04}; \cite[Definition 2.13]{LWX}}]
  \label{cor:level-independence}
  \uses{thm:lwx-prop217-h}
  \lean{LWX.discHeckeCharPowerSeries_eq_of_levels,
        LWX.discHeckeCharPowerSeries_zero_eq_heckeCharPowerSeries}
  \leanok
  The characteristic power series of $U_p$ on $S^{D,\dagger,m}$ is independent of $m$: both
  levels compute $\Char(P)(T_0)$.
\end{corollary}

This is a statement the $m=1$ seam could not even express, and it is what makes
\cite[Definition 2.13]{LWX}'s spectral curve well defined.

\section{The quaternionic case}

Chapter~\ref{chap:qmf}'s quaternionic instantiation works over a number field $F$ at a place $v$
with $K = F_v$; the integral model needs $\Zp$-entries.  At $F=\Q$, $v = (p)$ the two are
compared by mathlib's isomorphism, once one knows it is an isometry.

\begin{lemma}
  \label{lem:adic-completion-isometry}
  \lean{Padic.absNorm_asIdeal_primesEquiv_symm, Padic.norm_adicCompletionEquiv}
  \leanok
  Mathlib's continuous $\Q$-algebra isomorphism $\Qp \simeq \Q_v$ at the place $v = (p)$ is an
  isometry for the finite-place norm.
\end{lemma}

\begin{proof}
  \leanok
  The absolute norm of the ideal $v$ is $p$, so $\lVert p\rVert = p^{-1}$ on both sides; and the
  closed unit balls correspond, by mathlib's comparison of $\Zp$ with the ring of integers of the
  completion.
\end{proof}

This generalises to every prime the computation made at $p=3$ in Chapter~\ref{chap:jacobs}.

\begin{theorem}[The spectral reading at a halo point]
  \label{thm:lwx-spectral}
  \uses{thm:lwx-prop217, thm:lwx-prop217-h, lem:adic-completion-isometry, thm:serre12,
        def:rigidification, thm:qmf-eigenform-criterion}
  \lean{LWX.thetaInt, LWX.IntFormsQ, LWX.intHeckeUpQ,
        LWX.specCharSeries_ofCerts_eq_heckeCharPowerSeriesQ,
        LWX.evalT_specCharSeries_eq_zero_iff,
        LWX.specCharSeries_ofCerts_eq_discHeckeCharPowerSeriesQ,
        LWX.evalT_specCharSeries_eq_zero_iff_disc,
        LWX.exists_level_evalT_specCharSeries_eq_zero_iff}
  \leanok
  Let $D$ be a definite quaternion algebra over $\Q$ split at $p$, $U$ a tame level, and suppose
  the level is neat.  Then $S^{D}_{\mathrm{int}}$ is the integral model of
  Definition~\ref{def:intforms} for $D$, and for every halo point $p^{-1}<\lVert T_0\rVert<1$
  there is an analyticity level at which, for $a \neq 0$,
  \[
    \Char(P)(T_0)(a) = 0
    \quad\Longleftrightarrow\quad
    a^{-1} \text{ is a } U_p\text{-eigenvalue on } S^{D,\dagger,m}_{\kappa_{T_0}}(U).
  \]
\end{theorem}

\begin{proof}
  \leanok
  \uses{thm:lwx-prop217-h, lem:adic-completion-isometry, thm:qmf-eigenform-criterion}
  Lemma~\ref{lem:adic-completion-isometry} identifies the integral component map with the
  $K$-side one of Chapter~\ref{chap:qmf}, so $S^{D}_{\mathrm{int}}$ and Buzzard's
  $S^{D}_{\kappa}(U)$ are the two sides of Theorem~\ref{thm:lwx-prop217-h}.  Composing that
  identity with the eigenvalue criterion of Theorem~\ref{thm:qmf-eigenform-criterion} --- itself
  Serre's Theorem~\ref{thm:serre12} in the block model --- gives the displayed equivalence.  The
  level is supplied by Lemma~\ref{lem:pow-sub-one}.
\end{proof}

Theorem~\ref{thm:lwx-spectral} is the fibre of \cite[Definition 2.13]{LWX}'s spectral curve over
a halo point: from here on, every statement proved about $\Char(P)$ in
Chapter~\ref{chap:lwx-halo} or Chapter~\ref{chap:lwx-slopes} is a statement about the genuine
$U_p$ on a genuine space of overconvergent quaternionic forms.


%%FILE blueprint/src/chapter/LWXSlopes.tex
\chapter{Slopes over the Boundary of Weight Space}
\label{chap:lwx-slopes}

We can now state what \cite{LWX} prove.  Their target is the Coleman--Mazur--Buzzard--Kilford
conjecture \cite[Conjecture 1.2]{LWX}, transposed from the eigencurve to the spectral curve
$\mathrm{Spc}_D$ of overconvergent automorphic forms for a definite quaternion algebra.  That
conjecture has three clauses: sufficiently close to the boundary,
\begin{enumerate}
  \item $\mathrm{Spc}^{>r}_D$ is a disjoint union of countably many connected components $Z_n$,
    each finite and flat over the annulus $\mathcal{W}^{>r}$ via the weight map;
  \item there are rationals $\alpha_1 \leq \alpha_2 \leq \cdots \to \infty$ with
    $\lvert a_p(z)\rvert = \lvert T_{\mathrm{wt}(z)}\rvert^{\alpha_n}$ for $z \in Z_n$;
  \item the sequence $\alpha_1, \alpha_2, \dots$ is a disjoint union of \emph{finitely many
    arithmetic progressions}, counted with multiplicity (at least for large indices).
\end{enumerate}
\cite[Theorem 1.3]{LWX} proves a coarse form of (1)--(2) over the full annulus
$\mathcal{W}^{>1/p}$: a decomposition indexed by the intervals $[n,n]$ and $(n,n+1)$, with
$v(a_p(x)) \in \varphi(q)\,v(T_{\mathrm{wt}(x)})\cdot I$ on the piece $X_I$, and explicit degrees.
\cite[Theorem 1.5]{LWX} proves (1)--(2) on the nose over a smaller annulus
$\mathcal{W}^{>\lambda}$, $\lambda = p^{-8/((p^{2}-1)t+8)}$, in the form
\[
  v\bigl(a_p(y)\bigr) \;=\; \varphi(q)\,v\bigl(T_{\mathrm{wt}(y)}\bigr)\,\alpha_i(\omega),
  \qquad y \in Y_{i,\omega},
  \tag{1.5.1}
\]
and then, in a second half beginning ``Moreover'', proves (3).  This is the \emph{spectral halo}:
approaching the boundary of weight space, the spectral curve resolves into pieces of a completely
explicit shape.

Theorem~\ref{thm:lwx318} gave a \emph{lower} bound polygon.  The content of this chapter is that,
under an explicit and checkable hypothesis, the bound is attained on a band around each of the
points $n_k = kpt$, and that this pins the whole polygon.  Three ingredients:
a combinatorial upper bound (\S1), a vertex analysis converting touching into $T$-free
arithmetic (\S2), and the analytic Claim of \cite[\S4.2]{LWX} which controls the indices in
between (\S3).  Section~4 assembles the polygon-level form of Theorem 1.5, and
\S\ref{sec:what-15-says} says precisely how much of \cite[Theorem 1.5]{LWX} that is.
Section~5 transports everything to the genuine Hecke operator, and
Section~\ref{sec:stepone} discharges the hypothesis (Steps I and III of \cite[\S3.23]{LWX}),
granted adelic Atkin--Lehner data, and proves the degree formulas of \cite[Theorem 1.3]{LWX},
\cite[Corollary 1.4]{LWX} and the second half of \cite[Theorem 1.5]{LWX} at the level of Newton
polygons.

\section{The upper bound polygon}

Everything here is $\N$-combinatorics of the exponent
$\lambda(n) = \sum_{k<n}\bigl(\lfloor k/t\rfloor - \lfloor k/pt\rfloor\bigr)$ of
Definition~\ref{def:lwxlambda}.

\begin{definition}
  \label{def:touchx}
  \uses{def:lwxlambda}
  \lean{LWX.touchX, LWX.two_mul_lwxLambda_touchX, LWX.lwxUpperTwice}
  \leanok
  Set $n_k = k\,p\,t$, the abscissa of the $k$-th \emph{touching vertex}.  Then
  \cite[(3.23.1)]{LWX}
  \[
    \lambda(n_k) \;=\; \frac{k^{2}p(p-1)t}{2},
  \]
  and the \emph{upper bound polygon} is the polygon through the points $(n_k, \lambda(n_k))$,
  linear of slope $(k+\tfrac12)(p-1)$ on $[n_k, n_{k+1}]$.  We work with twice its exponent, so
  as to stay in $\N$.
\end{definition}

\begin{lemma}[{\cite[Lemma 4.1]{LWX}}]
  \label{lem:lwx41}
  \uses{def:touchx}
  \lean{LWX.two_mul_lwxLambda_le_lwxUpperTwice,
        LWX.lwxUpperTwice_sub_two_mul_lwxLambda_le}
  \leanok
  The upper and lower polygons agree at every $n_k$ and differ by at most
  $(p^{2}-1)t/8$ in between:
  \[
    0 \;\leq\; \mathrm{upper}(n) - \lambda(n) \;\leq\; \frac{(p^{2}-1)t}{8}
    \qquad\text{for all } n .
  \]
\end{lemma}

So the two polygons are uniformly close; the whole strategy is to show the Newton polygon is
squeezed between them and touches the lower one at the $n_k$.

\begin{lemma}[The band identities, {\cite[p.~25]{LWX}}]
  \label{lem:band}
  \uses{def:touchx}
  \lean{LWX.lwxLambda_touchX_add, LWX.lwxLambda_touchX_sub, LWX.lwxLambda_touchX_add_ge,
        LWX.lwxLambda_touchX_sub_ge}
  \leanok
  $\lambda(n_k\pm i) = \lambda(n_k) \pm k(p-1)i$ for $0\leq i\leq t$, and beyond the band of
  width $t$ the excess is at least $\lvert i\rvert - t$.
\end{lemma}

The band identity is the reason the lower bound polygon is \emph{linear} of slope $k(p-1)$ across
$[n_k-t, n_k+t]$: the increments $\lfloor n/t\rfloor - \lfloor n/pt\rfloor$ are constant there.

\section{Vertices and Theorem 1.3}

Write $c_n = \sum_m b_{n,m}T^{m}$ as in \S\ref{chap:lwx-halo}, and recall from
Theorem~\ref{thm:lwx-sharp} that the polygon touches the lower bound at $n$ exactly when
$b_{n,\lambda(n)}$ is a unit.  This condition does not mention $T$.

\begin{definition}
  \label{def:unitband}
  \uses{thm:lwx-sharp, def:touchx}
  \lean{LWX.IsUnitCoeff, LWX.leftIndex, LWX.rightIndex, LWX.HasUnitBand}
  \leanok
  Say $n$ is a \emph{unit index} if $b_{n,\lambda(n)}\in\Zp^{\times}$.  Let
  $n_k^{-}$ be the least unit index in $[n_k-t, n_k]$ and $n_k^{+}$ the greatest in
  $[n_k, n_k+t]$.  The \emph{touching hypothesis} $\mathrm{HasUnitBand}(k)$ asserts that both
  exist.
\end{definition}

\begin{lemma}
  \label{lem:unitband-from-touching}
  \uses{def:unitband, thm:lwx318, thm:lwx-sharp}
  \lean{LWX.hasUnitBand_of_height_eq, LWX.height_specCharSeries_touchX}
  \leanok
  $\mathrm{HasUnitBand}(k)$ holds if and only if the Newton polygon of $\Char(P)(T_0)$ touches
  the lower bound polygon at $(n_k, \lambda(n_k)v(T_0))$ for one --- equivalently every --- halo
  point $T_0$.
\end{lemma}

This equivalence is the pivot of the whole chapter.  Touching is a statement at one weight; the
unit band is a statement about the integral matrix and holds for all weights at once.  So a
single weight at which one can verify touching --- which is what \cite[\S3.23 Step I]{LWX}
supplies from Atkin--Lehner theory and classicality --- propagates to the entire annulus.

\begin{theorem}[{\cite[Theorem 1.3]{LWX}}, the vertex reading]
  \label{thm:lwx13-vertices}
  \uses{def:unitband, thm:lwx318, thm:lwx-sharp, lem:band, def:isnewtonpolygonof}
  \lean{LWX.bandLine, LWX.height_specCharSeries_eq_bandLine,
        LWX.bandLine_lt_height_specCharSeries_of_lt_leftIndex,
        LWX.bandLine_lt_height_specCharSeries_of_rightIndex_lt}
  \leanok
  Assume $\mathrm{HasUnitBand}(k)$.  Then at every halo point $T_0$ the points
  $\bigl(n_k^{-}, \lambda(n_k^{-})v(T_0)\bigr)$ and $\bigl(n_k^{+}, \lambda(n_k^{+})v(T_0)\bigr)$
  are consecutive vertices of the Newton polygon of $\Char(P)(T_0)$, joined by the segment of
  slope $k(p-1)v(T_0)$ through $(n_k,\lambda(n_k)v(T_0))$; strictly outside $[n_k^-, n_k^+]$ the
  polygon lies strictly above that line.
\end{theorem}

\begin{proof}
  \leanok
  \uses{thm:lwx-sharp, lem:band, thm:lwx318, def:isnewtonpolygonof}
  On the band, Lemma~\ref{lem:band} says the lower bound polygon is the straight line of slope
  $k(p-1)v(T_0)$ through $(n_k, \lambda(n_k)v(T_0))$, and Theorem~\ref{thm:lwx-sharp} says the
  valuation of $c_n(T_0)$ equals its lower bound exactly at the unit indices --- in particular at
  $n_k^{\pm}$.  So the polygon meets the line at those two abscissae, hence coincides with it in
  between by convexity.  Outside the band, the excess of Lemma~\ref{lem:band} plus the definitive
  margin of Theorem~\ref{thm:lwx-sharp} at non-unit indices puts every point strictly above the
  line; the greatest-convex-minorant characterisation of
  Definition~\ref{def:isnewtonpolygonof} then gives the strict inequality for the polygon.
\end{proof}

\begin{corollary}[{\cite[(3.23.3)--(3.23.4)]{LWX}}]
  \label{cor:lwx13-slopes}
  \uses{thm:lwx13-vertices}
  \lean{LWX.unitSlope_specCharSeries_eq_iff, LWX.unitSlope_specCharSeries_mem_Ioo}
  \leanok
  The $j$-th slope of the polygon equals $k(p-1)v(T_0)$ exactly for
  $n_k^{-}\leq j < n_k^{+}$, and lies strictly between $k(p-1)v(T_0)$ and $(k+1)(p-1)v(T_0)$ for
  $n_k^{+}\leq j < n_{k+1}^{-}$.  In particular the multiplicities
  $\deg X_k = n_k^{+}-n_k^{-}$ and $\deg X_{(k,k+1)} = n_{k+1}^{-}-n_k^{+}$ are independent of
  $T_0$.
\end{corollary}

The independence is what \cite[Theorem 1.3]{LWX} needs for its component decomposition: the
pieces of the eigencurve over the annulus have degrees that do not vary.

\section{The Claim of \S4.2}

Corollary~\ref{cor:lwx13-slopes} controls the slopes on the bands and says the remaining slopes
lie strictly between consecutive values.  To pin those too, \cite[\S4.2]{LWX} restrict to a
smaller annulus, $0 < v(T) < 8/\bigl((p^{2}-1)t+8\bigr)$, and prove:

\begin{theorem}[{\cite[\S4.2, Claim]{LWX}}]
  \label{thm:lwx-claim}
  \uses{lem:lwx41, thm:lwx-sharp, def:touchx}
  \lean{LWX.unitIndex, LWX.IsBelowUpper, LWX.exists_isUnit_and_coeffVal_eq,
        LWX.isBelowUpper_of_coeffVal_lt, LWX.coeffVal_specCharSeries_eq_unitIndex_mul,
        LWX.le_coeffVal_specCharSeries_of_not_isBelowUpper}
  \leanok
  Let $m(l)$ be the least $m$ with $b_{l,m}\in\Zp^{\times}$.  If $(l, v(c_l(T_0)))$ lies strictly
  below the upper bound polygon for one $T_0$ in the small annulus, then $m(l)\geq\lambda(l)$,
  the same holds for every $T$ in the annulus, and
  \[
    v\bigl(c_l(T)\bigr) \;=\; m(l)\cdot v(T) \qquad\text{for all such } T .
  \]
  If it does not, then $v(c_l(T))\geq \mathrm{upper}(l)\cdot v(T)$ for every $T$.
\end{theorem}

\begin{proof}
  \leanok
  \uses{lem:lwx41, thm:lwx-sharp}
  This is the analytic core (4.2.1)--(4.2.4) of \cite{LWX}.  Being strictly below the upper bound
  forces, by Lemma~\ref{lem:lwx41}, that the dominant term in
  $c_l(T_0) = \sum_m \psi(b_{l,m})T_0^{m}$ occurs at an index $m$ with
  $\lambda(l)\leq m$ and $2m < \mathrm{upper}$ doubled; the smallness of the annulus makes the
  competition between the factors $p^{-v(b_{l,m})}$ and $\lVert T_0\rVert^{m}$ strict, so the
  dominant index is unique and is characterised $T$-freely as the least $m$ with $b_{l,m}$ a
  unit.  Uniqueness of the dominant term gives the valuation identity by the ultrametric
  principle, at every $T$ in the annulus simultaneously, since the characterisation of $m(l)$ no
  longer mentions $T$.
\end{proof}

Both the hypothesis and the conclusion have been made $T$-free: ``strictly below the upper bound
polygon'' becomes the predicate $2\,m(l) < \mathrm{upper}_{2}(l)$, both polygons scaling with
$v(T)$.  That is what allows the next section to define a single $T$-free polygon.

\section{Theorem 1.5}

\begin{definition}
  \label{def:shapepolygon}
  \uses{thm:lwx-claim, def:unitband, def:isnewtonpolygonof}
  \lean{LWX.shapeVal, LWX.shapePolygon, LWX.slopeRatio, LWX.isNewtonPolygonOf_shapePolygon}
  \leanok
  The \emph{shape polygon} is the Newton polygon of the $T$-free valuation sequence
  \[
    n \longmapsto
    \begin{cases}
      \lambda(n_k) & n = n_k, \\
      m(n) & n \text{ strictly below the upper polygon}, \\
      \infty & \text{otherwise},
    \end{cases}
  \]
  and $\mathrm{slopeRatio}(j)$ is its $j$-th unit slope.  It is well defined unconditionally.
\end{definition}

\begin{theorem}[{\cite[Theorem 1.5]{LWX}}, (1.5.1)]
  \label{thm:lwx15}
  \uses{def:shapepolygon, thm:lwx-claim, thm:lwx13-vertices, cor:ofslopes-transport}
  \lean{LWX.height_specCharSeries_eq_smul_shape,
        LWX.unitSlope_specCharSeries_eq_slopeRatio}
  \leanok
  Assume $\mathrm{HasUnitBand}(k)$ for every $k$.  Then for every $T$ in the small annulus
  $0 < v(T) < 8/\bigl((p^{2}-1)t+8\bigr)$ and every $j$, the $j$-th slope of the Newton polygon
  of $\sum_n c_n(T)X^{n}$, counted with multiplicity, is
  \[
    \mathrm{slope}_j(T) \;=\; v(T)\cdot\mathrm{slopeRatio}(j),
  \]
  with $\mathrm{slopeRatio}$ independent of $T$.  In \cite{LWX}'s normalisation
  $\mathrm{slopeRatio}(j) = \varphi(q)\,\tilde\alpha_j(\omega)$, the factor $\varphi(q)$ of
  (1.5.1) being absorbed into our ratio and $\tilde\alpha$ being the sequence $\alpha$ repeated
  according to the degrees $\deg Y_{i,\omega}$.
\end{theorem}

\begin{proof}
  \leanok
  \uses{thm:lwx-claim, thm:lwx13-vertices, def:shapepolygon}
  The Newton polygon of $\Char(P)(T)$ is the convex hull of the points $(n, v(c_n(T)))$.  At
  $n = n_k$ the height is $\lambda(n_k)v(T)$ by Theorem~\ref{thm:lwx13-vertices}; at $n$ strictly
  below the upper polygon it is $m(n)v(T)$ by Theorem~\ref{thm:lwx-claim}; and every other point
  lies above the upper polygon and so above the hull of those two families.  Thus the two hulls
  agree after scaling by $v(T)$, and reading unit slopes off a scaled polygon multiplies them by
  the same scalar.
\end{proof}

\begin{theorem}[The remaining clauses of {\cite[Theorem 1.5]{LWX}}]
  \label{thm:slope-growth}
  \uses{def:shapepolygon, thm:lwx318}
  \lean{LWX.monotone_slopeRatio, LWX.tendsto_slopeRatio_atTop}
  \leanok
  The ratios $\mathrm{slopeRatio}(0)\leq\mathrm{slopeRatio}(1)\leq\cdots$ increase, and tend to
  infinity.  Both statements are unconditional: they need neither the small annulus nor the
  touching hypothesis.
\end{theorem}

\begin{proof}
  \leanok
  \uses{thm:lwx318, def:shapepolygon}
  Monotonicity is convexity of the shape polygon.  For the growth: unit coefficients occur only
  at heights $\geq\lambda(n)$ by Theorem~\ref{thm:lwx318}, so the shape sequence dominates
  $\lambda$, whose increments $\lfloor k/t\rfloor-\lfloor k/pt\rfloor$ are unbounded.  Given $M$,
  the line $y = Mx - M^{2}pt$ lies on or below every point of the shape sequence, hence on or
  below the shape polygon by the supporting-line lemma; comparing with the height as the sum of
  the first $n$ slopes and using monotonicity gives a slope $\geq M-1$.
\end{proof}

The supporting-line lemma used here is a spec-level consequence of the greatest-convex-minorant
characterisation, recorded separately.

\begin{lemma}[Supporting lines]
  \label{lem:supporting-line}
  \uses{def:isnewtonpolygonof, def:ofslopes}
  \lean{IsNewtonPolygonOf.twoSlope_le_height, IsNewtonPolygonOf.line_le_height,
        NewtonPolygon₀.unitSlope_eq_of_height_eq, NewtonPolygon₀.monotone_toReal_unitSlope}
  \leanok
  A two-slope broken line lying on or below every point of an admissible valuation sequence lies
  on or below its Newton polygon; with equal slopes this is the supporting-line statement for a
  single line.  Moreover a unit slope is the corresponding height increment, and the unit slopes
  of a constructed polygon are monotone and never $\bot$.
\end{lemma}

\section{What Theorem 1.5 says, and what \S4 proves}
\label{sec:what-15-says}

Theorem~\ref{thm:lwx15} is a statement about Newton polygons; \cite[Theorem 1.5]{LWX} is a
statement about a rigid analytic space.  It is worth being exact about the gap.  \cite{LWX}'s
theorem reads:

\begin{quote}
  Let $\omega : \Delta \to \Zp^{\times}$ be a character.  Then there exists
  $\lambda \in (0,1)$ such that there exists a sequence of rational numbers
  $\alpha_0(\omega), \alpha_1(\omega), \dots$ in increasing order and tending to infinity such
  that $\mathrm{Spc}^{>\lambda}_{D,\omega}$ is a disjoint union
  $\coprod_{i\geq 0} Y_{i,\omega}$ of rigid analytic spaces finite and flat over
  $\mathcal{W}^{>\lambda}_\omega$ via $\mathrm{wt}$, such that
  \[
    v\bigl(a_p(y)\bigr) = \varphi(q)\,v\bigl(T_{\mathrm{wt}(y)}\bigr)\,\alpha_i(\omega)
    \tag{1.5.1}
  \]
  for every $y \in Y_{i,\omega}$.  More precisely, if the tame level is neat, then we can take
  $\lambda = p^{-8/((p^{2}-1)t+8)}$ for $p > 2$ \dots
  \emph{Moreover}, let $M$ be a positive integer so that $p^{-q/p^{M-1}(p-1)} > \lambda$
  \dots  Then if we extend the sequence $\alpha_0(\omega), \alpha_1(\omega),\dots$ into
  $\tilde\alpha_0(\omega), \tilde\alpha_1(\omega), \dots$ with each $\alpha_i(\omega)$ appearing
  with multiplicity $\deg Y_{i,\omega}$, then
  $\tilde\alpha_0(\omega), \tilde\alpha_1(\omega), \dots$ is a disjoint union of
  $(p-1)p^{M-1}t/2$ arithmetic progressions with (the same) common difference
  $\varphi(q)p^{M}/2q^{2}$.  More precisely,
  \[
    \tilde\alpha_{j + q^{-1}p^{M}t}(\omega\omega_0^{2})
      = \tilde\alpha_j(\omega) + \frac{p^{M}}{q^{2}}
    \qquad\text{for any } j \geq 0.
    \tag{1.5.2}
  \]
\end{quote}

So the theorem is really three assertions, and the arithmetic progressions are the third.

\paragraph{(a) The polygon identity.}  This is what Theorem~\ref{thm:lwx15} proves, and it is
\cite{LWX}'s own reduction: the proof in \cite[\S4.2]{LWX} opens by saying that, the tame level
being neat, ``it suffices to show that for $T \in \Cp$ with
$0 < v(T) < \tfrac{8}{(p^{2}-1)t+8}$, the ratios to $v(T)$ of the slopes (counted with
multiplicity) of the Newton polygon of $\sum_{n\geq0} c_n(T)X^{n}$ are independent of the choice
of $T$.''  Theorem~\ref{thm:lwx15} is that sentence, with the ratios named
$\mathrm{slopeRatio}(j) = \varphi(q)\tilde\alpha_j(\omega)$ and produced as the unit slopes of an
explicit $T$-free polygon (Definition~\ref{def:shapepolygon}), on exactly \cite{LWX}'s annulus.
Theorem~\ref{thm:slope-growth} adds the two clauses ``in increasing order and tending to
infinity''.  We prove it over any complete ultrametric field receiving $\Zp$ isometrically, not
only $\Cp$, and --- through Theorem~\ref{thm:lwx-slopes-seam} --- for the genuine $U_p$ rather
than for the integral matrix.  As stated it is conditional on the touching hypothesis
$\mathrm{HasUnitBand}$; Theorem~\ref{thm:lwx-unitband} discharges that hypothesis at every vertex
and on every disc, granted the adelic data of Definition~\ref{def:al-data}, so that for the genuine
$U_p$ (1.5.1) holds with no further hypothesis.

\paragraph{(b) The geometric decomposition.}  The passage from ``the slope ratios are
$T$-independent'' to ``$\mathrm{Spc}^{>\lambda}_{D,\omega} = \coprod_i Y_{i,\omega}$ with each
$Y_{i,\omega}$ finite flat over the annulus'' is not formalised, and is out of scope for this
thesis: the spectral curve is a rigid analytic space, and no rigid geometry is developed here.
What we do have is its fibrewise shadow.  Theorem~\ref{thm:lwx-spectral} identifies, at each
weight in the annulus, the zeros of $\Char(P)(T_0)$ with the reciprocal $U_p$-eigenvalues on
$S^{D,\dagger,m}_{\kappa_{T_0}}(U)$, so (1.5.1) holds pointwise for every point of the spectral
curve lying over the annulus; and Theorem~\ref{thm:halo-riesz} builds, at each halo vertex, the
finite locally free $U_p$-stable summand over the \emph{whole} annulus $\Lhalo[1/T]$ whose
Fredholm determinant is the corresponding polynomial factor --- which is the algebra underlying
the finite flat piece.  Assembling those into a rigid space, and identifying the summands with
connected components, is the missing step.

\paragraph{(c) The arithmetic progressions.}  This is the ``Moreover'' clause, (1.5.2).  It is
\emph{not} a consequence of Theorem~\ref{thm:lwx15}: (1.5.2) relates the sequence at the
nebentypus $\omega$ to the sequence at $\omega\omega_0^{2}$ --- two \emph{different} characters,
shifted by the square of the Teichm\"uller character --- whereas Theorem~\ref{thm:lwx15} is a
statement at a single fixed $\omega$, one weight disc at a time.  \cite[Remark 1.6(4)]{LWX} says
where it comes from: ``The second half of Theorem 1.5 follows from the first half by classicality
results and Atkin--Lehner theory'' (an argument found independently by Bergdall and Pollack).
That is the route formalised in Theorem~\ref{thm:lwx152}, at the level of Newton polygons.  The
Atkin--Lehner input of \cite[\S3.23 Step I]{LWX} is supplied at conductor $p^{h+1}$ for every
level $h \geq 1$ and at every classical weight at once, by \emph{families} of adelic data
(Definition~\ref{def:al-data}); the slopes of the classical subspace are identified without the
classicality theorem \cite[Proposition 2.15]{LWX}; and the slope reflection, (4.2.5),
(4.2.6) $=$ (1.5.2) and the arithmetic progressions are then bookkeeping in the characters of
$\Delta$.  The same comparison of two weight discs is behind \cite[Corollary 1.4]{LWX}, the
periodicity of $\deg X_{I,\omega}$ modulo $\varphi(q)/2$, which is formalised as
Corollary~\ref{cor:lwx14}.

\paragraph{Degrees.}  Corollary~\ref{cor:lwx13-slopes} gives the multiplicities as
$\deg X_k = n_k^{+} - n_k^{-}$ and $\deg X_{(k,k+1)} = n_{k+1}^{-} - n_k^{+}$, which is what the
vertex analysis produces and is enough for $T$-independence.  \cite[Theorem 1.3]{LWX} additionally
evaluates them, at a neat tame level, as
$\deg X_{n,\omega} = r_{\mathrm{ord}}(\omega^{-1}\omega_0^{2n-2}) + r_{\mathrm{ord}}(\omega\omega_0^{-2n})$
and $\deg X_{(n,n+1),\omega} = qt - r_{\mathrm{ord}}(\omega^{-1}\omega_0^{2n}) -
r_{\mathrm{ord}}(\omega\omega_0^{-2n})$ in terms of dimensions of ordinary subspaces.  Those
identifications are Step III, and they too involve several nebentypi at once.  They are formalised
at the coefficient level, with $r_{\mathrm{ord}}$ the largest index of a unit coefficient of the
characteristic series --- which is how \cite{LWX} prove their Theorem 3.19 --- as
Theorem~\ref{thm:lwx13-degrees}; only their reading as degrees of rigid analytic components is
not, as in (b).

\section{Reading it on the genuine $U_p$}

Chapter~\ref{chap:lwx-model} identified $\Char(P)(T_0)$ with $\det(1-X\,U_p)$ on
$S^{D,\dagger,m}$ at every analyticity level (Theorem~\ref{thm:lwx-prop217-h}).  Composing gives
each of the statements above for the genuine Hecke operator; nothing new is proved, and no
neatness hypothesis is needed, because none of the polygon statements uses one.

\begin{theorem}
  \label{thm:lwx-slopes-seam}
  \uses{thm:lwx-prop217-h, thm:lwx318, cor:lwx13-slopes, thm:lwx15,
        lem:unitband-from-touching}
  \lean{LWX.isBelow_newtonPolygon_discHeckeCharPowerSeries,
        LWX.hasUnitBand_of_height_discHeckeCharPowerSeries_eq,
        LWX.unitSlope_discHeckeCharPowerSeries_eq_iff,
        LWX.unitSlope_discHeckeCharPowerSeries_mem_Ioo,
        LWX.unitSlope_discHeckeCharPowerSeries_eq_slopeRatio,
        LWX.exists_level_unitSlope_discHeckeCharPowerSeries_eq_slopeRatio}
  \leanok
  For $\det(1-X\,U_p)$ on $S^{D,\dagger,m}$, at every point of the boundary annulus
  $p^{-1}<\lVert T_0\rVert<1$: the Newton polygon lies on or above the lower bound polygon of
  Theorem~\ref{thm:lwx318}; if it touches it at $n_k$ at one halo point then
  $\mathrm{HasUnitBand}(k)$ holds; the slope readings of Corollary~\ref{cor:lwx13-slopes} hold;
  and, on the small annulus, so does (1.5.1) of Theorem~\ref{thm:lwx15}.
\end{theorem}

The last item deserves emphasis: \emph{the touching hypothesis has to be verified only once, at a
single weight}, and Lemma~\ref{lem:unitband-from-touching} then supplies it everywhere.

\section{Theorem 1.3: Steps I and III}
\label{sec:stepone}

\cite[Theorem 1.3]{LWX} --- the statement that the spectral curve over the boundary annulus is a
disjoint union of pieces finite and flat over it, with the slope of $a_p$ on the piece $X_I$
lying in $\varphi(q)v(T)\cdot I$ and with explicit degrees --- needs the touching hypothesis to be
discharged.  \cite[\S3.23]{LWX} do this in three steps.  \emph{Step I} compares the slopes on the
whole space with those on the finite-dimensional classical subspace, and uses Atkin--Lehner
theory to pin their total; this is what supplies the touching.  \emph{Step II} is the vertex
analysis of \S2 above.  \emph{Step III} computes the degrees $\deg X_{I,\omega}$ in terms of
dimensions of ordinary subspaces, using Hida theory and the theta exact sequence.  All three steps
are formalised, at the level of characteristic series and Newton polygons, granted the adelic
data of Definition~\ref{def:al-data}; \S\ref{sec:ledger} records what each of them bought.

No step uses Jacquet--Langlands.  \cite{LWX}'s proofs of their Proposition 3.22 (Atkin--Lehner)
and Proposition 2.15 (classicality) both appeal to it; here the first is replaced by a
double-coset computation (Proposition~\ref{prop:al-identity}) and the second is not needed.  A
per-result audit of every use of Jacquet--Langlands, and of our avoidance of it, is maintained
alongside the sources.

\begin{definition}[The conjugate, partner and target nebentypi]
  \label{def:conjchar}
  \lean{LWX.invChar, LWX.teichChar, LWX.partnerChar, LWX.targetChar}
  \leanok
  Let $\omega$ be a character of $\Delta$, $\omega_0$ the Teichm\"uller character and $k \geq 0$,
  and let $(k,\psi)$ be a classical weight on the disc $\omega$, so that
  $\psi|_\Delta = \omega\omega_0^{-k}$.  The \emph{conjugate} nebentypus is
  $\omega^{-1} : u \mapsto \omega(u)^{-1}$; the \emph{partner} nebentypus
  $\omega^{-1}\omega_0^{2k}$ is the disc of the Atkin--Lehner partner $(k,\psi^{-1})$; and the
  \emph{target} nebentypus $\omega\omega_0^{-2k-2}$ is the disc of the theta target
  $(-k-2,\psi)$.  Since every statement of Chapters~\ref{chap:lwx-halo}--\ref{chap:lwx-slopes} is
  proved for an arbitrary $\omega$, these further instances cost nothing, and the same $U_p$-datum
  serves all of them, the nebentypus entering only through the entry streams.
\end{definition}

\begin{proposition}[Atkin--Lehner, reduced]
  \label{prop:atkin-lehner}
  \uses{lem:charpoly-pairing}
  \lean{LWX.atkinLehner, LWX.atkinLehnerConj, LWX.atkinLehner_mul_atkinLehnerConj,
        LWX.atkinLehnerConj_mem_Iw, LWX.roots_charpoly_atkinLehner,
        LWX.norm_roots_charpoly_atkinLehner}
  \leanok
  The element $w = \begin{pmatrix}0&1\\-p^{m}&0\end{pmatrix}$ has determinant $p^{m}$, normalises
  the Iwahori subgroup, inverts the nebentypus, and conjugates $U_p$ to $U'_p$ --- all expressed
  by the inverse-free identity $w\,\gamma^{w} = \gamma\,w$.  \emph{Granting} the operator
  identity $U_p\circ U'_p = p^{k+1}$ on the classical space,
  Lemma~\ref{lem:charpoly-pairing} gives \cite[Proposition 3.22]{LWX}: the $U_p$-slopes at
  nebentypus $\psi$ and at $\psi^{-1}$ pair up to $k+1$.
\end{proposition}

The proof of \cite[Proposition 3.22]{LWX} on \cite[p.~24]{LWX} applies Jacquet--Langlands and
classifies the local component as a principal series; neither is available in Lean.  The operator
identity is instead proved directly, from adelic data of the following kind, by expanding the
composite over the coset representatives of $U_p$.

\begin{definition}[Adelic Atkin--Lehner data]
  \label{def:al-data}
  \lean{LWX.AtkinLehnerData, LWX.AtkinLehnerFamily, LWX.AtkinLehnerDataH,
        LWX.AtkinLehnerFamilyH}
  \leanok
  \emph{Atkin--Lehner data} of conductor $p^{h+1}$ at a nebentypus $\psi$ consist of a section
  $\iota_p : \mathrm{GL}_2(\mathbb{Q}_p) \to G$ of the $p$-component such that the level contains
  the lifts of the Iwahori subgroup, the central element $p$ at $p$ is a global central element
  times an element of the level, and the Atkin--Lehner element of level $p^{h+1}$ normalises the
  disc-$0$ part of the level; together with a Hecke character restricting to $\psi$ on the level.
  A \emph{family} of such data is a single section with a Hecke character for every disc $\omega$
  and every weight $k$.  The $U_p$-representatives are built from the section alone, so one
  $U_p$-datum serves every classical weight of the family at once.
\end{definition}

\begin{proposition}[The Atkin--Lehner identity]
  \label{prop:al-identity}
  \uses{def:al-data, prop:atkin-lehner, prop:theta}
  \lean{LWX.ClassicalDiscForms, LWX.atkinLehnerEquiv, LWX.discHeckeCl_comp_atkinLehner,
        LWX.atkinLehnerHypothesis_of_atkinLehnerData, LWX.discHeckeClH_comp_atkinLehnerH,
        LWX.atkinLehnerHypothesis_of_atkinLehnerDataH}
  \leanok
  Grant Atkin--Lehner data of conductor $p^{h+1}$ at $\psi$.  On the classical forms
  $S^{D}_{k+2}(K^{p}\Iw_{p^{h+1}};\psi)$ --- in the disc model, the forms with values in the
  locally polynomial functions of degree $\leq k$ --- the Atkin--Lehner element induces an
  isomorphism $W$ onto the classical forms at $\psi^{-1}$, and
  \[
    U_p \circ W^{-1} \circ U_p^{(\psi^{-1})} \circ W \;=\; p^{k+1} .
  \]
  Hence the hypothesis of Proposition~\ref{prop:atkin-lehner} holds at the classical points, and
  the $U_p$-slopes at $\psi$ and at $\psi^{-1}$ pair up to $k+1$.
\end{proposition}

\begin{proof}
  \leanok
  \uses{prop:atkin-lehner}
  Expand the composite over the coset representatives of $U_p$.  The term with trivial
  representative is central and contributes $p^{k+1}$.  After conjugation by the Atkin--Lehner
  element every other term is a sum of the nebentypus over the cosets of $1 + p^{h+1}\Zp$ in
  $1 + p^{h}\Zp$, which vanishes because the conductor is exactly $p^{h+1}$.
  Proposition~\ref{prop:atkin-lehner} then gives the pairing.
\end{proof}

\begin{proposition}[The theta operator and the classical subspace]
  \label{prop:theta}
  \uses{def:discforms, def:blockmaps}
  \lean{LWX.thetaDisc, LWX.thetaBlock, LWX.thetaDisc_eq_zero_iff, LWX.thetaBlock_eq_zero_iff,
        LWX.finrank_locPolyDegSubmodule, LWX.bol, LWX.thetaDisc_comp_discHeckeBlock_of_autFactor,
        LWX.thetaBlock_comp_discHeckeBlockOp_of_autFactor}
  \leanok
  In the disc model an element \emph{is} the family of Taylor coefficients on each disc, so
  $\theta^{r}$ --- $r$-fold differentiation, \cite[\S7]{Bu04} --- is the shift
  $(a,j)\mapsto (j+1)\cdots(j+r)\,c(a,j+r)$: a diagonal of integers composed with a shift,
  bounded by $1$ and column-finite, hence continuous, and block-diagonal over the discs.  Its
  kernel at $r = k+1$ is exactly the locally polynomial functions of degree $\leq k$, which is
  the first sentence of \cite[Proposition 4]{Bu04}, of dimension $(k+1)p^{h}$ per disc --- the
  local half of \cite[(3.21.1)]{LWX}.  If the automorphy factors of two weights have the classical
  shape of exponent $k$ and the target shape of exponent $k+2$, with the same constants, then
  $\theta^{k+1}$ intertwines $U_p$ at the first weight with $p^{k+1}U_p$ at the second.
\end{proposition}

The intertwining is \cite[\S7]{Bu04}'s displayed identity: Bol's identity
$\partial^{r}\bigl(u^{j+r}L^{-(j+1)}\bigr) = (\det\gamma)^{r}\,(j+r)_{r}\,u^{j}L^{-(j+r+1)}$, for
$u$ and $L$ the numerator and the denominator of the M\"obius transformation of $\gamma$, summed
over the coset decomposition of $U_p$.  Its hypotheses are placed on the automorphy factors, not on
the characters: a version in which the two weights shared an arbitrary finite part is false, since
$\theta$ differentiates that part when it is not locally constant.  Nothing here depends on
Jacquet--Langlands: \cite[Proposition 4]{Bu04}'s first sentence and its small-slope half are
Jacquet--Langlands-free, and its converse half, which does use it, is never imported.

\begin{proposition}[The theta sequence, in determinant form]
  \label{prop:theta-exact}
  \uses{prop:theta, lem:diag-intertwine, def:conjchar}
  \lean{LWX.IsThetaExact, LWX.TargetData, LWX.charPowerSeries_comp_one_sub_truncation_eq_rescale,
        LWX.isThetaExact_classicalData, LWX.targetData_classicalPoint}
  \leanok
  At the classical point of weight $k$ on the disc $\omega$, the Fredholm determinant of $U_p$ on
  the complement of the classical subspace equals the Fredholm determinant of $p^{k+1}U_p$ at the
  theta target $(-k-2,\psi)$, a halo point on the disc of the target nebentypus
  $\omega\omega_0^{-2k-2}$.
\end{proposition}

This is all that \cite[\S3.23 Step III]{LWX} uses of the exact sequence
$0 \to S^{D}_{k+2} \to S^{D,\dagger}_{(k,\psi)} \xrightarrow{(d/dz)^{k+1}}
S^{D,\dagger}_{(-k-2,\psi)} \to 0$.  Writing $\theta^{k+1}$ as a diagonal composed with a shift,
and choosing a section of the shift, turns the intertwining of Proposition~\ref{prop:theta} into
a \emph{diagonal} intertwining on the complement, and Lemma~\ref{lem:diag-intertwine} gives the
identity of determinants.  The operator $\theta^{k+1}$ is not surjective on the Banach model ---
the inverse diagonal is unbounded --- so the right-exactness of the sequence is neither proved nor
needed.

\begin{proposition}[The finite factor]
  \label{prop:finite-factor}
  \uses{prop:det-summand, lem:charpolyrev}
  \lean{TateFredholm.charPowerSeries_eq_mul_polynomial, LWX.charPowerSeries_eq_mul_of_stable}
  \leanok
  If a compactoid $u$ commutes with an idempotent whose range is finite free, its Fredholm
  determinant splits as the reversed characteristic polynomial of the finite piece times the
  determinant of the complement.
\end{proposition}

Step I uses exactly this, the classical subspace being a finite-dimensional $U_p$-stable piece.

\begin{theorem}[{\cite[\S3.23 Step I]{LWX}}: the touching]
  \label{thm:lwx-stepone}
  \uses{prop:finite-factor, prop:atkin-lehner, prop:theta, thm:lwx318,
        lem:unitband-from-touching}
  \lean{LWX.ClassicalData, LWX.classicalData, LWX.isStepOneTouching_of_atkinLehnerHypothesis,
        LWX.hasUnitBand_of_atkinLehnerHypothesis}
  \leanok
  Let $T_0$ be the classical point of a weight $(k,\psi)$ with $\psi$ of conductor $p^{2}$, and
  suppose that the hypothesis of Proposition~\ref{prop:atkin-lehner} holds between $T_0$ and its
  Atkin--Lehner partner.  Then the Newton polygon of $\Char(P)(T_0)$ passes through
  $\bigl(n_{k+1}, \lambda(n_{k+1})v(T_0)\bigr)$; hence $\mathrm{HasUnitBand}(k+1)$.
\end{theorem}

\begin{proof}
  \leanok
  \uses{prop:finite-factor, prop:atkin-lehner, thm:lwx318, lem:unitband-from-touching}
  By Proposition~\ref{prop:finite-factor} the characteristic series at $T_0$ is the determinant of
  the complement times the reversed characteristic polynomial of $U_p$ on the classical subspace,
  of dimension $n_{k+1} = (k+1)pt$ by Proposition~\ref{prop:theta}; so the height of the polygon at
  $n_{k+1}$ is at most $v(\det U_p)$ on that subspace, and the same holds at the partner point.
  The hypothesis gives $v(\det U_p) + v(\det U'_p) = n_{k+1}(k+1)v(p) = 2\lambda(n_{k+1})v(T_0)$,
  the partner point having the same norm as $T_0$.  With the lower bound of
  Theorem~\ref{thm:lwx318} at both points, both inequalities are equalities, and
  Lemma~\ref{lem:unitband-from-touching} converts touching into $\mathrm{HasUnitBand}(k+1)$.
\end{proof}

\begin{theorem}[The touching hypothesis, discharged]
  \label{thm:lwx-unitband}
  \uses{thm:lwx-stepone, prop:al-identity, def:al-data, thm:lwx-slopes-seam}
  \lean{LWX.hasUnitBand_of_atkinLehnerData, LWX.hasUnitBand_of_atkinLehnerFamily,
        LWX.unitSlope_discHeckeCharPowerSeries_eq_slopeRatio_of_atkinLehnerFamily}
  \leanok
  Grant a family of Atkin--Lehner data of conductor $p^{2}$.  Then $\mathrm{HasUnitBand}(k)$ holds
  for every $k$ and on every disc $\omega$.  Consequently Corollary~\ref{cor:lwx13-slopes} and
  Theorems~\ref{thm:lwx15} and \ref{thm:lwx-slopes-seam} hold with the touching hypothesis replaced
  by the family of data; in particular (1.5.1) holds for the genuine $U_p$ at every analyticity
  level.
\end{theorem}

\begin{proof}
  \leanok
  \uses{thm:lwx-stepone, prop:al-identity}
  At $k = 0$ the band is automatic.  For $k+1$, the classical point of weight $k$ on the disc
  $\omega$ lies in the halo, and Proposition~\ref{prop:al-identity} at the member of the family
  of weight $k$ on the disc $\omega$ supplies the hypothesis of Theorem~\ref{thm:lwx-stepone}
  there; the $U_p$-datum does not depend on $(\omega,k)$ because the family has a single section.
\end{proof}

\begin{theorem}[{\cite[Theorem 1.3]{LWX}}, the degrees]
  \label{thm:lwx13-degrees}
  \uses{thm:lwx-unitband, prop:theta-exact, prop:al-identity, def:conjchar, cor:lwx13-slopes}
  \lean{LWX.ordDim, LWX.ordDim_le_card, LWX.degX, LWX.degXint, LWX.degX_zero, LWX.degX_succ,
        LWX.degXint_zero, LWX.degX_succ_classicalPoint, LWX.degX_succ_of_atkinLehnerData,
        LWX.degX_succ_of_atkinLehnerFamily, LWX.degXint_of_atkinLehnerFamily,
        LWX.degXint_pos_of_atkinLehnerFamily}
  \leanok
  Let $r_{\mathrm{ord}}(\omega)$ be the largest index $d$ such that the $d$-th coefficient of the
  characteristic series on the disc $\omega$ is a unit of the halo ring; it is at most $t$.  With
  $\deg X_{k,\omega} = n_k^{+} - n_k^{-}$ and $\deg X_{(k,k+1),\omega} = n_{k+1}^{-} - n_k^{+}$ as
  in Corollary~\ref{cor:lwx13-slopes}, $\deg X_{0,\omega} = r_{\mathrm{ord}}(\omega)$; and, granted
  a family of Atkin--Lehner data of conductor $p^{2}$, for every disc $\omega$ and every $n \geq 0$
  \[
    \deg X_{n+1,\omega} = r_{\mathrm{ord}}(\omega^{-1}\omega_0^{2n})
      + r_{\mathrm{ord}}(\omega\omega_0^{-2n-2}),
    \qquad
    \deg X_{(n,n+1),\omega} = pt - r_{\mathrm{ord}}(\omega^{-1}\omega_0^{2n})
      - r_{\mathrm{ord}}(\omega\omega_0^{-2n}) \;>\; 0 .
  \]
\end{theorem}

\begin{proof}
  \leanok
  \uses{thm:lwx-unitband, prop:theta-exact, prop:al-identity}
  At $n_{k+1}$ the unit band splits the degree into the left gap $n_{k+1} - n_{k+1}^{-}$ and the
  right gap $n_{k+1}^{+} - n_{k+1}$, and the finite factor splits the polygon: the classical
  slopes are at most $(k+1)v(p)$ and those of the complement at least $(k+1)v(p)$.  The left gap
  counts the classical slopes equal to $(k+1)v(p)$; the Atkin--Lehner pairing of
  Proposition~\ref{prop:al-identity} matches them with the classical slopes $0$ at the partner
  point, which are all its ordinary slopes, so the left gap is
  $r_{\mathrm{ord}}(\omega^{-1}\omega_0^{2k})$.  The right gap counts the slopes of the complement
  equal to $(k+1)v(p)$; by Proposition~\ref{prop:theta-exact} these are the ordinary slopes of the
  theta target, so the right gap is $r_{\mathrm{ord}}(\omega\omega_0^{-2k-2})$.  The interval
  degrees are $pt$ minus the two gaps at consecutive vertices, an exact subtraction because
  $n_k^{+} \leq n_{k+1}^{-}$; positivity is $r_{\mathrm{ord}} \leq t$ with $p \geq 3$.
\end{proof}

\begin{corollary}[{\cite[Corollary 1.4]{LWX}}]
  \label{cor:lwx14}
  \uses{thm:lwx13-degrees, def:conjchar}
  \lean{LWX.degX_succ_mul_teichChar_sq, LWX.degXint_mul_teichChar_sq,
        LWX.degX_succ_mul_teichChar_pow, LWX.degXint_mul_teichChar_pow, LWX.degX_succ_add_period,
        LWX.degXint_add_period}
  \leanok
  Granted a family of Atkin--Lehner data of conductor $p^{2}$,
  $\deg X_{I,\omega} = \deg X_{I+1,\omega\omega_0^{2}}$ for $I = (0,1), 1, (1,2), 2, \dots$, and
  hence $\deg X_{I,\omega}$ is periodic in $I$, on the same range, modulo
  $\varphi(q)/2 = (p-1)/2$.
\end{corollary}

\begin{proof}
  \leanok
  \uses{thm:lwx13-degrees}
  Substitute $(n,\omega) \mapsto (n+1,\omega\omega_0^{2})$ in the formulas of
  Theorem~\ref{thm:lwx13-degrees}: the partner and target nebentypi of
  Definition~\ref{def:conjchar} do not change.  Iterate, and use $\omega_0^{p-1} = 1$.
\end{proof}

The source's list of intervals starts at $(0,1)$ and omits $I = 0$, and rightly: the identity
fails there in general, since $\deg X_{0,\omega} = r_{\mathrm{ord}}(\omega)$ while
$\deg X_{1,\omega\omega_0^{2}} = r_{\mathrm{ord}}(\omega^{-1}\omega_0^{-2}) +
r_{\mathrm{ord}}(\omega)$.

\begin{theorem}[{\cite[Theorem 1.5]{LWX}}, second half, at the polygon level]
  \label{thm:lwx152}
  \uses{thm:lwx-unitband, prop:al-identity, def:al-data, def:conjchar, def:shapepolygon}
  \lean{LWX.atkinLehnerHypothesis_of_atkinLehnerDataH,
        LWX.slopeRatio_add_slopeRatio_partnerChar_of_atkinLehnerFamilyH,
        LWX.slopeRatio_partnerChar_succ, LWX.slopeRatio_mul_teichChar_sq,
        LWX.slopeRatio_mul_teichChar_pow, LWX.slopeRatio_add_period}
  \leanok
  Let $h \geq 1$ satisfy $(p^{2}-1)t + 8 < 8p^{h-1}(p-1)$, and grant families of Atkin--Lehner
  data of conductors $p^{2}$ and $p^{h+1}$ over a single section.  Then for every disc $\omega$,
  every $k \geq 0$, every $i < N = t(k+1)p^{h}$ and every $j \geq 0$,
  \begin{align*}
    \mathrm{slopeRatio}_{\omega^{-1}\omega_0^{2k}}(N-1-i) + \mathrm{slopeRatio}_{\omega}(i)
      &= (k+1)\,p^{h-1}(p-1), \\
    \mathrm{slopeRatio}_{\omega\omega_0^{2}}(j + p^{h}t)
      &= \mathrm{slopeRatio}_{\omega}(j) + p^{h-1}(p-1), \\
    \mathrm{slopeRatio}_{\omega}\bigl(j + \tfrac{p-1}{2}\,p^{h}t\bigr)
      &= \mathrm{slopeRatio}_{\omega}(j) + \tfrac{p-1}{2}\,p^{h-1}(p-1).
  \end{align*}
  With $M = h+1$ and $\mathrm{slopeRatio} = \varphi(q)\tilde\alpha$, the second line is (1.5.2),
  and the third says that $\tilde\alpha_0(\omega), \tilde\alpha_1(\omega), \dots$ is a disjoint
  union of $(p-1)p^{M-1}t/2$ arithmetic progressions with common difference
  $\varphi(q)p^{M}/2q^{2}$.
\end{theorem}

\begin{proof}
  \leanok
  \uses{thm:lwx-unitband, prop:al-identity}
  At a classical point of conductor $p^{h+1}$ in the slope-reading region, the classical subspace
  carries exactly the first $N$ slopes: they are at most $(k+1)v(p)$ by
  Proposition~\ref{prop:al-identity} and $\lVert U_p\rVert \leq 1$, while those of the complement
  are at least $(k+1)v(p)$, through the theta target.  The Atkin--Lehner pairing reverses them
  against the partner nebentypus, which with (1.5.1) at both discs
  (Theorem~\ref{thm:lwx-unitband}) is the first line.  The first line at $(\omega,k)$ and at
  $(\omega,k+1)$, subtracted, is \cite[(4.2.5)]{LWX}; taking the disc $\omega^{-1}\omega_0^{2k}$
  gives (4.2.6), the second line; and $\omega_0^{p-1} = 1$ iterates it into the third.
\end{proof}

\section{The dependency ledger}
\label{sec:ledger}

It is worth recording, once, which input bought which conclusion, since the statements have
genuinely different requirements and the obvious guess --- that Steps I and III together give
everything --- is not right.

\begin{center}
\small
\begin{tabular}{p{0.27\textwidth}p{0.40\textwidth}p{0.24\textwidth}}
  \hline
  \textbf{Target} & \textbf{Input required} & \textbf{Status here} \\
  \hline
  $\mathrm{HasUnitBand}(k)$ for all $k$; hence unconditional
    Corollary~\ref{cor:lwx13-slopes}, Theorems~\ref{thm:lwx15}, \ref{thm:slope-growth},
    \ref{thm:lwx-slopes-seam}
  & Step I at conductor $q^{2}$: Atkin--Lehner \cite[Prop.~3.22]{LWX}, the dimension count
    \cite[(3.21.1)]{LWX}, and the initial-segment comparison (classicality
    \cite[Prop.~2.15]{LWX} turns out not to be needed)
  & formalised, granted adelic data: Theorem~\ref{thm:lwx-unitband} \\
  \hline
  \cite[Thm.~1.3]{LWX}'s degree formulas, and \cite[Cor.~1.4]{LWX}
  & Step III: Hida's \cite[Cor.~3.21]{LWX} at the coefficient level, and the theta exact sequence
    $0 \to S^{D}_{k+2} \to S^{D,\dagger}_{(k,\psi)} \xrightarrow{(d/dz)^{k+1}}
     S^{D,\dagger}_{(-k-2,\psi)} \to 0$ in determinant form
  & formalised at the coefficient level, granted adelic data: Theorem~\ref{thm:lwx13-degrees},
    Corollary~\ref{cor:lwx14} \\
  \hline
  \cite[Thm.~1.5]{LWX}'s second half: (1.5.2) and the arithmetic progressions
  & Step I's Atkin--Lehner input again, but at conductor $p^{M}$ for every admissible $M$,
    together with statements ranging over the $\omega_0^{2}$-orbit of nebentypi
  & formalised at the polygon level, granted adelic data: Theorem~\ref{thm:lwx152} \\
  \hline
  The decomposition $\mathrm{Spc}^{>\lambda}_{D,\omega} = \coprod_i Y_{i,\omega}$, finite flat
  & rigid analytic geometry
  & out of scope for this thesis; see \S\ref{sec:what-15-says}(b) \\
  \hline
\end{tabular}
\end{center}

Two remarks on the third row, since it is the one that is easy to misjudge.  First, (1.5.2) is a
statement about $\tilde\alpha$, the slope-ratio sequence counted with multiplicity --- our
$\mathrm{slopeRatio}$ --- so it is a \emph{polygon-level} statement and needs none of the rigid
geometry of the fourth row.  Second, it does not follow from Steps I and III as such: it uses no
degree formula, but it needs the classical space and the Atkin--Lehner identity at
\emph{arbitrary} level $p^{M}$ rather than at $q^{2}$, and it compares the sequence at $\omega$
with the sequence at $\omega\omega_0^{2}$, whereas
Chapters~\ref{chap:lwx-halo}--\ref{chap:lwx-slopes} work at a single fixed $\omega$.

Concretely, five objects had to be built on top of those chapters, and all five now exist:

\begin{enumerate}
  \item the classical space $S^{D}_{k+2}(K^{p}\Iw_{p^{m}};\psi)$ itself, as the forms with values
    in the locally polynomial functions of degree $\leq k$, at conductor $p^{2}$ and at every
    conductor $p^{h+1}$, with its dimension (Proposition~\ref{prop:theta});
  \item the operator identity $U_p \circ W^{-1} U_p^{(\psi^{-1})} W = p^{k+1}$ on that space
    (Proposition~\ref{prop:al-identity}), which replaces the standing hypothesis of
    Proposition~\ref{prop:atkin-lehner};
  \item the theta operator's intertwining, in the form of Proposition~\ref{prop:theta}, and the
    theta sequence in determinant form (Proposition~\ref{prop:theta-exact});
  \item the initial-segment comparison, which needs no new theory: faces of Newton polygons add
    under products (Theorem~\ref{thm:np-product}), and the classical and complementary slopes are
    separated by $(k+1)v(p)$;
  \item a nebentypus-varying layer: the classical points of conductor $p^{h+1}$ and their halo
    weights, the partner and target nebentypi (Definition~\ref{def:conjchar}), families of adelic
    data (Definition~\ref{def:al-data}), and statements over the $\omega_0^{2}$-orbit
    (Corollary~\ref{cor:lwx14}, Theorem~\ref{thm:lwx152}).
\end{enumerate}

What is \emph{not} formalised is the adelic data itself: Definition~\ref{def:al-data} is carried as
a hypothesis.  For the finite ad\`eles of a definite quaternion algebra each of its fields is class
field theory or a local computation at $p$, but that instantiation has not been carried out.

A remark on the prime $2$.  \cite{LWX} \emph{state} Theorems~1.3 and~1.5 for $p = 2$ as well
(with $q = 4$, $\Delta = (\mathbb{Z}/4)^\times$, and the analyticity thresholds $m \ge 2$,
$M \ge 4$), but every proof in the paper is written for $p > 2$ only --- ``we assume $p > 2$, the
case $p = 2$ being similar'' (\S4.1 and \S4.2).  There is therefore no written argument to
transcribe for $p = 2$, and it is used essentially throughout this development from the
exponential and logarithm onwards; $p \ne 2$ is a standing hypothesis here, and the results
formalised are the cases the paper actually proves, which are narrower than its stated theorems.

Items (1)--(4) are Step I and Step III; item (5) is what the second half of
\cite[Theorem 1.5]{LWX} and \cite[Corollary 1.4]{LWX} need on top of them.

To summarise the status of \cite{LWX}: \S2 and \S3's estimates
(Theorems~\ref{thm:lwx316}, \ref{thm:lwx318}, \ref{thm:lwx-sharp}), the model identifications
(Theorems~\ref{thm:prop31}, \ref{thm:lwx-prop217}, \ref{thm:lwx-prop217-h}), the slope theorems
(Corollary~\ref{cor:lwx13-slopes}, Theorems~\ref{thm:lwx15}, \ref{thm:slope-growth},
\ref{thm:lwx-slopes-seam}), the touching hypothesis (Theorem~\ref{thm:lwx-unitband}), the degree
formulas of \cite[Theorem 1.3]{LWX} (Theorem~\ref{thm:lwx13-degrees}), \cite[Corollary 1.4]{LWX}
(Corollary~\ref{cor:lwx14}) and the second half of \cite[Theorem 1.5]{LWX}
(Theorem~\ref{thm:lwx152}) are formalised and free of \texttt{sorry}, for $p \neq 2$, granted the
adelic data of Definition~\ref{def:al-data}, and with no use of Jacquet--Langlands.  What remains
is the instantiation of that data for a definite quaternion algebra, and the passage from Newton
polygons to rigid analytic spaces.


%%FILE blueprint/src/chapter/New/RPS.tex
\chapter{Restricted Power Series}

Recall formal multivariate power series over a set $K$ in $n$ variables, written
$K [ \! [ x_1, \dots, x_n ] \! ]$ are infinite sums
\[
  f(x) = \sum a_{(i_1, \cdots, i_n)} x_1^{i_1} \cdots x_n^{i_n},
\]
where the sum is taken over all $n$-tuples $(i_1, \cdots i_n)$ with $i_j \in \mathbb{N}$ and
$a_{(i_1, \cdots, i_n)} \in K$.

When $K$ is a ring, formal multivariate power series form a ring, and we shall refer to them as the
ring of power series over $K$.  We again denote such rings using $K [ \! [ x_1, \dots, x_n ] \! ]$.

Note, henceforth we shall drop the word multivariate and just say power series for simplicity.

We define a subset of $K [ \! [ x_1, \dots, x_n ] \! ]$ that consists of power series whose coefficients
approach zero in the cofinite filer.

\begin{definition}
  Let $f$ be a power series over a nonarchimedean normed ring $R$ in $n$ variables, write
  $f(x) = \sum a_{(i_1, \cdots, i_n)} x_1^{i_1} \cdots x_n^{i_n}$, and let
  $c = (c_1, \cdots, c_n)$ be an $n$-tuple of real numbers.
  We say $f$ is a restricted power series of parameter $c$ over $R$ in $n$ variables if
  \[
    \lim_{n \to \infty} \left\lvert a_{(i_1, \cdots, i_n)} \right\rvert \prod_{j=1}^{n} c_j^{i_j} = 0
  \]
  in the cofinite topology.

  We denote the ring of restricted power series of parameter $c$ over $R$ in $n$ variables by
  $R_c \langle \! x_1, \dots, x_n \! \rangle$.
\end{definition}

The key specification of this will be when $c = 1$ and $R$ is a nonarchimedean normed field. We will
call such rings the Tate algebra and denote it as $T_n$.

Note there is an alternative definition \cite{...} %%Def from wikipedia for now... think Judiths
%% notes could be used for tis

\begin{definition}
  Let $A$ be a linearly topologised ring, seperated and complete and $\{ I_\lamda \}$ the fundamental
  system of open ideals. Then the ring of restricted power series is defined as the projective limit
  of the polynomial rings over $A / I_\lambda$:
  \[
    A \langle \! x_1, \dots, x_n \! \rangle = \varprojlim_\lambda A / I_\lambda [x_1, \dots, x_n].
  \]
\end{definition}

It is not hard to see that we can identify such rings with our initial definition since each $I_\lambda$
contains all but finitely many coefficients. We decide to formalise the definition in Lean as the former
as this will be more convenient for computations. But API saying it is equivalent to this projective
limit should be added to Mathlib eventually.

In literature, the ring of restricted power series are sometimes only defined for non-negative $c$-tuples,
\cite{Gouvea}, however we note by the following lemma that this does not matter and keeping with
Mathlib generality we allow any $c$.

\begin{lemma}
  A power series $f$ is restricted for parameter $c$ if and only if it is restricted for parameter
  $\lvert c \rvert$.
\end{lemma}

\section{Structure}

In the definition we call restricted power series a ring, so we better hope they are a ring. Most of
the ring axioms follow immediately from the definition of being restricted, the exceptions are
closure of the operations.
Addition follows from the usual squeeze theorem, but multiplication is of
more interest.
The nuance comes from coefficients of the product being antidiagonal sums of coefficients
of the two power series, and it may not be immediately clear that we will get the wanted restricted
condition.

To prove closure of multiplication we will use the two more general lemmas.

We first give the definition of the antidiagonal property on additive monoids.

\begin{definition} % HasAntidiagonal Algebra.Order.Antidiag
  An additive monoid, $M$ is said to have the antidiagonal property if there exists a function
  $ \phi : M \to {S | S \subset M \times M, \; S \text{ finite}}$ such that for $(s_1, s_2) \in \phi (m)$
  $s_1 + s_2 = m$.
\end{definition}

A typical example of a monoid having the antidiagonal property is $\mathbb{N}$, but it should be
immediate to see that $\mathbb{Z}$ does not. This property's importance can be seen in definitions
of Laurent series and Hahn series. %% Could explain what I mean by this


%% I think I may need a few definitions discussing how Lean handles filters and convergence
%% It is worth remarking that with Leans vast filter library we rarely need to resort to hard
%% analysis
\begin{lemma}
  Let $M$ be an additive monoid that has the antidiagonal property and $f : M \times M \to R$ be
  a function to a linearly ordered set, that has a filter $F$. If for every neighbourhood
  $a$ in $F$, the preimage of $a$ under $f$ is a neighbourhood hood in the cofinite filter,
  %% If I give a definition of tendsto I can avoid this
  then for every neighbourhood of $a$ in $F$, the primeage of $a$ under the image of the supremum of the
  antidiagonal of $a$ under $f$ is a neighbourhood in the cofinite filter.
\end{lemma}

\begin{lemma}
  Let $M$ be an additive monoid with the antidiagonal property; $S$ a nonarchimedean normed
  ring; $C : M \to \R$ be a function such that $C (a + b) = C (a) * C (b)$; and $f, g : M \to S$.
  Then if we have the functions $i \mapsto \lvert f(i)  \rvert * C(i)$ and
  $i \mapsto \lvert g(i) \rvert * C(i)$ tend to zero under the cofinite filter (that is for all
  neighbourhoods of 0 the preimage of the function is a neighbourhood in the cofinite filter). Then
  the function $a \mapsto \lvert \sum_{p \in \text{antidiagonal} a} f (p.1) g (p.2) \rvert * C (a)$
  tends to zero under the cofinite filter.
\end{lemma}

The closure of multiplication is then immediate taking $C$ to be $i \mapsto \prod c_i^i$ and $f, g$
the coefficient function mapping a power series to its coefficients, since the conditions about tending
to zero are just the definition of $f,g$ being restricted power series.

\section{Remarks}

As with many things there are decisions on typing in Mathlib, because we want to work with the ring
of restricted power series as their own thing, and not just power series that are restricted, we
first show the structure properties on power series with a predicate saying that they are restricted.
%% e.g. add the IsRestricted def here
After we show they are a subring of power series we then promote them to their own type. However,
note this means given something of type power series we have a tuple of information: a power series;
and a proof that it is restricted. Therefore, if we want to apply results on power series to the
type of restricted power series we can just apply them on the first term.


%%FILE blueprint/src/chapter/New/intro.tex
\chapter{Introduction}


%% Definitely need to explain this section better, but it serves as motivation for what we are doing
A key area of research in modern algebraic number theory are the so called eigenvarieties, these
are rigid analytic spaces that lie above a weight space, whose points are isomorphic to
generalised eigenvectors of Hecke operators.

A point in weight space is a character, and its fiber will be family of forms of weight given by the
character. The isomorphism comes from the points in the eigenvariety being characters, which
determine the eigenvalues for the form with respect to the $T_n$ operator. Fixing a specific operator,
say $U_p$, we can look at the fiber of a character and then get a family of eigenvectors for $U_p$.

Formalising the required eigenvariety machinery is out of scope for this thesis. However, studying
what the family of eigenforms looks like for a given character is an interesting question. It turns
out this question can have surprising answers - and our goal of this thesis is to formalise
such an answer.

To do this we will restrict ourselves to quaternionic modular forms over definite quaternion algebras.
Jacquet-Langlands suggests that these should relate to certain families of classical modular forms.
However, compared to the analytical nature of the classical forms, quaternionic modular forms are
much nicer to work with, being combinatorial in nature.

The aim is to define quaternionic modular forms and their Hecke operators in Lean, and then following
\cite{Jacobs} show that the slopes (eigenvalues) of compact Hecke operators are in arithmetic
progression. That is, without defining eigenvarieties, we are able to discuss what the geometry
if a fiber looks like in the "quaternionic" eigenvariety.


To prove these statements we need to formalise a lot of mathematics from the ground up. This blueprint
will consist of numerous distinct sections that we need to ultimately formalise our goal. A rough
outline of what this looks like is given below:

%% I want to generate a flow chart of what the dependencies look like... for now just writing them as a list
%% This could be generated by my blueprint eventually... but think I want a basic version first

\begin{enumerate}
  \item Restricted power series and the gauss norm
  \item Weierstrass preparation
  \item Newton polygons
  \item Nonarchimedean compact operators and Fredholm theory
  \item Quaternionic modular forms and Fujisakis lemma
  \item Hecke operators
  \item the main result
\end{enumerate}


%%FILE blueprint/src/chapter/NewtonPolygons.tex
\chapter{Newton Polygons}
\label{chap:np}

The Newton polygon of a power series over a nonarchimedean field records, in a single convex
picture, the valuations of its zeros.  It is the tool that converts a statement about a
characteristic power series into a statement about eigenvalues, and it is therefore the
combinatorial heart of what follows.

There are two ways to define it and we need both.  The first is geometric: the polygon is the
lower boundary of the convex hull of the points $(i, v(a_i))$.  That description is the one that
makes theorems true, but it is not something one can compute with directly.  The second is
algorithmic: walk from the leftmost finite point, repeatedly taking the segment of least slope.
That description computes, but proving anything about it means reasoning about the walk.

We formalise both and prove they agree.  The geometric description becomes a
\emph{specification}, a predicate $\mathrm{IsNewtonPolygonOf}$ that a candidate polygon may
satisfy; the algorithm produces a candidate; and existence plus uniqueness say the algorithm's
output is the unique polygon meeting the specification.  Every later theorem is then proved
against whichever description is convenient.

\section{One-sided polygons and the specification}

\begin{definition}
  \label{def:newtonpolygon0}
  \lean{NewtonPolygon₀}
  \leanok
  A \emph{one-sided polygon} over an ordered value group $\Gamma$ consists of a starting vertex
  in $\Z \times \Gamma$, a sequence of segment lengths, and a monotone sequence of slopes,
  together with the bookkeeping that says where the polygon stops (it may be finite, may end in a
  ray, or may continue forever).  Monotonicity of the slopes is convexity, built into the
  structure rather than imposed afterwards.
\end{definition}

Given a sequence $v : \N \to \Gamma \cup \{\infty\}$ — think $v_i = v(a_i)$, with $\infty$ where
the coefficient vanishes — we ask which polygon deserves to be called \emph{the} Newton polygon of
$v$.

\begin{definition}
  \label{def:isnewtonpolygonof}
  \uses{def:newtonpolygon0}
  \lean{IsNewtonPolygonOf}
  \leanok
  A one-sided polygon $P$ is \emph{the Newton polygon of} $v$ when
  \begin{enumerate}
    \item there are no points of $v$ strictly left of the starting vertex, and the starting
      vertex is itself a point of $v$;
    \item $P$ lies on or below every point: $P(k) \leq v_k$ for all $k \in \N$; and
    \item $P$ is the greatest such: any polygon $Q$ with the same starting abscissa lying on or
      below every point lies on or below $P$.
  \end{enumerate}
\end{definition}

Conditions (ii) and (iii) say ``greatest convex minorant through the anchor'', which is the lower
convex hull; convexity is already in Definition~\ref{def:newtonpolygon0}, so it does not need
restating.  The anchoring in (iii) is not cosmetic: a competitor starting further left is
unconstrained at indices where $v$ is infinite and could stand arbitrarily high there, so
maximality over unanchored competitors is unsatisfiable.

\begin{theorem}[Uniqueness]
  \label{thm:np-unique}
  \uses{def:isnewtonpolygonof}
  \lean{IsNewtonPolygonOf.height_eq}
  \leanok
  Two polygons satisfying Definition~\ref{def:isnewtonpolygonof} for the same $v$ have the same
  height at every integer.
\end{theorem}

\begin{proof}
  \leanok
  Each is below the other by (iii), applied in both directions.
\end{proof}

%% Note the shape of uniqueness: the NewtonPolygon₀ *structures* need not be equal, since a
%% segment may be split into two collinear pieces. Height equality is the geometric content.

\section{Existence: the algorithm meets the specification}

Not every sequence has a Newton polygon.  Two degeneracies must be excluded: if $v$ is
identically infinite there are no points at all, and if the slopes out of the anchor are
unbounded below then no polygon lies below every point.

\begin{definition}
  \label{def:isadmissible}
  \lean{IsAdmissible}
  \leanok
  $v$ is \emph{admissible} if for every finite point of $v$ the set of slopes from that point to
  later points is bounded below.
\end{definition}

\begin{lemma}
  \label{lem:isadmissible-affine}
  \uses{def:isadmissible}
  \lean{isAdmissible_of_affine_bound}
  \leanok
  If the points of $v$ lie on or above a single affine line then $v$ is admissible.
\end{lemma}

\begin{theorem}[Existence]
  \label{thm:np-exists}
  \uses{def:isnewtonpolygonof, def:isadmissible}
  \lean{isNewtonPolygonOf_newtonPolygon₀OfSeq}
  \leanok
  If $v$ has at least one finite value and is admissible, the polygon $\mathrm{newtonPolygon}_0(v)$
  produced by the step algorithm satisfies Definition~\ref{def:isnewtonpolygonof}.
\end{theorem}

Together with Theorem~\ref{thm:np-unique} this justifies speaking of \emph{the} Newton polygon,
and lets us prove a property of it by exhibiting any polygon with that property and checking the
specification --- a technique we use twice below.

\section{The polygon of a power series}

To pass from sequences to series we need the coefficient valuation.  We use the unnormalised one:

\begin{definition}
  \label{def:coeffval}
  \lean{coeffVal}
  \leanok
  For $f = \sum a_i T^i$ over an ultrametric field, $\mathrm{coeffVal}(f)_i = -\log \lVert a_i
  \rVert$, with the value $\infty$ where $a_i = 0$.  The Newton polygon of $f$ is the Newton
  polygon of this sequence.
\end{definition}

The base is $e$, so norms are recovered as $\lVert a \rVert = \exp(-\mathrm{coeffVal}(f)_i)$.
Rescaling to a chosen uniformizer $\varpi$ divides every slope by $-\log\lVert \varpi \rVert$;
we do that only when comparing with a source that normalises $v(\varpi) = 1$.

\begin{lemma}
  \label{lem:np-restricted}
  \uses{def:coeffval, thm:np-exists, lem:isadmissible-affine, def:isrestricted}
  \lean{isAdmissible_coeffVal_of_isRestricted, isNewtonPolygonOf_coeffVal_of_isRestricted}
  \leanok
  A restricted power series of parameter $c > 0$ has admissible coefficient valuations --- the
  Gauss bound $\lVert a_k \rVert c^k \leq B$ \emph{is} an affine floor --- hence has a Newton
  polygon.  The same holds for a nonzero polynomial.
\end{lemma}

\section{Prescribed slopes}

The following construction is the ``exhibit a witness and check the specification'' technique in
its purest form, and it is what we use to read slopes off a computed valuation sequence.

\begin{definition}
  \label{def:ofslopes}
  \uses{def:newtonpolygon0}
  \lean{NewtonPolygon₀.ofSlopes}
  \leanok
  For a monotone sequence of reals $s : \N \to \R$ and a starting height $y_0$, let
  $\mathrm{ofSlopes}(s, y_0)$ be the polygon anchored at $(0, y_0)$ with infinitely many segments
  of length $1$, the $j$-th having slope $s_j$.
\end{definition}

\begin{theorem}
  \label{thm:ofslopes}
  \uses{def:ofslopes, def:isnewtonpolygonof}
  \lean{isNewtonPolygonOf_ofSlopes}
  \leanok
  If $v_m = y_0 + \sum_{i < m} s_i$ for every $m$, then $\mathrm{ofSlopes}(s, y_0)$ is the Newton
  polygon of $v$.
\end{theorem}

\begin{proof}
  \leanok
  \uses{def:isnewtonpolygonof}
  The polygon passes through every point, so condition (ii) holds with equality; and any
  competitor lying below the points lies below the polygon at every integer, which is where
  heights are compared.  Admissibility is Lemma~\ref{lem:isadmissible-affine} through the minimum
  slope $s_0$.
\end{proof}

\begin{corollary}
  \label{cor:ofslopes-transport}
  \uses{thm:ofslopes, thm:np-unique, thm:np-exists}
  \lean{unitSlope_newtonPolygon₀OfSeq_ofSlopes}
  \leanok
  Consequently the polygon built by the algorithm from such a $v$ has unit slopes exactly $s_j$.
\end{corollary}

\section{First breaks, purity, and the link to preparation}

The first segment of the polygon controls the smallest zeros, and it is the object that connects
this chapter to Chapter~\ref{chap:wp}.

\begin{definition}
  \label{def:hasfirstbreak}
  \uses{def:coeffval}
  \lean{HasFirstBreak, IsPureSeries}
  \leanok
  $f$ \emph{has first break $(i, m)$} when the first segment of its polygon has slope $m$ and
  ends at abscissa $i$.  $f$ is \emph{pure of slope $m$} when its polygon is a single segment of
  slope $m$.
\end{definition}

\begin{proposition}
  \label{prop:firstbreak-distinguished}
  \uses{def:hasfirstbreak, def:ismuldistinguished}
  \lean{isMulDistinguished_of_hasFirstBreak}
  \leanok
  A first break at $(i, m)$ makes $f$ distinguished of order $i$ at the radius $e^{m}$, in the
  sense of Definition~\ref{def:ismuldistinguished}.
\end{proposition}

This is the bridge: the polygon locates a distinguished degree, Weierstrass preparation then
factors the series there, and the factor is a polynomial whose roots are the zeros of $f$ on the
corresponding disc.

\begin{theorem}[Weierstrass factorisation along the polygon]
  \label{thm:np-factorisation}
  \uses{prop:firstbreak-distinguished, thm:wprep}
  \lean{exists_weierstrass_factorisation}
  \leanok
  A power series over a complete field factors, along each break of its polygon, as a polynomial
  of degree the break abscissa times a unit of the corresponding restricted algebra.
\end{theorem}

\section{Zeros and slopes}

\begin{theorem}
  \label{thm:np-zeros}
  \uses{thm:np-factorisation}
  \lean{hasSum_zero_iff_aeval_eq_zero}
  \leanok
  A point of the closed disc is a zero of the (convergent) series exactly when it is a root of
  the polynomial factor supplied by Theorem~\ref{thm:np-factorisation}.
\end{theorem}

\begin{theorem}[The slope reading]
  \label{thm:np-slope-reading}
  \uses{thm:np-zeros}
  \lean{norm_eq_exp_slope_of_hasSum_zero}
  \leanok
  If $z$ is a zero of $f$ then $\lVert z \rVert = \exp(m)$ for $m$ a slope of the Newton polygon
  of $f$; and each slope $m$ occurs for exactly as many zeros, with multiplicity, as the
  horizontal length of the corresponding segment.
\end{theorem}

This is the statement the rest of the thesis consumes.  Applied to a Fredholm determinant
$\det(1 - Tu)$, whose zeros are the reciprocal eigenvalues of $u$, it converts the polygon of the
determinant into the list of valuations of the eigenvalues of $u$; that is the subject of
Chapter~\ref{chap:slopes}.

\begin{proposition}
  \label{prop:np-radius}
  \uses{def:coeffval, def:isrestricted}
  \lean{radiusOfConvergence, isRestricted_of_lt_radiusOfConvergence}
  \leanok
  The radius of convergence of $f$ is the supremum of the radii at which $f$ is restricted, and
  $f$ is restricted at every smaller radius.  In polygon terms it is $\exp$ of the supremum of
  the slopes.
\end{proposition}

\section{The polygon of a product}

Newton's own observation is that the polygon of a product is the \emph{Minkowski sum} of the two
polygons: sort the segments of both by slope and concatenate.  We prove it in the height form the
rest of the development reads, and in the counting form that says slope multiplicities add.

Two auxiliary notions carry the proof.  The first is the \emph{face} of a given slope: the
(possibly degenerate) segment where the supporting line of that slope touches the polygon.  The
second is a finiteness condition without which the product formula is simply false.

\begin{definition}
  \label{def:slopesunbounded}
  \uses{def:newtonpolygon0}
  \lean{NewtonPolygon₀.SlopesUnbounded}
  \leanok
  A polygon has \emph{unbounded slopes} if every real number is exceeded by one of its unit
  slopes, the junk value $\infty$ past a finite polygon counting.  Polynomials qualify through
  that junk tail, and entire series because their coefficients lie above a line of every slope.
\end{definition}

The hypothesis cannot be dropped: $(1 - X)\sum_i X^i = 1$, where the two factors have polygons
whose Minkowski sum is a horizontal ray while the product's polygon is a single point.  It is the
polygon-level form of the convergence hypothesis Koblitz imposes for the same reason.

\begin{definition}
  \label{def:face}
  \uses{def:slopesunbounded}
  \lean{NewtonPolygon₀.faceLeft, NewtonPolygon₀.faceRight}
  \leanok
  For a real $\sigma$, $\mathrm{faceLeft}(\sigma)$ is the number of unit slopes strictly below
  $\sigma$ and $\mathrm{faceRight}(\sigma)$ the number at most $\sigma$.  These are the two
  endpoints of the face of slope $\sigma$; the polygon is the line of slope $\sigma$ between them
  and strictly above it outside.
\end{definition}

\begin{proposition}
  \label{prop:vertex-is-point}
  \uses{def:isnewtonpolygonof, def:face}
  \lean{IsNewtonPolygonOf.height_eq_pointHeight_of_unitSlope_lt}
  \leanok
  Where the unit slope strictly increases, the polygon meets the point sequence.  In particular
  both endpoints of every face are points.  (Textbooks take this as the definition of a vertex;
  against the specification it is a theorem, proved by raising a supporting line of an
  intermediate slope and contradicting maximality.)
\end{proposition}

\begin{definition}
  \label{def:minkowskiheight}
  \uses{def:newtonpolygon0}
  \lean{NewtonPolygon₀.minkowskiHeight}
  \leanok
  The height at $n$ of the Minkowski sum of two polygons is
  $\min_{i + j = n}\bigl(P(i) + Q(j)\bigr)$, the infimal convolution of the two height functions.
\end{definition}

\begin{theorem}[The polygon of a product]
  \label{thm:np-product}
  \uses{def:minkowskiheight, def:face, prop:vertex-is-point, def:coeffval}
  \lean{IsEntireNewtonPolygonOf.height_mul}
  \leanok
  For $f$ and $g$ with nonzero constant coefficient and unbounded slopes, the Newton polygon of
  $fg$ is the Minkowski sum of the polygons of $f$ and $g$.
\end{theorem}

\begin{proof}
  \leanok
  \uses{prop:vertex-is-point}
  Both inequalities come from the face of a common slope.  The ultrametric bound puts every point
  of $fg$ on or above the Minkowski sum, since $c_n = \sum_{i+j=n} a_i b_j$.  A minimising split
  of $n$ admits a slope $\sigma$ supporting both factors there, so the sum of the two supporting
  lines lies below every point of $fg$ and hence below its polygon; that is one inequality.  For
  the other, at the two endpoints of the face of slope $\sigma$ the minimising split is
  \emph{unique}, so the ultrametric bound is an equality there and the endpoints are genuine
  points of $fg$.  The polygon of $fg$ is therefore below the chord joining them, and that chord
  is exactly the Minkowski sum on the face.
\end{proof}

\begin{corollary}[Multiplicities add]
  \label{cor:np-product-multiplicity}
  \uses{thm:np-product}
  \lean{IsEntireNewtonPolygonOf.faceRight_mul, IsEntireNewtonPolygonOf.faceLeft_mul}
  \leanok
  The number of slopes of $fg$ that are at most $\sigma$ is the sum of the corresponding counts
  for $f$ and for $g$; likewise for the slopes strictly below $\sigma$.  Equivalently, the slope
  multiset of $fg$ is the union of the slope multisets of $f$ and $g$.
\end{corollary}

\begin{corollary}[The initial segment]
  \label{cor:np-product-initial}
  \uses{thm:np-product}
  \lean{IsEntireNewtonPolygonOf.height_mul_of_forall_le,
    height_newtonPolygon₀OfPowerSeries_mul_coe}
  \leanok
  If the first $n$ slopes of $g$ are at most every slope of $f$, then the polygon of $fg$ agrees
  with that of $g$ up to $n$.  Applied to a Fredholm determinant factored as a polynomial times
  the determinant of the complement, this says the finite factor supplies the initial segment of
  the polygon --- the form used in Chapter~\ref{chap:slopes}.
\end{corollary}


%%FILE blueprint/src/chapter/QMF.tex
\chapter{Quaternionic Modular Forms}
\label{chap:qmf}

Let $F$ be a totally real number field and $D$ a totally definite quaternion algebra over $F$,
split at a chosen finite place $v$.  Because $D$ is definite, $D^\times$ is discrete and
cocompact in the finite adelic points, and the ``modular forms'' attached to $D$ are not
functions on a symmetric space but functions on a finite set --- a double coset space.  That is
what makes them tractable: everything analytic in the classical theory is replaced by finite
linear algebra over the coefficients, and the only genuine analysis left is in the coefficients
themselves.

This chapter defines the forms at a general locally analytic weight, constructs the Hecke
operators, proves that the $U_\varpi$-type operators are compact, and computes their Fredholm
determinants.  The pattern throughout is: define the object adelically, transport it to a block
model through evaluation at class-set representatives, and do the analysis in the model.

\section{The adelic setting}

\begin{definition}
  \label{def:dfx}
  \lean{QMF.Dfx, QMF.globalUnits}
  \leanok
  Write $D_f = D \otimes_F \A_F^\infty$ and $D_f^\times$ for its unit group.  The subgroup
  $\Gamma = \mathrm{globalUnits}$ is the image of $D^\times$ in $D_f^\times$.
\end{definition}

\begin{definition}
  \label{def:rigidification}
  \uses{def:dfx}
  \lean{QMF.RigidificationAt, QMF.toLocal, QMF.toMatrix}
  \leanok
  A \emph{rigidification at $v$} is a fixed isomorphism $D \otimes_F F_v \cong M_2(F_v)$; it
  exists exactly when $D$ splits at $v$.  Composing the local component map $\mathrm{toLocal}_v :
  D_f \to D \otimes_F F_v$ with it gives $\mathrm{toMatrix}_v : D_f^\times \to M_2(F_v)$.
\end{definition}

\begin{definition}
  \label{def:sigma0}
  \lean{Sigma0'}
  \leanok
  For $\gamma < 1$ in the value group let
  \[
    \Sigma_0(\gamma) = \left\{ g = \begin{pmatrix} a & b \\ c & d\end{pmatrix} \in M_2(F_v) \;:\;
    v(g_{ij}) \leq 1, \; v(c) \leq \gamma, \; v(d) = 1, \; \det g \neq 0 \right\},
  \]
  the matrices that are integral, congruent to upper-triangular modulo $\gamma$, and with unit
  lower-right entry.  A subgroup $U \subseteq D_f^\times$ has \emph{wild level $\geq \gamma$} if
  $\mathrm{toMatrix}_v(U) \subseteq \Sigma_0(\gamma)$; this is Buzzard's condition
  \cite[\S9]{Buzzard}, and $\Delta = \mathrm{toMatrix}_v^{-1}(\Sigma_0(\gamma))$ is the monoid on
  which the weight will act.
\end{definition}

\section{Weights}

A weight is a locally analytic character $\kappa$, and the object it produces is an action of
$\Sigma_0(\gamma)$ on the Tate algebra:
\[
  (f \|_\kappa \gamma)(z) = \frac{\kappa(cz+d)}{(cz+d)^2}\, f\!\left(\frac{az+b}{cz+d}\right).
\]
Making this precise requires knowing that $\kappa(cz+d)$ expands as a power series in $z$ with
suitable coefficient bounds.  Rather than assume the expansion exists abstractly --- which would
force a choice, and a choice inside a definition propagates into every downstream statement ---
we carry it as data.

\begin{definition}
  \label{def:levelbounds}
  \uses{def:sigma0}
  \lean{QMF.LevelBounds}
  \leanok
  $\mathrm{LevelBounds}\, S\, \rho$ records the norm estimates extracted from an acting monoid
  $S$: entries integral, $\lVert c \rVert \leq \rho$, and $\lVert d \rVert = 1$.
\end{definition}

\begin{definition}
  \label{def:weightseries}
  \uses{def:levelbounds}
  \lean{QMF.WeightSeries}
  \leanok
  A \emph{weight datum} $W$ at $(S, \rho)$ consists of the level bounds, the expansion
  $\kappa\mathrm{col}(c,d) \in K[[z]]$ of $\kappa(cz+d)$ as a function of the lower row, the
  decay estimate $\lVert \mathrm{coeff}_m \rVert \leq \rho^m$, absolute summability, and the
  cocycle identity making $\|_\kappa$ an action.
\end{definition}

\begin{definition}
  \label{def:analyticweight}
  \uses{def:weightseries}
  \lean{QMF.ExpansionData, QMF.AnalyticWeight}
  \leanok
  An \emph{analytic weight} is a genuine character $\kappa : U_K \to K^\times$ together with an
  expansion datum for it.  The expansion is unique on the level, so the pair is no more data than
  the character; but carrying it makes the weight action a definition rather than a choice.
\end{definition}

\begin{definition}
  \label{def:kappaslash}
  \uses{def:analyticweight}
  \lean{QMF.AnalyticWeight.kappaSlash}
  \leanok
  $\kappa\text{-}\mathrm{slash}(g)$ is the continuous $K$-linear operator on the Tate algebra
  determined by $z^k \mapsto \kappa(cz+d)(cz+d)^{-2}\left(\frac{az+b}{cz+d}\right)^k$.
\end{definition}

\begin{proposition}
  \label{prop:kappaslash-matrix}
  \uses{def:kappaslash, prop:tate-algebra-onable, def:matrixcoeff}
  \lean{QMF.WeightSeries.matrixCoeff_kappaSlash, QMF.WeightSeries.norm_matrixCoeff_kappaSlash_le}
  \leanok
  Under the identification of Proposition~\ref{prop:tate-algebra-onable}, the matrix of
  $\kappa\text{-}\mathrm{slash}(g)$ is the coefficient array of the two-variable generating
  function $\kappa\mathrm{col}\cdot(cx+d)^{-1}(cx+d-axy-by)^{-1}$; and if $\lVert a \rVert \leq
  \sigma$ with $\rho \leq \sigma$ then
  \[
    \lVert \mathrm{matrixCoeff}(\kappa\text{-}\mathrm{slash}(g))_{j i} \rVert \leq \sigma^{\,j}.
  \]
\end{proposition}

The row-decay estimate is [Jacobs, Lemma 2.7]; it is where the analysis of the weight is spent,
and everything after it is formal.

\section{The forms}

\begin{definition}
  \label{def:forms}
  \uses{def:analyticweight, def:sigma0}
  \lean{QMF.Weight.levelMonoidOf, QMF.Weight.Forms}
  \leanok
  Let $\theta$ be the component map, $S$ an acting monoid, $\kappa$ an analytic weight, $U$ a
  subgroup with $\theta(U) \subseteq S$, and $\chi$ a twisting character.  Then
  \[
    S_\kappa(U) = \left\{ \varphi : D_f^\times \to K\langle z\rangle \;:\; \varphi(dgu) =
    \chi(u_v)\, \varphi(g) \|_\kappa u_v \ \ \forall d \in \Gamma,\ u \in U \right\}.
  \]
\end{definition}

\begin{definition}
  \label{def:heckeoperator}
  \uses{def:forms}
  \lean{AbstractHeckeOperatorSlash.heckeOperatorSlash, QMF.Weight.heckeOperator}
  \leanok
  For $\eta$ in the acting monoid, write $U \eta U = \coprod_i U x_i$ and set $[U\eta U]\varphi =
  \sum_i \varphi \|_\kappa x_i$.  This is independent of the representatives and preserves
  $S_\kappa(U)$.
\end{definition}

The definition presupposes that the double coset is a finite union.  That is automatic at a
compact open level:

\begin{theorem}
  \label{thm:hecke-finiteness}
  \uses{def:heckeoperator}
  \lean{AbstractHeckeOperatorSlash.finite_image_doubleCoset_of_isOpen_of_isCompact}
  \leanok
  If $U$ is open and compact then $U \eta U$ is a finite union of right cosets of $U$.
\end{theorem}

\begin{proposition}
  \label{prop:compact-open-levels}
  \uses{thm:hecke-finiteness}
  \lean{Units.isCompact_of_isCompact, QMF.isCompact_integralTensor, QMF.continuous_toMatrix}
  \leanok
  The standard levels are compact open.  Concretely: the units of a compact multiplicatively
  closed set form a compact set of units, because $u \mapsto (u, u^{-1})$ is a closed embedding;
  and the integral order $\calO_D \otimes \Zhat \subseteq D_f$ is compact, being the image of a
  product of compact local integer rings under a continuous coordinate map.
\end{proposition}

\begin{theorem}[Fujisaki]
  \label{thm:fujisaki}
  \uses{def:dfx}
  \lean{QMF.finite_classSet}
  \leanok
  For $U$ open the class set $\Gamma \backslash D_f^\times / U$ is finite.
\end{theorem}

\section{The block model}

Fix representatives $c_1, \dots, c_n$ of the class set and let $\Gamma_i = c_i^{-1}\Gamma c_i
\cap U$ be the stabilisers.  A form is determined by its values at the $c_i$, and those values
are $\Gamma_i$-invariant.

\begin{definition}
  \label{def:evalatreps}
  \uses{def:forms, thm:fujisaki}
  \lean{QMF.Weight.evalAtReps}
  \leanok
  $\mathrm{evalAtReps}(\varphi) = (\varphi(c_i))_i$, a map $S_\kappa(U) \to c(\iota \times \N,
  K)$.
\end{definition}

\begin{theorem}
  \label{thm:evalatreps}
  \uses{def:evalatreps}
  \lean{QMF.Weight.evalAtReps_injective, QMF.Weight.mem_range_evalAtReps_iff,
        QMF.Weight.bijective_evalAtReps}
  \leanok
  $\mathrm{evalAtReps}$ is injective, and its image is exactly the tuples whose $i$-th block is
  $\Gamma_i$-invariant.  If the level is \emph{neat} --- every $\Gamma_i$ acts trivially on the
  Tate algebra --- it is bijective, so $S_\kappa(U) \cong c(\iota \times \N, K)$.
\end{theorem}

Neatness is [Jacobs, Lemma 2.2] and holds in the case of interest.  Without it the forms are
still a direct summand of the model, which is property (Pr); see \S\ref{sec:pr}.

\begin{definition}
  \label{def:heckeblock}
  \uses{def:evalatreps, def:kappaslash}
  \lean{QMF.Weight.heckeBlock, QMF.Weight.heckeBlockOp}
  \leanok
  Given right coset representatives $v_t$ of $U\eta U$ and, for each $(i,t)$, a factorisation
  $c_i v_t^{-1} = d_{i,t}\, c_{\sigma(i,t)}\, u_{i,t}$ with $d_{i,t} \in \Gamma$, $u_{i,t} \in
  U$, the $(i,j)$ \emph{block} is
  \[
    \varepsilon_{ij} = \sum_{t \,:\, \sigma(i,t) = j} \chi(w_{i,t})\,
    \kappa\text{-}\mathrm{slash}(w_{i,t}), \qquad w_{i,t} = \theta(u_{i,t} v_t),
  \]
  and $\mathrm{heckeBlockOp}$ is the operator on $c(\iota \times \N, K)$ they assemble.
\end{definition}

\begin{theorem}[The transport]
  \label{thm:transport}
  \uses{def:heckeblock, def:heckeoperator}
  \lean{QMF.Weight.heckeOperator_apply_rep, QMF.Weight.evalAtReps_heckeOperator}
  \leanok
  $\bigl([U\eta U]\varphi\bigr)(c_i) = \sum_j \varepsilon_{ij}\bigl(\varphi(c_j)\bigr)$; that is,
  $\mathrm{evalAtReps}$ intertwines $[U\eta U]$ with $\mathrm{heckeBlockOp}$.
\end{theorem}

\section{Compactness and the Fredholm determinant}

\begin{theorem}[{\cite[Lemma 12.2]{Buzzard}}, {[Jacobs, Lemma 2.7]}]
  \label{thm:qmf-compact}
  \uses{thm:transport, prop:kappaslash-matrix, prop:compactoid-criterion}
  \lean{QMF.Weight.isCompactoid_heckeBlock, QMF.Weight.isCompactoid_heckeBlockOp}
  \leanok
  Suppose $\lVert \det \theta(\eta) \rVert \leq \sigma$ with $\rho \leq \sigma < 1$.  Then every
  block, and hence the block operator, is compactoid.
\end{theorem}

\begin{proof}
  \leanok
  \uses{prop:kappaslash-matrix}
  Each certificate element $u_{i,t}v_t$ has determinant norm at most that of $\eta$, the outer
  factors being integral; a small determinant with integral entries and unit $d$ forces
  $\lVert a \rVert \leq \sigma$; and then Proposition~\ref{prop:kappaslash-matrix} gives the row
  decay $\sigma^j$, which is Proposition~\ref{prop:compactoid-criterion}'s hypothesis.  A finite
  sum of compactoid operators is compactoid.
\end{proof}

\begin{definition}
  \label{def:heckecharpowerseries}
  \uses{thm:qmf-compact, def:minor}
  \lean{QMF.Weight.heckeCharPowerSeries}
  \leanok
  $\det(1 - T\,[U\eta U])$, computed in the block model.
\end{definition}

\begin{theorem}[Eigenforms are reciprocal roots]
  \label{thm:qmf-eigenform-criterion}
  \uses{def:heckecharpowerseries, thm:serre12, thm:evalatreps}
  \lean{QMF.Weight.evalT_heckeCharPowerSeries_eq_zero_iff}
  \leanok
  At a neat level, $a \neq 0$ is a zero of $\det(1 - T[U\eta U])$ if and only if $a^{-1}$ is an
  eigenvalue of $[U\eta U]$ on $S_\kappa(U)$.
\end{theorem}

\begin{theorem}
  \label{thm:qmf-riesz}
  \uses{thm:qmf-eigenform-criterion, thm:riesz}
  \lean{QMF.Weight.exists_riesz_decomposition_forms}
  \leanok
  Each zero $a$ of the determinant, of multiplicity $n$, gives a decomposition $S_\kappa(U) = N
  \oplus F$ into $[U\eta U]$-stable summands with $(1 - a[U\eta U])^n = 0$ on $N$, invertible on
  $F$, and $\dim_K N = n$.
\end{theorem}

\section{Non-neat levels}
\label{sec:pr}

\begin{proposition}
  \label{prop:pr}
  \uses{thm:evalatreps, def:onable}
  \lean{QMF.Weight.stabProj, QMF.Weight.heckeCharPowerSeriesPr}
  \leanok
  Without neatness, averaging over the finite stabilisers gives projectors realising
  $S_\kappa(U)$ as a direct summand of the model, so it satisfies property (Pr).  The
  determinant of $[U\eta U]$ on the summand is well defined and independent of the projector
  chosen \cite[Lemma 2.12]{Buzzard}.
\end{proposition}

\section{Base change and the Hecke algebra}

\begin{proposition}
  \label{prop:qmf-basechange}
  \uses{def:heckecharpowerseries}
  \lean{QMF.Weight.heckeCharPowerSeries_map}
  \leanok
  If $\kappa'$ over $L$ is compatible with $\kappa$ over $K$ along an isometric embedding $K \to
  L$, then the two determinants correspond \cite[Lemma 2.13]{Buzzard}.  This is what lets a
  computation made over a small field be read over $\bbC_p$.
\end{proposition}

Following Buzzard \cite[\S5]{Buzzard} we take the commuting structure of the Hecke algebra as
data rather than proving it; for $T_w, T_{w'}$ at distinct places it is a Satake computation
outside our scope.  What we prove is what the eigenvariety machine consumes from that datum.

\begin{proposition}
  \label{prop:hecke-comm-criterion}
  \uses{def:heckeoperator}
  \lean{AutomorphicFunction.heckeOperatorSlash_comm_of_reps}
  \leanok
  Two double coset operators commute as soon as their representative families can be matched
  pairwise, up to the level, with the two products taken in the opposite orders.
\end{proposition}

%% The naive "disjoint support implies commuting" criterion is false here: representatives are
%% η·u with a global u, so they do not commute elementwise.

\begin{proposition}
  \label{prop:qmf-eigensystem}
  \uses{thm:qmf-riesz, prop:hecke-comm-criterion}
  \lean{QMF.Weight.mapsTo_ker_of_commute, QMF.Weight.exists_eigensystem_of_riesz}
  \leanok
  Let $\mathbb{T}$ be a commuting family of endomorphisms of $S_\kappa(U)$, each commuting with $\varphi
  = [U\eta U]$, and let $N$ be the finite-slope subspace of Theorem~\ref{thm:qmf-riesz}.  Then
  every element of $\mathbb{T}$ preserves $N$, and over an algebraically closed coefficient field $N$
  contains a nonzero simultaneous generalised eigenvector: a $\varphi$-finite system of
  eigenvalues.
\end{proposition}

\begin{proof}
  \leanok
  \uses{thm:qmf-riesz}
  Stability is the identity $(1-a\varphi)^n t = t (1-a\varphi)^n$ applied to a vector killed by
  the left factor.  $N$ is finite dimensional by Theorem~\ref{thm:qmf-riesz} and nonzero, and a
  commuting family of endomorphisms of a nonzero finite-dimensional space over an algebraically
  closed field is simultaneously triangularisable.
\end{proof}


%%FILE blueprint/src/chapter/RestrictedPowerSeries.tex
\chapter{Restricted Power Series and the Gauss Norm}
\label{chap:rps}

Everything in this thesis takes place inside rings of power series that converge on a closed
disc.  Over $\bbC$ one would speak of holomorphic functions and take convergence for granted;
in the nonarchimedean world convergence of $\sum a_t x^t$ on the polydisc of polyradius $c$ is
simply the condition
\[
  \lVert a_t \rVert c^t \longrightarrow 0,
\]
because a nonarchimedean series converges exactly when its terms tend to $0$.  A power series
satisfying this is called \emph{restricted} for $c$, and the collection of them is the ring we
work in.  The supremum $\sup_t \lVert a_t\rVert c^t$ is then finite and is the \emph{Gauss
norm}, which makes the ring into a Banach ring.

Two design decisions shape this chapter and are worth stating at the outset.

First, we work multivariate from the start and treat the one-variable case as the specialisation
to a one-element index type.  This is not generality for its own sake: the multivariate
Weierstrass theory of Chapter~\ref{chap:wp} is proved by splitting off one variable and applying
the one-variable theory \emph{over the ring of restricted series in the remaining variables}, so
the two cases must be literally the same definition rather than parallel developments.

Second, the radius $c$ is an arbitrary real parameter, not $1$.  Fixing $c = 1$ would give the
Tate algebra and nothing else; but the Newton polygon of Chapter~\ref{chap:np} reads a series at
\emph{every} radius, and the Weierstrass theorems have to hold at each of them.

Part of what follows is already in mathlib --- the predicate $\mathrm{IsRestricted}$ and the
Gauss norm as a bare function on $\mathrm{MvPowerSeries}$ --- and we say so where that is the
case.  What this project adds is the ring, its norm, its completeness, and the structure theory
of its units and power-bounded elements.

\section{Restricted power series}

\begin{definition}
  \label{def:isrestricted}
  \lean{MvPowerSeries.IsRestricted}
  \leanok
  Let $R$ be a normed ring, $\sigma$ an index type and $c : \sigma \to \R$.  A power series
  $f = \sum_t a_t x^t \in R[[x_i : i \in \sigma]]$ is \emph{restricted for $c$} if
  \[
    \lVert a_t \rVert \prod_{i} c_i^{\,t_i} \longrightarrow 0
  \]
  along the cofinite filter on the multidegrees $t$.
\end{definition}

Convergence along the cofinite filter, rather than along a sequence, is the right formulation:
the index set $\sigma \to_0 \N$ has no preferred enumeration, and the cofinite filter says
precisely ``all but finitely many terms are small'', which is what an unordered sum needs.  It
is also what mathlib's summability API consumes, so no bespoke analysis is required.

\begin{lemma}
  \label{lem:isrestricted-abs}
  \uses{def:isrestricted}
  \lean{MvPowerSeries.isRestricted_abs_iff}
  \leanok
  $f$ is restricted for $c$ if and only if it is restricted for $\lvert c \rvert$.
\end{lemma}

Some sources define restricted series only for nonnegative radii.  Lemma~\ref{lem:isrestricted-abs}
says nothing is lost by allowing arbitrary real $c$, and keeping the definition unconstrained
avoids a positivity side condition in every downstream statement.  Positivity is imposed later,
where it is genuinely needed --- for the Gauss norm to be a norm rather than a seminorm.

\begin{lemma}
  \label{lem:isrestricted-finite-support}
  \uses{def:isrestricted}
  \lean{MvPowerSeries.isRestricted_of_finite_support, MvPolynomial.isRestricted_toMvPowerSeries}
  \leanok
  A power series with finitely many nonzero coefficients is restricted for every $c$; in
  particular every polynomial is.
\end{lemma}

\begin{proposition}
  \label{prop:restricted-ring}
  \uses{def:isrestricted}
  \lean{MvPowerSeries.Restricted}
  \leanok
  The restricted series for a fixed $c$ form a subring of $R[[x_i]]$.  We write $R_c\langle x_i :
  i \in \sigma\rangle$, or $\mathrm{Restricted}\ R\ c$, for it.
\end{proposition}

\begin{proof}
  \leanok
  \uses{lem:isrestricted-finite-support}
  Closure under subtraction is the usual squeeze.  Multiplication is the interesting axiom: the
  coefficients of a product are antidiagonal sums,
  \[
    (fg)_t = \sum_{u + w = t} a_u b_w ,
  \]
  and it is not immediate that these tend to $0$.  They do, because for each $t$ the
  antidiagonal is finite and the norm of the sum is bounded by the maximum of $\lVert a_u
  \rVert c^u \cdot \lVert b_w \rVert c^w$ over it; given $\varepsilon$, all but finitely many
  $u$ have $\lVert a_u\rVert c^u$ small and all but finitely many $w$ have $\lVert b_w \rVert
  c^w$ small, and a multidegree $t$ outside a suitable finite set cannot be split with both
  parts in the exceptional sets.
\end{proof}

%% This is where HasAntidiagonal matters: the argument needs the antidiagonal of t to be finite,
%% which is why N-indexed multidegrees work and Z-indexed ones (Laurent series) do not.

\begin{definition}
  \label{def:restricted-elements}
  \uses{prop:restricted-ring}
  \lean{MvPowerSeries.Restricted.monomial, MvPowerSeries.Restricted.C,
        MvPowerSeries.Restricted.X, MvPolynomial.toRestricted}
  \leanok
  Monomials, constants and variables are restricted, giving $\mathrm{monomial}$, $C$ and $X$ in
  $R_c\langle x\rangle$, and the inclusion $\mathrm{toRestricted}$ of the polynomial ring, which
  is an injective ring homomorphism.
\end{definition}

\begin{definition}
  \label{def:restricted-map}
  \uses{prop:restricted-ring}
  \lean{MvPowerSeries.Restricted.map}
  \leanok
  A norm-nonincreasing ring homomorphism $R \to S$ induces $R_c\langle x\rangle \to S_c\langle
  x\rangle$ coefficientwise.
\end{definition}

Functoriality in the coefficients is what makes base change of the Fredholm determinant work in
Chapter~\ref{chap:qmf}: the norm-nonincreasing hypothesis is exactly what preserves
restrictedness.

\begin{definition}
  \label{def:tate-algebra}
  \uses{prop:restricted-ring}
  \lean{PowerSeries.IsRestricted, PowerSeries.Restricted}
  \leanok
  The one-variable case is the specialisation of Definition~\ref{def:isrestricted} to a
  one-element index type.  When moreover $R = K$ is a complete nonarchimedean field and $c = 1$,
  the ring $K\langle x \rangle$ is the \emph{Tate algebra} $T_1$.
\end{definition}

%% In Lean, PowerSeries.Restricted R c is *by definition* MvPowerSeries.Restricted R (fun _ : Unit => c),
%% so every instance transfers rather than being reproved. The N-indexed lemmas are restatements.

\section{The Gauss norm}

\begin{definition}
  \label{def:gaussnorm}
  \uses{def:isrestricted}
  \lean{MvPowerSeries.gaussNorm, MvPowerSeries.HasGaussNorm}
  \leanok
  The \emph{Gauss norm} of $f$ at radius $c$ is
  \[
    \lVert f \rVert_c = \sup_t \; \lVert a_t \rVert \prod_i c_i^{\,t_i},
  \]
  when the supremum is finite; the predicate $\mathrm{HasGaussNorm}$ records finiteness.
\end{definition}

For a restricted series the supremum is automatically finite --- the terms tend to $0$, so all
but finitely many are below any given bound --- and it is moreover \emph{attained}:

\begin{definition}
  \label{def:achievesgaussnorm}
  \uses{def:gaussnorm}
  \lean{MvPowerSeries.AchievesGaussNorm}
  \leanok
  $f$ \emph{achieves} its Gauss norm at $t$ if $\lVert a_t \rVert c^t = \lVert f \rVert_c$.
\end{definition}

Attainment is not a technicality; it is the property that makes the Gauss norm behave
multiplicatively rather than merely submultiplicatively, and everything in Chapter~\ref{chap:wp}
rests on it.  The following is the master lemma of this chapter.

\begin{lemma}[Dominant terms on the antidiagonal]
  \label{lem:gaussnorm-dominant}
  \uses{def:achievesgaussnorm}
  \lean{MvPowerSeries.exists_achievesGaussNorm_dominant,
        MvPowerSeries.IsRestricted.exists_achievesGaussNorm_dominant}
  \leanok
  Let $f, g$ be restricted with nonzero Gauss norms.  Fix any well-order on $\sigma$ and let
  $u$, $w$ be the achieving multidegrees of $f$ and $g$ that are largest for the induced
  lexicographic order.  Then the term $a_u b_w$ strictly dominates every other term of the
  antidiagonal of $u + w$.
\end{lemma}

\begin{proof}
  \leanok
  Any other splitting $u' + w' = u + w$ has $u' \neq u$, so by maximality either $u'$ is not
  achieving --- and then its term is strictly smaller --- or $u'$ is achieving and
  lexicographically smaller, forcing $w'$ to be lexicographically larger than $w$ and hence
  non-achieving.  Either way one factor is strictly below its Gauss norm.
\end{proof}

\begin{theorem}
  \label{thm:gaussnorm-mul}
  \uses{lem:gaussnorm-dominant}
  \lean{MvPowerSeries.gaussNorm_mul_le, MvPowerSeries.gaussNorm_mul_eq_mul}
  \leanok
  The Gauss norm is always submultiplicative, and if the norm on $R$ is multiplicative and both
  series attain their Gauss norms, then
  \[
    \lVert fg \rVert_c = \lVert f \rVert_c \, \lVert g \rVert_c .
  \]
\end{theorem}

\begin{proof}
  \leanok
  \uses{lem:gaussnorm-dominant}
  Submultiplicativity is the ultrametric bound on the antidiagonal sum.  For equality, the
  coefficient of $u + w$ from Lemma~\ref{lem:gaussnorm-dominant} is an ultrametric sum with one
  strictly dominant term, so its norm is $\lVert a_u\rVert \lVert b_w \rVert c^{u+w}$, which is
  the product of the two Gauss norms.
\end{proof}

\begin{theorem}
  \label{thm:restricted-normedring}
  \uses{thm:gaussnorm-mul, prop:restricted-ring}
  \lean{MvPowerSeries.Restricted.gaussNorm, MvPowerSeries.Restricted.gaussNormRingNorm,
        MvPowerSeries.Restricted.isNonarchimedean_norm}
  \leanok
  For strictly positive radii the Gauss norm is a ring norm on $R_c\langle x \rangle$, and it is
  nonarchimedean whenever the norm on $R$ is.  If the norm on $R$ is multiplicative then so is
  the Gauss norm.
\end{theorem}

%% The NormedRing instance built from gaussNormRingNorm is the canonical source of the norm,
%% metric, uniformity and topology on Restricted R c; downstream structure is derived from it and
%% never rebuilt, which is what keeps the topologies definitionally equal.

\begin{lemma}
  \label{lem:gaussnorm-values}
  \uses{thm:restricted-normedring}
  \lean{MvPowerSeries.Restricted.norm_monomial, MvPowerSeries.Restricted.norm_C,
        MvPowerSeries.Restricted.norm_X, MvPolynomial.norm_toRestricted}
  \leanok
  $\lVert a x^t \rVert_c = \lVert a \rVert c^t$, $\lVert C(a) \rVert_c = \lVert a \rVert$ and
  $\lVert X_i \rVert_c = c_i$; and the norm of a polynomial, viewed in $R_c\langle x\rangle$, is
  its polynomial Gauss norm --- a finite maximum.
\end{lemma}

\begin{proposition}
  \label{prop:polynomial-gaussnorm}
  \uses{def:gaussnorm}
  \lean{Polynomial.gaussNorm_mul, Polynomial.gaussNorm_isAbsoluteValue}
  \leanok
  On polynomials over a field with multiplicative norm the Gauss norm is multiplicative, hence
  an absolute value.  This is Gauss's lemma in its normed form, and the reason the norm carries
  his name.
\end{proposition}

\section{Completeness}

\begin{theorem}
  \label{thm:restricted-complete}
  \uses{thm:restricted-normedring}
  \lean{MvPowerSeries.Restricted.cauchySeq_coeff,
        MvPowerSeries.Restricted.tendsto_of_eventually_norm_coeff_sub_le}
  \leanok
  If $R$ is complete then so is $R_c \langle x \rangle$: it is a Banach ring.
\end{theorem}

\begin{proof}
  \leanok
  A Cauchy sequence is coefficientwise Cauchy, since each coefficient functional is
  norm-nonincreasing, so the coefficients converge in $R$ and assemble into a candidate limit.
  The point requiring care is that the candidate is itself restricted: the Cauchy estimate is
  \emph{uniform} in the multidegree, so a tail of the sequence has all its coefficients within
  $\varepsilon$ of the limit's, and restrictedness of any one member transfers to the limit.
  Convergence in Gauss norm then follows from the same uniform estimate.
\end{proof}

Completeness is what Chapter~\ref{chap:wp} spends: the Weierstrass division algorithm produces a
contracting sequence of approximate quotients, and the limit exists only because this ring is
Banach.

\section{Power-bounded elements, units and reduction}

The last three sections of theory are what turn $R_c\langle x\rangle$ from a Banach ring into an
object of rigid geometry: its power-bounded subring, its units, and its reduction.

\begin{proposition}
  \label{prop:powerbounded}
  \uses{thm:restricted-normedring}
  \lean{MvPowerSeries.Restricted.isPowerBounded_iff_forall_isPowerBounded_coeff,
        MvPowerSeries.Restricted.isTopologicallyNilpotent_iff_forall_isTopologicallyNilpotent_coeff}
  \leanok
  At radius $c = 1$, a restricted power series is power bounded if and only if all its
  coefficients are, and topologically nilpotent if and only if all its coefficients are.
\end{proposition}

Both statements fail at general radii in one direction, and the file records the two halves
separately with their genuine hypotheses ($c \geq 1$ for one implication, $c \leq 1$ for the
other); the clean iff is the case $c = 1$.  We write $T^\circ$ for the power-bounded subring.

\begin{theorem}
  \label{thm:restricted-units}
  \uses{thm:restricted-complete, def:achievesgaussnorm}
  \lean{MvPowerSeries.Restricted.constantCoeff,
        MvPowerSeries.Restricted.isUnit_of_norm_lt_norm_constantCoeff,
        MvPowerSeries.Restricted.isUnit_iff}
  \leanok
  Let $R$ be complete.  A restricted power series $f$ is a unit if and only if its constant
  coefficient is a unit of $R$ and every other coefficient is strictly smaller in norm:
  the Gauss norm is achieved at the constant term alone.
\end{theorem}

\begin{proof}
  \leanok
  \uses{thm:restricted-complete}
  If $\lVert f - C(f_0)\rVert_c < \lVert f_0 \rVert$ then $f = C(f_0)(1 + \delta)$ with $\lVert
  \delta \rVert_c < 1$, and $1 + \delta$ is invertible by the geometric series, which converges
  by Theorem~\ref{thm:restricted-complete}.  Conversely if $f$ is a unit then
  $\lVert f \rVert_c \lVert f^{-1}\rVert_c \geq 1$ forces the norm to be attained at the
  constant term, and no other coefficient can attain it.
\end{proof}

This description of the units is the analytic input to Weierstrass preparation: the ``unit
factor'' produced there is recognised as a unit exactly by
Theorem~\ref{thm:restricted-units}.

\begin{proposition}
  \label{prop:reduction}
  \uses{prop:powerbounded}
  \lean{MvPowerSeries.Restricted.residueRingHom, MvPowerSeries.Restricted.residueEquiv}
  \leanok
  For an ideal $I$ of $R^\circ$ containing the topologically nilpotent elements, reduction gives
  a surjection $T^\circ \to (R^\circ/I)[x_i]$ onto the \emph{polynomial} ring, whose kernel is
  the series with all coefficients in $I$.  In particular the reduction of a restricted power
  series is a polynomial.
\end{proposition}

That the reduction is a polynomial rather than a power series is the whole content: it is the
formal counterpart of ``a restricted series has only finitely many coefficients of size close to
the norm'', and it is what makes the residue ring noetherian.

\section{Splitting off a variable}

\begin{theorem}
  \label{thm:finsuccequiv}
  \uses{prop:restricted-ring, thm:restricted-normedring}
  \lean{MvPowerSeries.Restricted.finSuccEquiv, MvPowerSeries.Restricted.finSuccIsometry}
  \leanok
  For a polyradius $c$ on $\mathrm{Fin}(n+1)$ there is a ring isomorphism
  \[
    R_c\langle x_0, \dots, x_n \rangle \;\cong\;
    \bigl(R_{(c_1,\dots,c_n)}\langle x_1, \dots, x_n\rangle\bigr)_{c_0}\!\langle x_0 \rangle,
  \]
  and it is an isometry for the Gauss norms.
\end{theorem}

\begin{proof}
  \leanok
  \uses{def:gaussnorm}
  The underlying bijection regroups the coefficient family, and restrictedness on either side is
  the same convergence statement grouped differently.  The isometry is immediate from
  Definition~\ref{def:gaussnorm}, since a supremum over pairs is the supremum of the inner
  suprema.
\end{proof}

Theorem~\ref{thm:finsuccequiv} is the engine of the multivariate theory.  It says the
$(n+1)$-variable Tate algebra \emph{is} the one-variable Tate algebra over the $n$-variable one,
with matching norms --- so any one-variable theorem valid over an arbitrary complete ultrametric
normed \emph{ring} applies, and transports back.  That is exactly why Chapter~\ref{chap:wp} is
stated at Martin's generality rather than over a field: the induction leaves the category of
fields at the first step.

\begin{lemma}
  \label{lem:isemptyequiv}
  \uses{prop:restricted-ring}
  \lean{MvPowerSeries.Restricted.isEmptyEquiv, MvPowerSeries.Restricted.isEmptyIsometry}
  \leanok
  Over an empty index type the restricted power series are the coefficients themselves,
  isometrically.  This is the base case of the induction.
\end{lemma}


%%FILE blueprint/src/chapter/RieszColeman.tex
\chapter{Entire Series, Coleman's Resultant, and the Riesz--Coleman Decomposition}
\label{chap:riesz-coleman}

Chapter~\ref{chap:compact} attached to a compactoid operator $u$ its Fredholm determinant
$\det(1-Tu)$ and proved Serre's theorems: $1-au$ is invertible exactly when $\det(1-Tu)$ is a
unit at $a$ (Theorem~\ref{thm:serre11}), a zero of the determinant is a reciprocal eigenvalue
(Theorem~\ref{thm:serre12}), and each zero carries a finite-dimensional Riesz summand
(Theorem~\ref{thm:riesz}).  Serre's decomposition is attached to a \emph{single} zero, and it
lives over a field.

The applications need more.  A slope decomposition splits off \emph{all} the eigenvalues of a
given valuation at once, which means splitting the determinant as $\det(1-Tu) = Q\cdot S$ with
$Q$ a polynomial and asking for the corresponding splitting of the space; and the coefficient
ring in the halo application of Chapter~\ref{chap:lwx-halo} is the Banach--Tate ring
$\Lhalo[1/T]$, not a field, and is not Noetherian.  This chapter formalises the general
statement over a Banach--Tate ring, following Johansson--Newton \cite{JN} Theorem~2.2.2, which
is Buzzard \cite{Buzzard} Theorem~3.3, which is Bella\"iche \cite{Bellaiche} Theorem~II.2.18,
and ultimately Coleman \cite{Coleman} Theorem~A4.3.  No Noetherian hypothesis appears anywhere:
none of the sources' proofs uses one.

The route is Bella\"iche's.  Given the factorisation $\det(1-Tu) = QS$, one replaces $u$ by the
auxiliary operator $\varphi' = 1 - Q^{*}(u)/Q^{*}(0)$ and computes \emph{its} Fredholm
determinant by a spectral mapping formula; that determinant turns out to have a zero of order
$\deg Q$ at $1$, and Serre's projector at that zero is the projector one wants.  So we need
three things first: the ring $R\{\{T\}\}$ of entire series and division inside it (\S1),
Coleman's resultant $D(B,P)$ and the spectral mapping formula (\S2), and then the decomposition
itself (\S3).  Section~4 adds the vertex factorisation that produces the polynomial $Q$ in the
first place.

\section{Entire series and division}

\begin{definition}
  \label{def:isentire}
  \uses{def:isrestricted}
  \lean{PowerSeries.IsEntire, PowerSeries.entireSubring}
  \leanok
  A power series $F = \sum a_n T^n$ over a normed ring $R$ is \emph{entire} if
  $\lVert a_n\rVert C^n \to 0$ for every $C > 0$; equivalently, if it is restricted at every
  radius in the sense of Definition~\ref{def:isrestricted}.  The entire series form a subring
  $R\{\{T\}\}$ of $R\llbracket T\rrbracket$.  A \emph{Fredholm series} is an entire series with
  constant term $1$.
\end{definition}

This is \cite[Definition II.1.16]{Bellaiche} and the ring written $R\{\{T\}\}$ in
\cite[\S2.1]{JN}; Theorem~\ref{thm:charpowerseries-entire} says the Fredholm determinant of a
compactoid operator lies in it.

\begin{theorem}[Euclidean division in $R\{\{T\}\}$, {\cite[Proposition II.2.8]{Bellaiche}}]
  \label{thm:entire-division}
  \uses{def:isentire, thm:wdivision, thm:wdivision-unique}
  \lean{PowerSeries.IsEntire.exists_eq_mul_add, PowerSeries.eq_mul_add_q_unique,
        PowerSeries.eq_mul_add_r_unique}
  \leanok
  Let $g \in R[T]$ have invertible leading coefficient and degree $n$.  Every entire $F$ can be
  written uniquely as $F = g\,q + r$ with $q$ entire and $r$ a polynomial of degree $< n$.
\end{theorem}

\begin{proof}
  \leanok
  \uses{thm:wdivision, thm:wdivision-unique}
  Existence and uniqueness are Martin's Weierstrass division (Theorem~\ref{thm:wdivision} and
  Theorem~\ref{thm:wdivision-unique}) applied at a radius large enough that $g$ dominates: the
  divisor of a distinguished division is precisely a polynomial whose leading coefficient
  achieves the Gauss norm, and for a polynomial with invertible leading coefficient this happens
  at all sufficiently large radii.  Entireness of the quotient is then the statement that the
  division is performed at every radius simultaneously, which uniqueness supplies.
\end{proof}

\begin{definition}
  \label{def:isentirecoprime}
  \uses{def:isentire}
  \lean{PowerSeries.IsEntireCoprime}
  \leanok
  Two entire series are \emph{relatively prime in $R\{\{T\}\}$} if the ideal they generate is
  the unit ideal (\cite[Definition 2.2.1]{JN}).
\end{definition}

\begin{lemma}[{\cite[Corollary II.2.9]{Bellaiche}}]
  \label{lem:coprime-remainder}
  \uses{def:isentirecoprime, thm:entire-division}
  \lean{PowerSeries.isEntireCoprime_iff_isCoprime_of_eq_mul_add}
  \leanok
  If $F = g\,q + r$ is the division of Theorem~\ref{thm:entire-division}, then $F$ and $g$ are
  relatively prime in $R\{\{T\}\}$ if and only if $r$ and $g$ are relatively prime in $R[T]$.
\end{lemma}

Division reduces coprimality questions to polynomial ones; that is the only use we make of it.

\begin{definition}
  \label{def:isgoodzero}
  \uses{def:isentire}
  \lean{PowerSeries.IsGoodZero, PowerSeries.hasseDeriv_mul}
  \leanok
  An entire $F$ has a \emph{good zero of order $s$} at $a$ if the divided derivatives
  $\Delta^{i}F$ vanish at $a$ for $i < s$ and $\Delta^{s}F(a)$ is a unit
  (\cite[Definition II.2.11]{Bellaiche}).  The divided derivatives obey the Leibniz rule
  $\Delta^{s}(FG) = \sum_{i+j=s}\Delta^{i}F\cdot\Delta^{j}G$, so good zeros multiply.
\end{definition}

\section{Coleman's $D(B,P)$ and the spectral mapping formula}

Coleman's construction \cite[\S A3--A4]{Coleman} takes a polynomial $B$ and a Fredholm series
$P$ and produces the series whose reciprocal roots are the $B$-images of those of $P$:
informally $D(B,P)(T) = \prod_i (1 - T\,B(t_i))$ over the roots $t_i$ of the reciprocal
polynomial $P^{*}$.  Written that way the definition needs roots; written as a resultant it
needs nothing.

\begin{definition}
  \label{def:dpoly}
  \lean{TateFredholm.Coleman.dPoly, TateFredholm.Coleman.bQ}
  \leanok
  For polynomials $B, P$ over a commutative ring and $N \geq \deg P$, set
  \[
    D_N(B,P) \;=\; \Res_X\bigl(P^{*}, \; 1 - T\,B(X)\bigr) \in R[T],
  \]
  the resultant of the reversed polynomial $P^{*} = \mathrm{reflect}_N P$ against
  $1 - T B(X)$.  We also write $b_Q = 1 - u^{-1}Q^{*}$ for a polynomial $Q$ with leading
  coefficient the unit $u$; this is Bella\"iche's auxiliary polynomial.
\end{definition}

\begin{lemma}[{\cite[Lemma II.2.13]{Bellaiche}}]
  \label{lem:dpoly-mul}
  \uses{def:dpoly}
  \lean{TateFredholm.Coleman.dPoly_mul}
  \leanok
  $D_N(B, PQ) = D_N(B,P)\,D_N(B,Q)$: the construction is multiplicative in its second argument.
\end{lemma}

\begin{proof}
  \leanok
  Multiplicativity of the resultant in its first argument, together with
  $\mathrm{reflect}(PQ) = \mathrm{reflect}(P)\,\mathrm{reflect}(Q)$ at compatible degrees.
\end{proof}

\begin{proposition}[The finite-dimensional case]
  \label{prop:resultant-charpoly}
  \uses{def:dpoly}
  \lean{Matrix.resultant_charpoly, Matrix.det_aeval_eq_prod_roots}
  \leanok
  For a square matrix $A$ over any commutative ring and any polynomial $g$,
  \[
    \Res(\mathrm{charpoly}\,A,\; g) \;=\; \det\bigl(g(A)\bigr).
  \]
\end{proposition}

\begin{proof}
  \leanok
  Over an algebraically closed field both sides equal $\prod_\lambda g(\lambda)$ over the
  eigenvalues with multiplicity.  The general case follows from the universal characteristic
  polynomial: both sides are polynomial identities in the matrix entries and the coefficients of
  $g$, so it suffices to verify them over $\Z[x_{ij}, y_k]$, which embeds in an algebraic closure
  of its fraction field.
\end{proof}

\begin{definition}
  \label{def:dseries}
  \uses{def:dpoly, def:isentire}
  \lean{TateFredholm.Coleman.dSeries, TateFredholm.Coleman.norm_coeff_dSeries_sub_le,
        PowerSeries.IsEntire.dSeries}
  \leanok
  For a Fredholm series $F$ over a Banach--Tate ring, $D(B,F)$ is the coefficientwise limit of
  $D_N(B, \mathrm{trunc}_N F)$.  The limit exists because each coefficient of $D_N$ is a fixed
  polynomial in the coefficients of the truncation, and that polynomial is Lipschitz in them;
  $D(B, \cdot)$ is multiplicative and carries entire series to entire series.
\end{definition}

Entireness uses Buzzard's rescaling trick \cite[p.~21]{Buzzard}: the identity
$\Res(Q,P) = \Res(u^{-n}Q(uT), P(uT))$ lets one renormalise the radius at which the estimate is
performed.

\begin{lemma}[{\cite[Lemma II.2.14 and Proposition II.2.15]{Bellaiche}}]
  \label{lem:dseries-goodzero}
  \uses{def:dseries, def:isentirecoprime, def:isgoodzero}
  \lean{TateFredholm.Coleman.isUnit_evalT_one_dSeries_bQ_iff,
        TateFredholm.Coleman.isGoodZero_dSeries_bQ}
  \leanok
  Let $Q$ be a polynomial with unit leading coefficient and constant term $1$, and $S$ a
  Fredholm series.  Then $D(b_Q, S)(1)$ is a unit if and only if $Q$ and $S$ are relatively
  prime in $R\{\{T\}\}$; and $D(b_Q, QS)$ has a good zero of order $\deg Q$ at $1$.
\end{lemma}

\begin{proof}
  \leanok
  \uses{lem:dpoly-mul, lem:coprime-remainder, def:isgoodzero}
  By Lemma~\ref{lem:dpoly-mul}, $D(b_Q, QS) = D(b_Q, Q)\,D(b_Q, S)$.  The first factor is
  computed directly: $b_Q$ was built so that $D(b_Q, Q)$ is $(1-T)^{\deg Q}$ up to a unit, which
  is a good zero of order $\deg Q$ at $1$.  The second factor is a unit at $1$ by the
  coprimality criterion, which is Lemma~\ref{lem:coprime-remainder} transported through the
  resultant.  Good zeros multiply by the Leibniz rule of Definition~\ref{def:isgoodzero}.
\end{proof}

\begin{theorem}[Spectral mapping, {\cite[Proposition II.2.16]{Bellaiche}}]
  \label{thm:spectral-mapping}
  \uses{def:dseries, prop:resultant-charpoly, def:iscompactoid}
  \lean{TateFredholm.Coleman.charPowerSeries_aeval}
  \leanok
  For a compactoid $\varphi$ on $c(I,R)$ and any polynomial $B$ with $B(0)=0$,
  \[
    \det\bigl(1 - T\,B(\varphi)\bigr) \;=\; D\bigl(B, \det(1-T\varphi)\bigr).
  \]
\end{theorem}

\begin{proof}
  \leanok
  \uses{prop:resultant-charpoly, def:dseries}
  For a finite matrix this is Proposition~\ref{prop:resultant-charpoly} rewritten in reversed
  form.  A compactoid operator is a limit of its truncations to finite blocks, both sides are
  continuous in the operator coefficientwise, and $D(B,\cdot)$ was defined as exactly that
  limit.
\end{proof}

\section{The Riesz--Coleman decomposition}

\begin{theorem}[{\cite[Theorem 2.2.2]{JN}; \cite[Theorem 3.3]{Buzzard};
                 \cite[Theorem A4.3]{Coleman}}]
  \label{thm:riesz-coleman}
  \uses{thm:spectral-mapping, lem:dseries-goodzero, thm:riesz, def:iscompactoid}
  \lean{TateFredholm.exists_rieszColemanProjection, TateFredholm.IsRieszColemanProjection}
  \leanok
  Let $R$ be a Banach--Tate ring, $u$ compactoid on $M = c(I,R)$, and suppose
  $\det(1-Tu) = Q\,S$ with $Q$ a polynomial with unit leading coefficient and constant term $1$,
  relatively prime to the Fredholm series $S$ in $R\{\{T\}\}$.  Then there is a continuous
  projector $p$, lying in the closure of $R[u]$, with $Q^{*}(u)\,w = p$ for some $w$ in that
  closure and $Q^{*}(u)^{\deg Q}(1-p) = 0$.
\end{theorem}

\begin{proof}
  \leanok
  \uses{thm:spectral-mapping, lem:dseries-goodzero, thm:riesz}
  Put $\varphi' = 1 - Q^{*}(u)/Q^{*}(0)$, a compactoid operator.  Theorem~\ref{thm:spectral-mapping}
  computes $\det(1 - T\varphi') = D(b_Q, \det(1-Tu))$, and Lemma~\ref{lem:dseries-goodzero}
  says that this has a good zero of order $\deg Q$ at $1$.  Serre's projector at that zero
  (the construction behind Theorem~\ref{thm:riesz}) is $p$.  Because the projector is produced
  as an operator-norm limit of polynomials in $\varphi'$, and $\varphi'$ is itself a polynomial
  in $u$, both $p$ and the inverse witness $w$ lie in the closure of $R[u]$ --- so they commute
  with $u$, not merely with $\varphi'$.
\end{proof}

The refinements of \cite[Proposition 3.2 and Theorem 3.3]{Buzzard} and the last sentence of
\cite[Theorem 2.2.2]{JN} follow.  We record them as the API of the predicate
$\mathrm{IsRieszColemanProjection}$.

\begin{theorem}[The full decomposition]
  \label{thm:riesz-coleman-full}
  \uses{thm:riesz-coleman}
  \lean{TateFredholm.IsRieszColemanProjection.finite,
        TateFredholm.IsRieszColemanProjection.projective,
        TateFredholm.IsRieszColemanProjection.rankAtStalk,
        TateFredholm.IsRieszColemanProjection.charPowerSeries_mul_one_sub,
        TateFredholm.IsRieszColemanProjection.charPowerSeries_mul,
        TateFredholm.IsRieszColemanProjection.ker_aeval_reverse,
        TateFredholm.IsRieszColemanProjection.isUnit_mul_one_sub_add,
        TateFredholm.IsRieszColemanProjection.eq_range_of_isTopCompl}
  \leanok
  With the notation of Theorem~\ref{thm:riesz-coleman}, write $N = \ker Q^{*}(u)$.  Then
  $M = N \oplus \mathrm{im}\,p$ topologically; $N$ is finitely generated projective of rank
  $\deg Q$ at every prime; $u$ is invertible on $N$; $\det(1-Tu \mid N) = Q$ and
  $\det(1-Tu\mid \mathrm{im}\,p) = S$; and $\mathrm{im}\,p$ is the unique $u$-stable closed
  complement on which $Q^{*}(u)$ is invertible.
\end{theorem}

\begin{lemma}[{\cite[Lemma 3.1]{Buzzard}}]
  \label{lem:coprime-iff-unit}
  \uses{def:isentirecoprime}
  \lean{TateFredholm.isEntireCoprime_iff_isUnit_aeval_reverse}
  \leanok
  $Q$ and $S$ are relatively prime in $R\{\{T\}\}$ if and only if $Q^{*}(u)$ is invertible on
  the summand where the determinant is $S$.
\end{lemma}

\begin{proposition}[Determinants of summands]
  \label{prop:det-summand}
  \uses{thm:riesz-coleman-full}
  \lean{TateFredholm.charPowerSeries_eq_mul_of_comm,
        TateFredholm.exists_polynomial_charPowerSeries_of_range_le,
        TateFredholm.exists_matrix_realisation}
  \leanok
  If $u$ commutes with a continuous projector then its determinant is the product of the
  determinants of the two pieces; on a finitely generated summand the determinant is a
  polynomial, computed as the reversed characteristic polynomial of any matrix realisation.
\end{proposition}

The matrix-level facts behind Proposition~\ref{prop:det-summand} are elementary but were not in
mathlib.

\begin{lemma}
  \label{lem:charpolyrev}
  \lean{Matrix.charpolyRev_map, Matrix.charpolyRev_mul_comm, Matrix.charpolyRev_eq_of_mul_eq,
        Matrix.exists_mul_eq_of_idempotent, Matrix.charpolyRev_eq_one_sub_pow_rank}
  \leanok
  Write $\mathrm{charpolyRev}\,M = \det(1 - XM)$.  It commutes with ring homomorphisms and
  satisfies Sylvester's identity $\det(1-XAB) = \det(1-XBA)$; an idempotent matrix with free
  range factors as $E = PQ$, $QP = 1$, and along such a factorisation
  $\mathrm{charpolyRev}\,M = \mathrm{charpolyRev}(QMP)$ whenever $ME = EM = M$.  Over a field,
  if moreover $M = E - N$ with $N$ nilpotent then
  $\mathrm{charpolyRev}\,M = (1-X)^{\mathrm{rank}\,E}$.
\end{lemma}

\section{The vertex factorisation}

Theorem~\ref{thm:riesz-coleman} takes the factorisation $\det(1-Tu) = QS$ as input.  Producing
it is a Newton-polygon statement: a vertex of the polygon at which the coefficient is a unit
splits off a polynomial factor of the corresponding degree.

\begin{definition}
  \label{def:dominant-index}
  \uses{def:isentire}
  \lean{PowerSeries.IsDominantIndex, PowerSeries.IsDominantPoly}
  \leanok
  The index $N$ is \emph{dominant} for $F = \sum a_k T^k$ at radius $\rho$ if
  $\lVert a_k\rVert\rho^k \leq \lVert a_N\rVert\rho^N$ for all $k$, with strict inequality for
  $k > N$; that is, $N$ is Bella\"iche's largest integer minimising $v(a_n) - n\nu$
  (\cite[Definition II.2.3]{Bellaiche}).  A polynomial of degree $N$ dominant at $\rho$ with
  invertible leading coefficient is a \emph{$\rho$-dominant polynomial}
  (\cite[Definition II.2.4]{Bellaiche}).
\end{definition}

\begin{lemma}
  \label{lem:dominant-distinguished}
  \uses{def:dominant-index, def:ismuldistinguished}
  \lean{PowerSeries.IsDominantPoly.isMulDistinguished, PowerSeries.IsDominantPoly.mul,
        PowerSeries.exists_isDominantPoly_of_isUnit_leadingCoeff}
  \leanok
  A $\rho$-dominant polynomial is exactly a Martin-distinguished element at radius $\rho$
  (Definition~\ref{def:ismuldistinguished}); dominant polynomials are closed under
  multiplication.
\end{lemma}

That identification is what lets Chapter~\ref{chap:wp} be reused here.

\begin{theorem}[The slope factorisation, {\cite[Theorem II.3.6]{Bellaiche}};
                {\cite[Definition 2.2.5]{JN}}]
  \label{thm:slope-factorisation}
  \uses{def:dominant-index, lem:dominant-distinguished, thm:wdivision, def:isentirecoprime}
  \lean{PowerSeries.exists_isDominantFactorization, PowerSeries.IsDominantFactorization,
        PowerSeries.IsDominantFactorization.isEntireCoprime}
  \leanok
  Let $F$ be entire over a Banach--Tate ring, with dominant index $N$ at radius $\rho$ and with
  $\lVert\cdot\rVert$ multiplicative on the dominant coefficient $a_N$.  Then $F = PG$ with $P$
  a $\rho$-dominant polynomial of degree $N$ and $G$ entire of dominant index $0$ at $\rho$;
  moreover $P$ and $G$ are relatively prime in $R\{\{T\}\}$.
\end{theorem}

\begin{proof}
  \leanok
  \uses{lem:dominant-distinguished, thm:wdivision, thm:wprep}
  Bella\"iche and Kedlaya \cite[Proposition 3.2.2]{Kedlaya} both run a Newton iteration
  $P_{l+1} = P_l + X_l$, $Q_{l+1} = Q_l + Y_l$ driven by the division estimates
  $\lVert f\rVert = \max(\lVert g\rVert\lVert q\rVert, \lVert r\rVert)$.  We instead observe that
  those division estimates \emph{are} Martin's, so the hypotheses say exactly that $F$ is
  distinguished of order $N$ at radius $\rho$ (Lemma~\ref{lem:dominant-distinguished}); dividing
  $T^N$ by $F$ (Theorem~\ref{thm:wdivision}) and subtracting the remainder from $T^N$ produces
  the monic factor.  This is Weierstrass preparation (Theorem~\ref{thm:wprep}) in the form
  needed.  For coprimality (the norm-level form of \cite[Lemma 2.2.7]{JN}, and where \cite{JN}
  argue with resultants over the Gelfand spectrum): $\lVert G-1\rVert_\rho < 1$, so $G$ is a unit
  of the Tate algebra of radius $\rho$; dividing $G^{-1}$ by $P$ there gives
  $1 = Pq' + rG$ with $\deg r < N$, and the cofactor is entire because the entire division of
  $1 - rG$ by $P$ has the same quotient, by uniqueness.  That is a B\'ezout identity for $P$ and
  $G$ in $R\{\{T\}\}$.
\end{proof}

Together with Theorem~\ref{thm:riesz-coleman-full}, Theorem~\ref{thm:slope-factorisation} says:
\emph{a vertex of the Newton polygon of $\det(1-Tu)$ whose coefficient is a multiplicative unit
splits off a $u$-stable summand of rank equal to the abscissa of the vertex, carrying exactly
the eigenvalues of smaller slope.}  This is the form in which the halo $U_p$ of
\S\ref{sec:tate-riesz} is decomposed.

\section{Auxiliary constructions}

Three further pieces of the model-space toolkit are used in
Chapters~\ref{chap:lwx-halo}--\ref{chap:lwx-slopes} and recorded here.

\begin{lemma}[Two-sided weight bounds]
  \label{lem:two-sided}
  \uses{lem:hadamard, def:minor}
  \lean{TateFredholm.norm_charCoeff_le_pow_two_sided}
  \leanok
  Suppose the matrix of $u$ satisfies $\lVert u_{ab}\rVert \leq \sigma^{\max(w(a) - w'(b),\,0)}$
  for weights $w, w'$.  Then $\lVert c_n(u)\rVert \leq \sigma^{f(n)}$ whenever
  $f(n) \leq \sum_{a\in S} w(a) - \sum_{a \in S} w'(a)$ for all $\lvert S\rvert = n$.
\end{lemma}

\begin{proof}
  \leanok
  \uses{lem:hadamard}
  A permutation $\pi$ of $S$ satisfies $\sum_{a\in S} w'(\pi a) = \sum_{a \in S} w'(a)$, so every
  Leibniz monomial of $\mathrm{minor}(u,S)$ carries the same column contribution and the row
  contribution is $\sum_{a\in S}w(a)$.  This is the Hadamard bound of Lemma~\ref{lem:hadamard}
  with the column weight subtracted; note that nothing is inverted, so the estimate runs over an
  arbitrary complete ultrametric ring rather than a field.
\end{proof}

This is the estimate behind \cite[Theorem 3.16]{LWX}, where $\sigma = \lVert T\rVert$ and the
conjugation by $\mathrm{diag}(T^{w})$ that \cite{LWX} performs explicitly is replaced by the
observation that principal minors are invariant under diagonal conjugation.

\begin{lemma}[Diagonal intertwining]
  \label{lem:diag-intertwine}
  \uses{def:minor}
  \lean{TateFredholm.minor_eq_of_diag_intertwine, TateFredholm.charCoeff_eq_of_diag_intertwine,
        TateFredholm.charPowerSeries_eq_of_diag_intertwine}
  \leanok
  If $D\,M_v = M_u\,D$ entrywise for a diagonal $D$ with unit entries then
  $\mathrm{minor}(v,S) = \mathrm{minor}(u,S)$ for every finite $S$, and hence
  $\det(1-Tv) = \det(1-Tu)$.  No compactness hypothesis and no inversion of $D$ is needed.
\end{lemma}

The point of avoiding $D^{-1}$ is that the diagonal occurring in \cite[Proposition 2.17]{LWX} is
$\mathrm{diag}(n!)$, whose inverse is unbounded.

\begin{definition}[Block-diagonal and rectangular block maps]
  \label{def:blockmaps}
  \uses{prop:blockop}
  \lean{TateFredholm.blockDiag, TateFredholm.diagBlockEquiv, TateFredholm.blockOpMap,
        TateFredholm.blockMap, TateFredholm.blockMapEquiv}
  \leanok
  An operator $f$ on $c(I,R)$ induces the block-diagonal operator $\mathrm{blockDiag}\,f$ on
  $c(\sigma\times I, R)$, an equivalence if $f$ is; more generally a map
  $f : c(I,R) \to c(I',R)$ changing the fibre induces $\mathrm{blockMap}\,f$, and a family
  $T_{ij} : c(I,R)\to c(I',R)$ induces the rectangular block operator
  $\mathrm{blockOpMap}\,T$, with the expected composition rules against
  Proposition~\ref{prop:blockop}.
\end{definition}

\begin{theorem}[Perturbed unitriangular matrices are isometries]
  \label{thm:unitriangular}
  \uses{def:cspace}
  \lean{TateFredholm.IsUnitriangularPerturbation, TateFredholm.ofPerturbation,
        TateFredholm.norm_ofPerturbation, TateFredholm.surjective_ofPerturbation,
        TateFredholm.equivOfPerturbation}
  \leanok
  Let $M$ be an infinite matrix over a complete ultrametric field with entries of norm
  $\leq 1$, unit diagonal, entries strictly below the diagonal of norm $\leq q < 1$, and
  finitely supported columns.  Then $M$ defines an isometry of $c(\N,K)$ with dense range, hence
  a continuous linear equivalence.
\end{theorem}

\begin{proof}
  \leanok
  Norm-preservation is a largest-index argument: in $Mx$ the coordinate at the largest index
  where $x$ attains its norm receives the unit diagonal contribution and nothing else of that
  size, the sub-diagonal entries being smaller by the factor $q$.  Surjectivity is $q$-adic
  successive approximation.
\end{proof}

This is the criterion behind Colmez's proof of Amice's theorem, which Colmez states through the
residue field of a discretely valued base; the argument above needs only completeness and the
ultrametric inequality, which is what \S\ref{sec:amice} requires.

\begin{lemma}[Pairing of characteristic roots]
  \label{lem:charpoly-pairing}
  \lean{Matrix.det_mul_det_of_mul_eq_smul, Matrix.C_det_mul_charpolyRev_eq,
        Matrix.roots_charpolyRev, Matrix.roots_charpoly_of_mul_eq_smul}
  \leanok
  If $AB = c\cdot 1$ then $A$ is invertible with $B = c\,A^{-1}$, $\det A\det B = c^{\,n}$, and
  over an algebraically closed field
  $(\mathrm{charpoly}\,B).\mathrm{roots} = (\mathrm{charpoly}\,A).\mathrm{roots}$ mapped by
  $x \mapsto c/x$.
\end{lemma}

The root statement is phrased as an identity of multisets rather than as a sorted indexing
$\alpha_i = k+1-\alpha_{n-1-i}$: summing recovers the slope total, and counting recovers the
per-slope multiplicities, which are the two consumers in \S\ref{sec:stepone}.


%%FILE blueprint/src/chapter/Slopes.tex
\chapter{Slopes of Compact Hecke Operators}
\label{chap:slopes}

We can now combine the two halves of the thesis.  Chapter~\ref{chap:compact} attaches to a
compact operator its Fredholm determinant and identifies the reciprocal roots with the
eigenvalues; Chapter~\ref{chap:np} reads the valuations of the roots of a power series off the
slopes of its Newton polygon.  Composing:

\begin{center}
  \emph{the slopes of the Newton polygon of $\det(1 - Tu)$ are the valuations of the eigenvalues
  of $u$.}
\end{center}

So to say something about the eigenvalues of a Hecke operator it suffices to bound the
coefficients of its determinant.  This chapter does that in general, and the next specialises to
the case where the bound is attained.

\section{From row decay to coefficient bounds}

Let $u$ be a compactoid operator whose matrix decays along a row weight $w$, meaning $\lVert
u_{ji} \rVert \leq \sigma^{w_j}$.  A minor of $u$ picks entries from $\lvert S \rvert$ distinct
rows, so by the Hadamard bound its norm is at most $\sigma$ to the power of the weight sum over
$S$; the coefficient $c_n$ is an ultrametric sum of $n$-element minors, hence bounded by the
worst such.  The content is therefore combinatorial: how small can $\sum_{j \in S} w_j$ be for
$\lvert S \rvert = n$?

\begin{lemma}
  \label{lem:minor-bound}
  \uses{lem:hadamard, def:iscompactoid}
  \lean{TateFredholm.norm_minor_le_pow_sum}
  \leanok
  If $\lVert u_{ji}\rVert \leq \sigma^{w_j}$ for all $i, j$ then $\lVert \mathrm{minor}(u, S)
  \rVert \leq \sigma^{\sum_{j \in S} w_j}$ for every finite $S$.
\end{lemma}

\begin{theorem}[The slope bound]
  \label{thm:slope-bound}
  \uses{lem:minor-bound, def:minor}
  \lean{TateFredholm.norm_charCoeff_le_pow}
  \leanok
  Let $0 \leq \sigma < 1$, let $u$ be compactoid with $\lVert u_{ji}\rVert \leq \sigma^{w_j}$,
  and let $f : \N \to \N$ satisfy $f(n) \leq \sum_{j \in S} w_j$ for every $S$ with $\lvert S
  \rvert = n$.  Then
  \[
    \lVert c_n(u) \rVert \leq \sigma^{f(n)} .
  \]
\end{theorem}

\begin{proof}
  \leanok
  \uses{lem:minor-bound}
  $c_n$ is $\pm$ the sum of the $n$-element minors, and an ultrametric infinite sum is bounded by
  the supremum of its terms; each term is bounded by Lemma~\ref{lem:minor-bound} and then by
  $\sigma^{f(n)}$ since $\sigma \leq 1$ and $f(n)$ is the smaller exponent.
\end{proof}

Two combinatorial inputs cover the cases we need.

\begin{lemma}
  \label{lem:choose-two}
  \lean{TateFredholm.choose_two_le_sum, TateFredholm.choose_two_lt_sum_of_ne_range}
  \leanok
  Among $n$-element subsets of $\N$, the sum $\sum_{i \in S} i$ is minimised uniquely by $S =
  \{0, 1, \dots, n-1\}$, with value $\binom{n}{2}$; every other $n$-element subset has sum at
  least $\binom{n}{2} + 1$.
\end{lemma}

\begin{lemma}
  \label{lem:block-weight}
  \lean{TateFredholm.sum_div_le_sum_block}
  \leanok
  On $\iota \times \N$ with $\lvert \iota \rvert = d$ and weight $(i, m) \mapsto m$, an
  $n$-element subset has weight sum at least $\sum_{k < n} \lfloor k/d \rfloor$: each value of
  $m$ can be used at most $d$ times.
\end{lemma}

\begin{proof}
  \leanok
  Peel off an element $M$ of maximal second coordinate.  Since $S \subseteq \iota \times \{0,
  \dots, M_2\}$, which has $d(M_2+1)$ elements, we get $n \leq d(M_2+1)$ and hence $\lfloor
  (n-1)/d\rfloor \leq M_2$.  Induct on the remaining $n-1$ elements.
\end{proof}

\begin{corollary}
  \label{cor:slope-bound-block}
  \uses{thm:slope-bound, lem:choose-two, lem:block-weight}
  \lean{TateFredholm.norm_charCoeff_le_pow_choose_two, TateFredholm.norm_charCoeff_le_pow_block}
  \leanok
  On $c(\N, K)$ with $\lVert u_{ji}\rVert \leq \sigma^j$ we get $\lVert c_n \rVert \leq
  \sigma^{\binom{n}{2}}$; on the block model $c(\iota \times \N, K)$ we get $\lVert c_n \rVert
  \leq \sigma^{\sum_{k<n} \lfloor k/\lvert\iota\rvert\rfloor}$.
\end{corollary}

\section{The slope theorem at a general weight}

Applying this to a Hecke operator is now immediate, because
Theorem~\ref{thm:qmf-compact} produced exactly the required row decay.

\begin{proposition}
  \label{prop:qmf-row-decay}
  \uses{thm:qmf-compact, def:heckeblock}
  \lean{QMF.Weight.norm_matrixCoeff_heckeBlock_le, QMF.Weight.norm_matrixCoeff_heckeBlockOp_le}
  \leanok
  With $\lVert \det\theta(\eta)\rVert \leq \sigma$, $\rho \leq \sigma$, and a twisting character
  of norm at most $1$, the block operator satisfies $\lVert \mathrm{matrixCoeff}_{(i,m),(j,r)}
  \rVert \leq \sigma^{m}$.
\end{proposition}

\begin{theorem}
  \label{thm:qmf-slope-bound}
  \uses{prop:qmf-row-decay, cor:slope-bound-block, def:heckecharpowerseries}
  \lean{QMF.Weight.norm_charCoeff_heckeCharPowerSeries_le}
  \leanok
  The coefficients of $\det(1 - T[U\eta U])$ satisfy
  \[
    \lVert c_n \rVert \leq \sigma^{\sum_{k < n}\lfloor k/\lvert \iota\rvert\rfloor}.
  \]
\end{theorem}

Translating a coefficient bound into a polygon statement is what
Definition~\ref{def:ofslopes} was built for.

\begin{definition}
  \label{def:blockslopes}
  \lean{QMF.Weight.blockSlopes, QMF.Weight.monotone_blockSlopes}
  \leanok
  The \emph{block slope sequence} is $s_k = \lfloor k/d \rfloor \cdot (-\log \sigma)$, where $d =
  \lvert \iota \rvert$ is the size of the class set: each of the values $0, -\log\sigma,
  -2\log\sigma, \dots$ repeated $d$ times, once per block.  It is monotone since $\sigma < 1$.
\end{definition}

\begin{theorem}[The slope theorem at a general analytic weight]
  \label{thm:qmf-slope-theorem}
  \uses{thm:qmf-slope-bound, def:blockslopes, def:ofslopes, def:isnewtonpolygonof}
  \lean{QMF.Weight.isBelow_newtonPolygon_heckeCharPowerSeries}
  \leanok
  The Newton polygon of $\det(1 - T[U\eta U])$ lies on or above the polygon anchored at the
  origin whose $k$-th unit slope is $\lfloor k/\lvert\iota\rvert\rfloor \cdot (-\log \sigma)$.
\end{theorem}

\begin{proof}
  \leanok
  \uses{thm:ofslopes, thm:np-exists, thm:qmf-slope-bound}
  The determinant has constant coefficient $1$, so its polygon is anchored at the origin, and all
  its coefficients have norm at most $1$, which is an affine floor and hence admissibility.  The
  comparison polygon is $\mathrm{ofSlopes}$ of the block slope sequence, whose height at $k$ is
  $\sum_{i<k} s_i$; the required inequality height $\leq -\log\lVert c_k\rVert$ is
  Theorem~\ref{thm:qmf-slope-bound} after taking logarithms.  Now apply the maximality clause of
  Definition~\ref{def:isnewtonpolygonof}.
\end{proof}

In words: on a class set of size $d$, the slopes of $U_\varpi$ at weight $\kappa$ are at least
$0, \dots, 0, -\log\sigma, \dots$, each value repeated $d$ times.  This is the general form of
[Jacobs, Theorem 2.12]; the theorem itself asserts equality, and equality needs a hypothesis.

\section{When the bound is attained}

\begin{theorem}[{[Jacobs, Theorem 2.12]}]
  \label{thm:unit-minors}
  \uses{thm:slope-bound, lem:choose-two}
  \lean{TateFredholm.norm_charCoeff_of_unit_minors, TateFredholm.val_charCoeff_of_unit_minors}
  \leanok
  Let $u$ on $c(\N,K)$ satisfy $\lVert u_{ji}\rVert \leq \lVert \varpi \rVert^j$ and suppose the
  rescaled matrix $N_{ji} = \varpi^{-j} u_{ji}$ has every top-left minor of norm $1$.  Then
  \[
    \lVert c_m(u) \rVert = \lVert \varpi \rVert^{\binom{m}{2}} \qquad \text{exactly}.
  \]
\end{theorem}

\begin{proof}
  \leanok
  \uses{lem:choose-two, lem:minor-bound}
  The principal block $S = \{0,\dots,m-1\}$ contributes exactly $\lVert\varpi\rVert^{\binom
  m2}$: factoring $\varpi^j$ out of the $j$-th row leaves the unit-minor matrix.  Every other
  $m$-element block contributes strictly less, since its row sum exceeds $\binom m2$ by
  Lemma~\ref{lem:choose-two}.  An ultrametric sum with one strictly dominant term has the norm of
  that term.
\end{proof}

\begin{corollary}[The slope reading]
  \label{cor:slope-reading}
  \uses{thm:unit-minors, thm:ofslopes, cor:ofslopes-transport}
  \lean{JacobsSlash.unitSlope_newtonPolygon₀OfPowerSeries_charPowerSeries}
  \leanok
  Under the hypotheses of Theorem~\ref{thm:unit-minors} the Newton polygon of $\det(1-Tu)$, in
  the normalisation $v(\varpi) = 1$, has unit slopes $0, 1, 2, 3, \dots$.  Consequently the
  eigenvalues of $u$ have valuations $0, 1, 2, \dots$: an arithmetic progression.
\end{corollary}

\begin{proof}
  \leanok
  \uses{thm:ofslopes}
  The valuations $v(c_m) = \binom m2$ are the partial sums of $0,1,2,\dots$, so
  Theorem~\ref{thm:ofslopes} identifies the polygon and
  Corollary~\ref{cor:ofslopes-transport} reads its slopes.  Passing to eigenvalues is
  Theorem~\ref{thm:np-slope-reading} together with Theorem~\ref{thm:serre12}.
\end{proof}

There is a second case, with slopes $\tfrac12, \tfrac32, \tfrac52, \dots$, which occurs when the
field is ramified over $\Qp$ so that half-integers are values of the valuation.  That is what
happens for Jacobs's $U_3$, and it is the subject of the next chapter.


%%FILE blueprint/src/chapter/ToRewrite/Affinoid.tex
\chapter{Rigid Analytic Geometry}

\section{Affiniod algebras}

In this section we aim to begin an exposition on rigid analytic geometry following on from our
work on restricted power series.

\begin{definition}
  Let $k$ be a field complete with respect to a non-trivial non-archimedean valuation
  $\lvert \cdot \rvert_k$. A $k$-Banach algebra is a commutive $k$-algebra equipped with a function
  $\lvert \cdot \rvert : A \to \mathbb{R}$ satisfying:
  \begin{itemize}
    \item $\lvert 1 \rvert \leq 1$, and $\lvert a \rvert = 0 \iff a = 0 $;
    \item $\lvert a + b \rvert \leq \max \left\{ \lvert a \rvert, \lvert b \rvert \right\}$;
    \item $\lvert a b \rvert \leq \lvert a \rvert \lvert b \rvert$;
    \item $\lvert \lambda a \rvert = \lvert \lamba \rvert_k \lvert a \rvert$;
  \end{itemize}
  and such that $A$ is compelte with respect to the metric induced by $\lvert \cdot \rvert$.
\end{definition}
%% Could rewrite this definition in a nicer way

Recall the Tate algebra, $T_n$ is defined as the algebra of all power series in $n$ variables over
$k$ whose coefficients form a zero sequence.
%% This has to be adapted to what ever form I choose in my other section.

\begin{definition}
  A $k$-Banach algebra $A$ is called affinod ($k$-affiniod) if there exists an integer $n \geq 0$
  and a continuous epimorphism $\alpha : T_n \to A$.
\end{definition}

By Banach's theorem %(find reference in BGR)
$\alpha$ is open; hence $A$ is isomorphic tas $k$-Banach algebra to the residue algebra
$T_n / \ker \alpha$

\section{Affiniod sets}

$T_n$ can be viewed in two ways: $k_a$-valued functions on $B^n(k_a)$; and the algebra of functions
on $\max T_n$.
%% Note this assumes c = 1 ... BGR has Tate algebra only defined for c = 1, and shows for more
%% general polydiscs they are an affiniod algebra.

The second case comes from the following. For $f \in T_n$ and $\mathfrak{m} \in \max T_n$, we can
define $f (\mathfrak{m})$ as the residue class of $f$ in $T_n / \mathfrak{m}$. Since
$T_n / \mathfrak{m}$ is a finite algeraic extension of $k$, it can be embedded as a subfield of
$k_a$, and thus $f (\mathfrak{m})$ can be viewed as an element of $k_a$. However, this embedding is
not unique, in general, thus $f (\mathfrak{m})$ is only determined up to conjugation over $k$.
Therefore, setting $\Gamma = \text{Gal}(k_a/k)$, any $f \in T_n$ can be viewed as a function on
$\max T_n$ taking values in $k_a / \Gamma$.

\begin{definition}
  Let $F$ be a subset of $T_n$. We denote by
  \begin{align*}
    V_a(F) & := \left\{ x \in B^n(k_a) | f(x) = 0 \forall f \in F \right\} \\
    V(F) & := \left\{ \mathfrak{m} \in \max T_n | f (\mathfrak{m}) = 0 \forall f \in F \right\}
  \end{align*}
  the zero sets for $F$, and call these \emph{affiniod subsets} of $B^n (k_a)$ and $\max T_n$, respectively.
\end{definition}

These sets are interchangeable, $V(F)$ is the quotient of $V_a(F) by \Gamma$, and $V_a(F)$ is the
image of $V(F)$ with respect to the projection $B^n(k_a) \to \max T_n$.

%% Decide what else I want to include from BGR

\begin{definition}
  A $k$-affiniod variety is a pair $\text{Sp} A = (\max A, A)$ where $A$ is a $k$-affiniod algebra.
\end{definition}


%%FILE blueprint/src/chapter/ToRewrite/CompactOperators.tex
\chapter{Compact Operators}



%% Will rewrite all of this in terms of Kevins Eigenvariety or Bellaiche's Eigenbook (probably the latter)

In this section we introduce the theory of comapct operators over Banach spaces and work towards
defining the Fredholmd determinant.

Let $p$ be a prime and $K$ a complete subfield of $\mathbb{C}_p$ with absolute value $|\cdot|_p$.
Write $\mathcal{O}$ for the valuation ring of $K$ and $m$ its maximal ideal and let $\nu_p$ denote
the valuation on $K$.

\begin{definition}
    A Banach space over $K$ is a complete, normed vector space over $K$ whose norm $|\cdot|$
    satisfies the non-archimedean property.
\end{definition}

\begin{remark}
    By our assumptions this means $K$ itself is a Banach space with norm $|\cdot| = | \cdot |_p.$
\end{remark}

\begin{definition}
    Given an index set $I$ and a set of elements $x_i \in K$, we say $x_i$ tends to $0$ and write
    $x_i \to 0$ if for all $\varepsilon > 0$, there exists finitely many $i \in I$ such that
    $|x_i| > \varepsilon.$

    We write $c(I)$ the set of sequences $x = (x_i)_{i \in I}$ with $x_i \in K$ such that
    $x_i \to 0$.
\end{definition}

\begin{remark}
    If $I$ is finite it is trivial that all sets of elements tends to $0$; thus
    $c(I) \simeq K^{|I|}$.
\end{remark}

\begin{lemma}\label{Norm}
    If we define a function $|\cdot|_{c(I)} : c(I) \to \R$ by
    \[
    |x|_{c(I)} = \sup_{i \in I} |x_i|_p
    \]
    then this is a norm and gives $c(I)$ the structure of a Banach space.
\end{lemma}

\begin{proof}
    It is not hard to see that $|x|_{c(I)}$ is a non-archimedean norm; every property follows from the fact $|x|_p$ is a non-archimedean norm and taking supremums.
    To be a Banach space it remains to show $c(I)$ is complete. So consider a Cauchy sequence $x_j = ((x_{j,i})_{i \in I})$, by which me mean
    \[
    \forall \varepsilon > 0, \; \exists M > 0 \text{ such that } \forall k,j > M \text{ we have }
    |x_j - x_k |_{c(I)} \leq \varepsilon.
    \]
    That is
    \[
    \sup_{i \in I} |x_{j,i} - x_{k,i} |_p \leq \varepsilon
    \]
    and so we have $\forall i$, $|x_{j,i} - x_{k,i} |_p \leq \varepsilon$. Thus $(x_{j,i})_{j \in \N}$ are Cauchy sequences in $K$ and since $K$ is complete they converge.
    So consider
    \[
    x = (x_i)_{i \in I}
    \]
    where $x_i = \lim_{j \to \infty} x_{j,i}$. We want to show this is the limit of the $x_j$ and is also in $c(I)$.
    Recall for all $j,k > M$ we have $|x_{j,i} - x_{k,i}|_p \leq \varepsilon$, letting $j \to \infty$ we see $|x_i - x_{k,i} | \leq \varepsilon$ for all $i$. Therefore,
    \[
    |x - x_j |_{c(I)} = \sup_i |x_i - x_{j.i}|_p < \varepsilon
    \]
    and indeed $x$ is the limit. Finally, they belong in $c(I)$ as $x_{j,i} \to 0$ for a fixed $j$, and taking limits on $j$ preserves inequalities, and so we also have $x_i \to 0$. We thus conclude the proof.
\end{proof}

We now want to give a basis to our Banach space.

\begin{definition}
    Let $E$ be a Banach space over $K$ and suppose there exists $\{ e_i : i \in I \}$ such that every $x \in E$ can be written uniquely as $\sum_{i \in I} x_i e_i$ with $x_i \to 0$ and $|x| = \sup_{i \in I} | x_i |$. Then we call $\{e_i\}_{i \in I}$ an orthonormal basis for $E$.
\end{definition}

Moreover, we can always find a basis.

\begin{lemma}
    For any Banach space $E$ over $K$ there exists an index set $I$ such that $E  \cong c(I)$ and we can identify an orthonormal basis $\{ e_i\}_{i \in I}$ of $E$.
\end{lemma}

This follows from \cite[Proposition~1]{Serre} and \cite[Lemma~1]{Serre}.

As Banach spaces are vector spaces, we can consider (continuous) linear maps from one space to another. From now on let us suppose $E$ and $F$ are Banach spaces over $K$.

\begin{definition}
     Write $\mathcal{L}(E,F)$ for the vector space of continuous linear maps from $E$ to $F$.
\end{definition}

\begin{remark}
    We call this a vector space; and indeed it is not hard to show it actually is. If $f,g \in \mathcal{L}(E,F)$ then by linearity of linear maps $f + g$ is a linear map, and continuous via addition of continuous of maps, hence $f + g \in \mathcal{L}(E,F)$. Scalar multiplication is equally straightforward, and it is not hard to see that we can complete the argument easily.
\end{remark}

\begin{lemma}
   Consider the following function, $| \cdot | : \mathcal{L}(E,F) \to \mathcal{L}(E,F)$
    \[
    |u| = \sup_{x\neq 0} \frac{|u(x)|}{|x|}.
    \]
    This function acts as a norm on $\mathcal{L}(E,F)$ and gives it the structure of a Banach space.
\end{lemma}

We omit the proof, as it follows similar to that of Lemma~\ref{Norm}.

Let us consider a basis of $E$ and $F$, respectively $\{e_i : i \in I \}$ and $\{ f_j : j \in J \}$ for indexing sets $I$ and $J$. That is $E \cong c(I)$ and $F \cong c(J)$. Consider $u \in \mathcal{L}(E,F)$, we can write this as a matrix $(n_{ji})$ with entries in $K$. This follows from writing $u(e) = u(\sum_{i \in I} a_i e_i) = \sum_{i \in I} a_i u(e_i)$ and noting $u(e_i) \in F$, that is $u(e_i) = \sum_{j \in J} b_j f_j$. Thus
\[
u(e) = \sum_{i \in I} \sum_{b \in J} a_i b_j f_j
\]
and we see $n_{ji} = a_i b_j$.
We now note the following proposition, \cite[Proposition~3]{Serre}.

\begin{proposition}
The map $u \mapsto (u(e_i))_{i \in I}$ is an isomorphism of $\mathcal{L}(E,F)$ to the vector space of sequences $(f_i)_{i \in I}$ of $F$, equipped with the norm $\sup_{i \in I} |f_i|$.
\end{proposition}

This implies for a fixed $i$ the $|b_j|$ is bounded, and since they form coefficients of an orthonormal basis they tend to 0. That is, we have $|n_{ji}|$ is bounded and $n_{ji} \to 0$ for fixed $i$. Moreover, we have
\[
| u | = \sup_{i,j} |n_{ji}|. %% may need more explanation
\]

\begin{remark}
    Note nowhere we have required $I,J$ to be finite indexing sets, therefore $n_{ji}$ could be an infinite matrix - this is where we will use the previous section.
\end{remark}

\begin{definition}
    Let $\mathcal{F}(E,F)$ denote the subspace of $\mathcal{L}(E,F)$ consisting of the continuous, linear maps that are of finite rank.

    We say a map $u$ is compact if $u$ belongs to the closure of $\mathcal{F}(E,F)$. Write $\mathcal{C}(E,F)$ for this closure.
\end{definition}

We then have the following lemma which relates the compact property of a linear map to our work above.

\begin{lemma}
    Define $r_j (u) = \sup_{i \in I} |n_{ji} |$. We have
    \[
    u \text{ is compact} \quad \Longleftrightarrow \quad r_j(u) \to 0.
    \]
\end{lemma}

\begin{proof}
    Immediate by the following equivalences.
    \begin{align*}
        u \text{ compact } \iff &u \in \mathcal{C}(E,F) \\
        \iff &\text{isomorphic to a matrix } (n_{ji}) \text{ such that for fixed $i$} \\ & (n_{ji}) \to 0 \text{ under the norm } \sup_i |n_{ji}| := r_j (u)
    \end{align*} \qedhere
\end{proof}

From this we have the following useful corollary.

\begin{corollary}
    Suppose the matrix $(n_{ji})$ of $u$ has all its entries in $\mathcal{O}_K$. If $D(\frac{1}{p}) \cdot(n_{ji})$ has all entries in $\mathcal{O}_K$ then $u$ is compact. Recall $D(\frac{1}{p})$ is the infinite matrix with $(\frac{1}{p})^n$ on the diagonal; and note we can extend $n_{ji}$ to an infinite matrix with zeroes outside of its defined range such that the multiplication makes sense.
\end{corollary}

\begin{proof}
    Write $D(1/p)(n_{ji})$ as $\sum n_{ji} (1/p)^j x^jy^i$. If all its entries are in $\mathcal{O}_K$ then $n_{ji} (1/p)^j \in \mathcal{O}_K$. Let us fix $j$. Then for all $i$, we have $p^j | n_{ji}$ by definition of $\mathcal{O}_K$. Thus $|n_{ji}| < p^{-j}$, and therefore
    \[
    r_j (u) = \sup_{i} |n_{ji}| < p^{-j}
    \]
    and clearly this goes to zero as $j \to \infty$. Therefore by the above lemma we conclude that $u$ is compact.
\end{proof}

We will say a matrix is compact if it is the matrix of a compact operator.

\subsection{Characteristic power series of an operator}

Let $E$ be a Banach space isomorphic to $c(I)$ for some indexing set $I$.
We now state the following, \cite[Proposition~7]{Serre}.

\begin{proposition}
    Let $u : E \to E$ be a compact operator.
    \begin{enumerate}
        \item If the matrix of $u$ with respect to an orthonormal basis $\{e_i : i \in I \}$ for some index set $I$ is $(n_{ji})$ then
        \[
        \det (1 - t u) = \sum_{m = 0}^{\infty} c_m t^m
        \]
        where
        \[
        c_m = (-1)^m \sum_{ S \subset I, \; |S| = m} \sum_{\sigma \in \text{Sym}(S)} \text{sgn}(\sigma) \prod_{i \in S} n_{\sigma(i)i}
        \]
        \item The series $\det(1-tu)$ is an entire function of $t$, that is it has infinite radius of convergence.
        \item If $u_n \to u$ with $u_n \in \mathcal{C}(E,E)$ for all $n$, then $\det (1 - tu_n)$ tends to $\det(1-tu)$.
    \end{enumerate}
\end{proposition}

\begin{definition}
    We call the power series $\det(1 - tu)$ the characteristic power series of $u$. We can also refer to it as the characteristic power series of the matrix $(n_{ji}).$
\end{definition}


%%FILE blueprint/src/chapter/ToRewrite/Fujisaki.tex
\chapter{Fujisaki's Lemma}

This section aims to prove a classical result about the finiteness of a certain double coset. This
lemma plays a key role in determining that quaternionic modular forms are finite dimensional.

This lemma was completed as part of Kevin Buzzard's FLT project; and we will use his blueprint
as reference \ref{FLT}.

\section{Prerequesites}

\begin{definition}
    A \emph{global field} is one of two types of fields:
    \begin{itemize}
        \item an algebraic number field; or
        \item an irreducible algebraic curve over a finite field, or equivalently a finite extension
        of $\mathbb{F}_q(T)$, the field of rational functions in one variable over a finite field
        with $q = p^n$ elements.
    \end{itemize}
\end{definition}

\begin{definition}
    A \emph{division algebra} is a ring in which every non-zero element has multiplicative inverse.
\end{definition}

\begin{remark}
    Non-split quaternion algebras are division algebras. To see this, recall split quaternion
    algebras are isomorphic to $M_2 (K)$ and consider $\begin{pmatrix} 1 & 0 \\ 0 & 0
    \end{pmatrix}$, which clearly has no inverse.
\end{remark}

For the lemma we let $B$ be a finite-dimensional division algebra over a global field $K$.

%% May need another chapter where I define quaternion algebras and create such prerequesites...

We will introduce the following notation:
\begin{align*}
    &\underline{B}  = B \otimes_K \mathds{A} \\
    &\underline{B}_f  = B \otimes_K \mathds{A}_f \\
    &\underline{B}^{(1)} = \left\{ (x_\nu)_\nu \in \underline{B} : \prod_\nu |
      \text{nrd}(x_\nu)|_\nu = 1\right\}
\end{align*}

The main result we need is (reference).%%\cite[Theorem~27.6.14(a)]{Voight}.

\begin{lemma}\label{Fujisaki}
    Let $B$ be a finite-dimensional division algebra over a global field $F$, then the following
    statements hold
    \begin{itemize}
        \item $B^\times \leq \underline{B}^{(1)}$ is cocompact
        \item If $U$ is an open subgroup of $(\underline{B}_f )^\times$
          then the double coset space $B^\times \backslash (\underline{B}_f)^\times /U$ is finite.
    \end{itemize}
\end{lemma}


\section{The proof}

Within the FLT project we proved this lemma via a series of lemmas - for now directly copied from
his blueprint.

%% TODO : make the notation consistent - e.g. B and D; differing notation for finite adeles!

%% Not formalised by me.
\begin{lemma}
  \label{NumberField.AdeleRing.DivisionAlgebra.Aux.existsE}
  %\lean{NumberField.AdeleRing.DivisionAlgebra.Aux.existsE}
  \leanok
  There's a compact subset $E$ of $D_{\A}$
  with the property that for all $x\in D_{\A}^{(1)}$,
  the obvious map $xE\to D\backslash D_{\A}$ is not injective.
\end{lemma}

\begin{proof}
  \leanok
  We know that if we pick a $\Q$-basis for $D$ of size $d$ then this identifies $D$ with $\Q^d$,
  $D_{\A}$ with $\A_{\Q}^d$, and $D\backslash D_{\A}$ with $(\Q\backslash\A_{\Q})^d$. Now $\Q$ is
  discrete in $\A_{\Q}$, and the quotient $\Q\backslash \A_{\Q}$ is compact.
  Hence, $D$ is discrete in $D_{\A}$ and the quotient $D\backslash D_{\A}$ is compact.

  Fix a Haar measure $\mu$ on $D_{\A}$ and push it forward to $D\backslash D_{\A}$; by compactness
  this quotient has finite and positive measure, say $m\in\R_{>0}$. Choose any compact
  $E\subseteq D_{\A}$ with measure $> m$ (for example, choose a $\Z$-lattice $L\cong\Z^d$ in
  $D\cong\Q^d$, define $E_f:=\prod_p L_p\in D\otimes_{\Q}\A_{\Q}^\infty$, and define
  $E_{\infty}\subseteq D\otimes_{\Q}\R\cong\R^n$ to be a huge closed ball, large enough to ensure
  the measure of $E:=E_f\times E_{\infty}$ is bigger than $m$). Then $\mu(xE)=\mu(E)>m$ so the map
  can't be injective.
\end{proof}

\begin{definition}
  \label{NumberField.AdeleRing.DivisionAlgebra.Aux.E}
  %\lean{NumberField.AdeleRing.DivisionAlgebra.Aux.E}
  \uses{NumberField.AdeleRing.DivisionAlgebra.Aux.existsE}
  \leanok
  We let $E$ denote any compact set satisfying the hypothesis of the previous lemma.
\end{definition}

\begin{definition}
  \label{NumberField.AdeleRing.DivisionAlgebra.Aux.X}
  %\lean{NumberField.AdeleRing.DivisionAlgebra.Aux.X}
  \uses{NumberField.AdeleRing.DivisionAlgebra.Aux.E}
  \leanok
Define $X:=E-E:=\{e-f:e,f\in E\}\subseteq D_{\A}$.
\end{definition}

\begin{definition}
  \label{NumberField.AdeleRing.DivisionAlgebra.Aux.Y}
  %\lean{NumberField.AdeleRing.DivisionAlgebra.Aux.Y}
  \uses{NumberField.AdeleRing.DivisionAlgebra.Aux.X}
  \leanok
Define $Y:=X.X:=\{xy:x,y\in X\}\subseteq D_{\A}$.
\end{definition}

\begin{lemma}
  \label{NumberField.AdeleRing.DivisionAlgebra.Aux.X_compact}
  %\lean{NumberField.AdeleRing.DivisionAlgebra.Aux.X_compact}
  \uses{NumberField.AdeleRing.DivisionAlgebra.Aux.X}
  \leanok
  $X$ is a compact subset of $D_{\A}$.
\end{lemma}
\begin{proof}
  \leanok
  It's the continuous image of the compact set~$E\times E$.
\end{proof}

\begin{lemma}
  \label{NumberField.AdeleRing.DivisionAlgebra.Aux.Y_compact}
  %\lean{NumberField.AdeleRing.DivisionAlgebra.Aux.Y_compact}
  \uses{NumberField.AdeleRing.DivisionAlgebra.Aux.Y}
  \leanok
  $Y$ is a compact subset of $D_{\A}$.
\end{lemma}
\begin{proof}
  \leanok
  It's the continuous image of the compact set~$X\times X$.
\end{proof}

\begin{lemma}
  \label{NumberField.AdeleRing.DivisionAlgebra.Aux.X_meets_kernel}
  %\lean{NumberField.AdeleRing.DivisionAlgebra.Aux.X_meets_kernel}
  \uses{NumberField.AdeleRing.DivisionAlgebra.Aux.X, NumberField.AdeleRing.DivisionAlgebra.Aux.E}
  \leanok
  If $\beta\in D_{\A}^{(1)}$ then
$\beta X\cap D^\times\not=\emptyset$.
\end{lemma}
\begin{proof}
  \leanok
The map $\beta E\to D\backslash D_{\A}$ isn't injective, so there are distinct
$\beta e_1,\beta e_2\in \beta E$ with $e_i\in E$ and $\beta e_1-\beta e_2=b\in D$. Now $b\not=0$
and $D$ is a division algebra, so $b\in D^\times$. And $e_1-e_2\in X$ so $b=\beta(e_1-e_2)\in
\beta X$, so we're done.
\end{proof}

\begin{lemma}
  \label{NumberField.AdeleRing.DivisionAlgebra.Aux.X_meets_kernel'}
  %\lean{NumberField.AdeleRing.DivisionAlgebra.Aux.X_meets_kernel'}
  \uses{NumberField.AdeleRing.DivisionAlgebra.Aux.X}
  \leanok
  Similarly, if $\beta\in D_{\A}^{(1)}$ then $X\beta^{-1}\cap D^\times\not=\emptyset$.
\end{lemma}
\begin{proof}
  \leanok
  \uses{NumberField.AdeleRing.DivisionAlgebra.Aux.E, addHaarScalarFactor.left_mul_eq_right_mul}
  Indeed, $\beta^{-1}\in D_{\A}^{(1)}$, and so left multiplication by $\beta^{-1}$
  doesn't change Haar measure on $D_{\A}$, so neither does right multiplication.
  So the same argument works: $E\beta^{-1}\to D\backslash D_{\A}$ is not
  injective so choose $e_1\beta^{-1}\not=e_2\beta^{-1}$ with difference $b\in D$
  and then $(e_1-e_2)\beta^{-1}\in D-{0}=D^\times$.
\end{proof}

\begin{definition}
  \label{NumberField.AdeleRing.DivisionAlgebra.Aux.T}
  %\lean{NumberField.AdeleRing.DivisionAlgebra.Aux.T}
  \uses{NumberField.AdeleRing.DivisionAlgebra.Aux.Y}
  \leanok
  Let $T:=Y\cap D^\times$.
\end{definition}

\begin{lemma}
  \label{NumberField.AdeleRing.DivisionAlgebra.Aux.T_finite}
  %\lean{NumberField.AdeleRing.DivisionAlgebra.Aux.T_finite}
  \uses{NumberField.AdeleRing.DivisionAlgebra.Aux.T}
  \leanok
  $T$ is finite.
\end{lemma}
\begin{proof}
  \uses{NumberField.AdeleRing.DivisionAlgebra.Aux.Y_compact}
  \leanok
  It suffices to prove that $Y\cap D$ is finite. But $D\subseteq D_{\A}$ is a discrete additive
  subgroup, and hence closed. And $Y\subseteq D_{\A}$ is compact. So $D\cap Y$ is compact and
  discrete, so finite.
\end{proof}

\begin{definition}
  \label{NumberField.AdeleRing.DivisionAlgebra.Aux.C}
  %\lean{NumberField.AdeleRing.DivisionAlgebra.Aux.C}
  \uses{NumberField.AdeleRing.DivisionAlgebra.Aux.T, NumberField.AdeleRing.DivisionAlgebra.Aux.X}
  \leanok
  Define $C:= (T^{-1}.X) \times X\subset D_{\A}\times D_{\A}$.
\end{definition}

\begin{lemma}
  \label{NumberField.AdeleRing.DivisionAlgebra.Aux.C_compact}
  %\lean{NumberField.AdeleRing.DivisionAlgebra.Aux.C_compact}
  \uses{NumberField.AdeleRing.DivisionAlgebra.Aux.C,
    NumberField.AdeleRing.DivisionAlgebra.Aux.T_finite,
    NumberField.AdeleRing.DivisionAlgebra.Aux.X_compact}
  \leanok
  $C$ is compact.
\end{lemma}
\begin{proof}
  \leanok
   $X$ is compact and $T$ is finite.
\end{proof}

\begin{lemma}
  \label{NumberField.AdeleRing.DivisionAlgebra.Aux.antidiag_mem_C}
  %\lean{NumberField.AdeleRing.DivisionAlgebra.Aux.antidiag_mem_C}
  \uses{NumberField.AdeleRing.DivisionAlgebra.Aux.C}
  \leanok
  For every $\beta\in D_{\A}^{(1)}$, there exists $b\in D^\times$
  and $\nu\in D_{\A}^{(1)}$ such that $\beta=b\nu$ and $(\nu,\nu^{-1})\in C.$
\end{lemma}
\begin{proof}
  \leanok
  \uses{NumberField.AdeleRing.DivisionAlgebra.Aux.E,
    NumberField.AdeleRing.DivisionAlgebra.Aux.X_meets_kernel,
    NumberField.AdeleRing.DivisionAlgebra.Aux.X_meets_kernel',
    NumberField.AdeleRing.DivisionAlgebra.Aux.C}
  We have $\beta X\cap D^\times\not=\emptyset$, and $X\beta^{-1}\cap D^\times\not=\emptyset$,
  so we can write $\beta x_1=b_1$ and $x_2\beta^{-1}=b_2$ with $b_i\in D^\times$ and $x_i\in X$.
  Note that $\beta\in D_{\A}^{(1)}$ and $b_i\in D^{\times}\subseteq D_{\A}^{(1)}$, so
  $x_i\in D_{\A}^{(1)}$ as well. In particular $x_i\in D_{\A}^\times$ so $x_1^{-1}$ makes sense.

  Multiplying the equations defining the $x_i$ and $b_i$ we deduce that
  $x_2x_1=b_2b_1\in Y\cap D^\times=T$ (recall that $Y=X.X$ and $T=Y\cap D^\times$ is finite); call
  this element $t$. Then $x_1^{-1}=t^{-1}x_2\in T^{-1}.X$, and $x_1\in X$, so if we set
  $\nu=x_1^{-1}\in D_{\A}^{(1)}$ and $b=b_1\in D^\times$ then we have $\beta=b\nu$ and
  $(\nu,\nu^{-1})\in C := (T^{-1}.X)\times X$. We are done!
\end{proof}

We can now prove the first part of Fujisaki's lemma.

\begin{proof}
  \proves{NumberField.AdeleRing.DivisionAlgebra.compact_quotient}
  \leanok
  \uses{MeasureTheory.ringHaarChar_continuous,
    NumberField.AdeleRing.DivisionAlgebra.Aux.antidiag_mem_C,
    NumberField.AdeleRing.DivisionAlgebra.Aux.C_compact}
  Indeed, if $M$ is the preimage of $C$ under the inclusion $D_{\A}^{(1)} \to D_{\A}\times D_{\A}$
  sending $\nu$ to $(\nu,\nu^{-1})$, then $M$ is a closed subspace of a compact space so it's
  compact (note that $\delta_{D_{\A}}$ is continuous, so $D_{\A}^{(1)}$ is a closed subset of
  $D_{\A}^\times$ which is itself a closed subset of $D_{\A}\times D_{\A}$). $M$ surjects onto
  $D^\times\backslash D_{\A}^{(1)}$ which is thus also compact.
\end{proof}

And now the second part.

\begin{theorem}
  \label{NumberField.FiniteAdeleRing.DivisionAlgebra.units_cocompact}
  %\lean{NumberField.FiniteAdeleRing.DivisionAlgebra.units_cocompact}
  \leanok
  $D^\times\backslash(D\otimes_K\A_K^\infty)^\times$ is compact.
\end{theorem}
\begin{proof}
  \uses{NumberField.AdeleRing.DivisionAlgebra.compact_quotient}
  \leanok
There's a natural map $\alpha$ from $D^\times\backslash D_{\A}^{(1)}$ to
  $D^\times\backslash (D\otimes_K \A_K^\infty)^\times$. We claim that it's
  surjective. Granted this claim, we are home, because if we put the quotient
  topology on $D^\times\backslash (D\otimes_K \A_K^\infty)^\times$ coming from
  $(D\otimes_K \A_K^\infty)^\times$ then it's readily verified that $\alpha$
  is continuous, and the continuous image of a compact space is compact.

  As for surjectivity: say $x\in (D\otimes_K \A_K^\infty)^\times$. We need to extend
  $x$ to an element $(x,y)\in (D\otimes_K \A_K^\infty)^\times\times(D\otimes_K K_\infty)^\times$
  which is in the kernel of $\delta_{D_{\A}}$. Because $\delta_{D_{\A}}(x,1)$ is some positive
  real number, it will suffice to show that if $r$ is any positive real number then we can
  find $y\in (D\otimes_K \A_K^\infty)^\times=(D\otimes_{\Q}\R)^\times$ with $\delta_{D_{\A}}(1,y)=r$,
  or equivalently (setting $D_{\R}=D\otimes_{\Q}\R$) that $\delta_{D_{\R}}(y)=r$.
  But $D\not=0$ as it is a division algebra,and hence $\Q\subseteq D$, meaning
  $\R\subseteq D_{\R}$, and if
  $x\in\R^\times\subseteq D_{\R}^\times$ then $\delta(x)=|x|^d$ with $d=\dim_{\Q}(D)$,
  as multiplication by $x$ is just scaling by a factor of $x$ on $D_{\R}\cong\R^d$.
  In particular we can set $x=y^{1/d}$.
\end{proof}

\begin{remark} In this generality the quotient might not be Hausdorff.
\end{remark}

\begin{theorem}
  \label{NumberField.FiniteAdeleRing.DivisionAlgebra.finiteDoubleCoset}
  %\lean{NumberField.FiniteAdeleRing.DivisionAlgebra.finiteDoubleCoset}
  \uses{NumberField.FiniteAdeleRing.DivisionAlgebra.units_cocompact}
  \leanok
  If $U$ is an open subgroup of $(D\otimes_K \A_K^\infty)^\times$
  then the double coset space $D^\times\backslash(D\otimes_K \A_K^\infty)^\times/U$ is finite.
\end{theorem}
\begin{proof}
  \leanok
  The double cosets give a disjoint open cover of $(D\otimes_K \A_K^\infty)$
  which descends to a disjoint open cover of the quotient space
  $D^\times\backslash(D\otimes_K \A_K^\infty)^\times$. However this space is compact
  by theorem~\ref{NumberField.FiniteAdeleRing.DivisionAlgebra.units_cocompact}.
\end{proof}


%%FILE blueprint/src/chapter/ToRewrite/GaussNorm.tex
\chapter{Gauss Norm}

Thus, far we have shown that for a non-archimedean normed ring $R$, $R_c \langle x_1, \dots, x_n \rangle$
is a ring for any parameter $c$. We now want to endow this with a non-archimedea absolute value.

Towards this let us consider the following function on multivariate power series:

\begin{definition}
Define $\lvert \cdot \rvert _c :  R_c \langle x_1, \dots, x_n \rangle \to \R$ where
\[
    \lvert f \rvert _c = \sup_{i \in \mathbb{N}^n} \lvert a_{(i_1, \cdots, i_n)} \rvert
    \prod_n \lvert c \rvert _n^{i_n},
\]
for $f(x) = \sum a_{(i_1, \cdots, i_n)} x_1^{i_1} \cdots x_n^{i_n}$ and call this the Gauss norm
with parameter $c$.
\end{definition}

We may think to ask is this a norm on the whole space of multivariate power series; but this
has a very easy answer. Simply, for a general power series
\[
  \sup_{i \in \mathbb{N}^n} \lvert a_{(i_1, \cdots, i_n)} \rvert
    \prod_n \lvert c \rvert _n^{i_n}
\]
need not exist - specifically, unless the set is bounded above this would be undefined in $\R$.

Moreover, even if we restrict to power series that have bounded
\[
  \lvert a_{(i_1, \cdots, i_n)} \rvert \prod_n \lvert c \rvert _n^{i_n},
\]
which we will refer to as "having gauss norm". A product of two such power series need not have gauss
norm.
For example, consider $c = 1$ and
\[
  f = \sum x^n \in \Q_p [ \! [x] \! ]
\]
then
\[
  f^2 = \sum (n + 1) x^n,
\]
and we have $\lvert n + 1 \rvert_p \tendsto \infty$.

Nevertheless, restricting to $R_c \langle x_1, \dots, x_n \rangle$ is sufficient to prove what we want.

\begin{theorem}
  For $n$ a natural number, $c$ a tuple of non-zero real numbers and $R$ a non-archimedean ring. Then
  the gauss norm is a non-archimedean absolute value on $R_c \langle x_1, \dots, x_n \rangle$.
\end{theorem}

\begin{remark}
  We could alternatively define the gauss norm without $\lvert \cdot \rvert$ around the $c$ and instead
  require that $c$ is strictly positive in the theorem. The latter is the choice we make in Lean; this is
  to match the definition of restricted power series - where we do not have norms around the $c$ either -
  and to prevent having $\lvert \cdot \rvert_a = \lvert \cdot \rvert_{-a}.$

  We could argue to also require $0 \leq c$ in the definition of restricted power series as well, but in
  interest of keeping definitions the most general we opt to allow negative entries in $c$.
\end{remark}

Note we really require no entries of $c$ to be zero, as it is immediate to see that in this case we
would have $\lvert \cdot \rvert_c \equiv 0$.

\section{Lean implementation}

To prove the theorem, for each property of an absolute value we actually generalise the conditions
on $R$ and $c$ as much as possible. For example, instead of $R$ being a normed ring we only impose the
following variables:

\begin{itemize}
  \item $R$ is a semiring; and
  \item $v$ is a function $R \to \R$.
\end{itemize}

Then when required we strengthen can strengthen $R$ to a ring and impose properties on $v$.

Moreover, although we established that the gauss norm is not a norm on the space of power series -
some properties can still be proved, and assuming the power series "has gauss norm" we can prove
everything except multiplicativity.

We will now state the harder statements we give towards showing the gauss norm is a
non-archimedean absolute value on $R_c \langle x_1, \dots, x_n \rangle$.

First, we generalise the definition of gauss norm into and formally define what we mean by a power
series "has gauss norm".

\begin{definition}
  We define the gauss norm of a power series $f \in R_c [ \! [ x_1, \dots, x_n ] \! ]$ as
  \[
    \lvert f \rvert_c = \sup_{(i_1, \cdots, i_n)} v (a_{(i_1, \cdots, i_n)})  \prod_n \lvert c \rvert _n^{i_n}
  \]
  where $f = \sum a_{(i_1, \cdots, i_n)} x_1^{i_1} \cdots x_n^{i_n}$.
\end{definition}

\begin{definition}
  Let $R$ be a semiring. We say a power series $f = \sum a_{(i_1, \cdots, i_n)} x_1^{i_1} \cdots x_n^{i_n}$
  has gauss norm if the set
  \[
    v (a_{(i_1, \cdots, i_n)})  \prod_n \lvert c \rvert _n^{i_n}
  \]
  is bounded above.
\end{definition}

Then, the non-archimedean property can be stated as.

\begin{lemma}
  If $R$ is a semiring, $c$ is non-negative and $v$ is a non-negative non-archimedean function then
  for two power series $f, g \in R_c [ \! [ x_1, \dots, x_n ] \! ]$ that have gauss norm, then
  \[
    \lvert f + g \rvert_c \leq \max \left\{ \lvert f \rvert_c, \lvert g \rvert_c \right\}.
  \]
\end{lemma}

%% I also want to state the gaussNorm_mul_le and gaussNorm_le_mul + antidiagonal_dominant statements
%% But need to decide how I want to do this


%%FILE blueprint/src/chapter/ToRewrite/Intro.tex
\chapter{Introduction}

Jacquet-Langlands suggests that certain classical modular forms should be related to modular forms
on quaternion algebras. This is because classical modular forms can be viewed as automorphic forms
for $M_2 (\Q)$; which is also a quaternion algebra. Compared to the analytic defintion of classical
modular forms, for definite quaternion algebras, the definition of quaternionic modular forms is
much nicer to work with - specifically, being of a concrete combinatorial nature. Thus, to study
(some) classical modular forms, it may be easier to study quaternionic modular forms.

In this project the aim is to define such quaternionic modular forms over definite quaternion
algebras and then prove some results surrounding them in Lean. Specifcally, we aim to formalise
Daniel Jacobs' thesis "slopes of compact Hecke operators", which says that $p$-adic valuations of the
eigenvalues of certain Hecke operators over quaternionic modular forms are in arithmetic
progression.

\subsection{What is being formalised?}

In attempting to formalise Daniel Jacobs' theis paper we need to first build up foundation material
before we can start working on his paper. This includes:

\begin{enumerate}
  \item restricted power series;
  \item the Gauss norm; and
  \item Newton polygons.
\end{enumerate}

Each of these can be found in its own chapter of the blueprint. Where we introduce the topic and
state what we need in Lean. This is by no means the entirity of the field, but should be sufficient
for what we need.

After what we are dubbing the foundational material, we need to start working on results on
quaternionic modular forms. This includes:
\begin{enumerate}
  \item their definition;
  \item a result on their finite dimensionality (Fujisaki's lemma); and
  \item their Hecke operators.
\end{enumerate}

Finally, with this we can hopefully start formalising the results from Jacobs' thesis.


%%FILE blueprint/src/chapter/ToRewrite/NP.tex
\chapter{Newton Polygon}

\section{Definitions}

In this section we will determine a way to check for zeroes of a power series via the construction
of Newton polygons. We consider power series over a non-archimedean valued field $R$ where we denote
the \emph{additive} valuation by $v$ and the associated multiplicative valuation by $\lvert \cdot \rvert$.
The typical example would be $\Q_p$.

These are formally given by the following definition.

\begin{definition}
    The \emph{Newton polygon} of $f = \sum a_i x^i$ is the lower boundary of the convex hull
    of the set of points $(i,\nu (a_i))$, where we ignore points with infinite $p$-adic valuation.
    Here by convex hull we mean the smallest convex polygon that contains all points.
\end{definition}

Given a Newton polygon we can give the following information.

\begin{definition}
    \begin{enumerate}
        \item The \emph{slopes} of the line segments appearing in the Newton polygon are called the
        slopes of $f(x)$.
        \item The \emph{length} of a slope is the length of the projection of the line segment with
        that slope onto the $x$-axis.
        \item The \emph{breaks} are the values of $i$ which are vertices of the Newton polygon.
    \end{enumerate}
\end{definition}

It turns out when working with Newton polygons the slopes and lengths are the information we care about;
and moreover given a starting point using this we can reconstruct all breaks. This provides motivation
for the definition we give in Lean.

\begin{definition}
  A Newton polygon is the following information:
  \begin{itemize}
    \item a natural number $n$, determining the number of slopes;
    \item a (potentially infinite) list of slopes that are strictly increasing except for the final
      slope (which may only be equal to the previous);
    \item a (potentially infinite) list of lengths; and
    \item a starting point.
  \end{itemize}
\end{definition}

It is clear that given a Newton polygon, in the formal sense, we can construct the following information;
however the reverse also holds by using the comment above the definition and the fact the slopes are
(strictly) increasing give the convexity condition.
%% Perhaps I could elaborate on this further?

We now want to construct a way to obtain this information given a power series, $f = \sum_{i=0}^\infty a_n x^i \in R[x]$.
Towards this we consider the following algorithm:

\begin{enumerate}
    \item Plot the points $\{(i,\nu (a_i) )\}$.
    \item Start with the vertical half-line down from the point with the smallest $i$.
    \item Rotate the line until one of the following happens:
    \subitem The line simultaneously hits infinitely many of the points we have plotted. In this
        case don't break the line and the polygon is complete.
    \subitem The line reaches a position where it can be rotated no further without leaving behind
        some points. In this case, the half lines forms the final segment of the polygon.
    \subitem The line hits a finite number of points. In this case, break at the last point that was
        hit and repeat the procedure.
\end{enumerate}

It should be clear that by construction this constructs the lower convex hull of the set of points
$\{(i,\nu (a_i) )\}$ and taking the slopes and lengths of the obtained segments we get our corresponding
definition in Lean.

As mentioned, the study of Newton polygons is deeply connected to studying the zeroes of $f(x)$,
thus we will assume $f(0) \neq 0$. That is, we factor out any powers of $x$ dividing $f$, which
means our initial vertical line starts on the $y$-axis.

Moreover, we will also assume $a_0 = 1$. This is because if we factor
$f(x) = \sum_{i = 0}^\infty a_i x^i$ as $f(x) = a_0 \cdot g(x)$, where
$g(x) = 1 + \sum_{i=1}^\infty \frac{a_i}{a_0}x^i$, it is not hard to see that the resulting Newton
polygon of $g(x)$ is the Newton polygon of $f(x)$ translated down by $\nu (a_0)$.
This fact follows from observing $\nu (\frac{a_i}{a_0}) = \nu(a_i) - \nu (a_0)$.

\section{Properties}

Let $K$ be a $p$-adic field, either a finite extension of $\mathbb{Q}_p$ or $\mathbb{C}_p$.
%% can this be generalised to any nonarch-valued field (probably?)

We want to consider a power series and consider its norms $\lvert f(x) \rvert_c$ for many different
$c$.

It is easy to see that
\[
c_1 > c_2 \implies \lvert f(x) \rvert_{c_1} \geq \lvert f(x) \rvert_{c_2}.
\]

Let us consider $f = 1 + a_1 x + \cdots a_n x^n$ a polynomial and consider its first segment. If this
segment has slope $m$, it connects the point $(0,0)$ to some other point $(i, mi)$ (i.e. slope $m$,
length $i$). Then we get the following properties:
\begin{itemize}
  \item $\lvert a_j \rvert \leq p^{-mj} = (p^{-m})^j$  for all $j$;
  \item $\lvert a_i \rvert = p^{-mi}$; and
  \item $\lvert a_j \rvert < p^{-mj}$ if $j > i$.
\end{itemize}

If we let $c = p^m$, we can read these conditions as
\begin{itemize}
  \item $\lvert f(x) \rvert_c = 1$; and
  \item $i$ is the largest integer such that $\lvert f(x) \rvert_c = \lvert a_i \rvert c^i.$
\end{itemize}

That is we see that $f$ is $i$-distinguished. In particular, this means, by Weierstrass preparation
that if $i$ is less than the degree of $f$ then, $f$ is divisible by a polynomial of degree $i$.

%% Note I actually need to formalise the statements about weirstrass div and prep that focuses on
%% if they are only polynomials now!

\begin{definition}
  A polynomial $f(x)$ is pure if its Newton polygon has only one slope. If this slope is $m$, we say
  $f$ is pure of slope $m$.
\end{definition}

\begin{proposition}
  Irreducible polynomials are pure.
\end{proposition}

\begin{proof}
  A break at $(i,mi)$ means a factor of degree $i$; hence, if there is a break at $i \neq 0, n$
  (so not pure) and hence the polynomial is reducible.
\end{proof}

\begin{proposition}
  A polynomial, of degree $n$, will be pure of slope $m$ if and only if it has the property that for
  $c = p^m$
  \[
  \lvert f(x) \rvert_c = \lvert b_n \rvert c^n = 1,
  \]
  i.e. it is $n$-distinguished.
\end{proposition}

\begin{theorem}
  Let $f(x) = 1 + a_1 x + \cdots + a_n x^n \in K[x],$ and assume the newton polygon of $f$ has its
  first break at $(i,mi)$. Then there exist polynomials $g(x), h(x) \in K[x]$ satisfying:
  \begin{itemize}
    \item $f(x) = g(x) h(x)$;
    \item $g(x)$ has degree $i$ and is pure of slope $m$; and
    \item $h(x)$ has no zeros in the closed ball of radius $p^m$ around 0.
  \end{itemize}
\end{theorem}

\begin{proof}
  Put everything above together.
\end{proof}

\begin{lemma}
  Let $f(x) = 1 + a_1 x + \cdots a_n x^n$, and assume the first break of the Newton polygon occurs
  at $(i,mi)$. Let $c$ be any positive real number less than $p^m$. Then we have
  \[
    \lvert f(x) \rvert_c = 1 \quad \text{and} \quad \lvert f(x) - 1 \rvert_c < 1.
  \]
\end{lemma}

\begin{proof}
  If $\lvert f(x) - 1 \rvert_c < 1$, then we must have $\lvert f (x) \rvert_c = 1$ by the usual all
  triangles are isosceles argument. Thus, we only need to prove the inequality.
  A line through the origin with slope $m_1 < m$ passes below all points on the polygon, touching it
  only at $(0,0)$. This means that $v_p(a_j) > m_1 j$ for every $j > 0$. Or in terms of absolute
  values $\lvert a_j \rvert < p^{-m_1 j}.$ Since $a_0 = 1$, the zeroth coefficient of $f(x) - 1$ is
  zero, and therefore also $\lvert a_0 - 1 \rvert < 1.$ It follows that $\lvert f(x) - 1 \rvert_c < 1$
  for $c < p^{m_1}$.
\end{proof}

\begin{lemma}
  Let $f(x) = 1 + a_1 x + \cdots + a_n x^n$, and assume that the first break of the newton polygon
  occurs at the point $(i,mi)$. Let $c$ be any positive real number less than $p^m$. Then $f$ has no
  zeros in the closed ball in $\mathbb{C}_p$ of radius $c$ around 0.
\end{lemma}

\begin{proof}
  The previous lemma says that $\lvert f(x) - 1 \rvert_c < 1$. From that, it follows for any $x$ such
  that $\lvert x \rvert \leq c$, we have $\lvert f(x) - 1 \rvert < 1$, which means $f(x) \neq 0$.
  Thus, $f(x)$ has no zeros in $\overline{B}(0,c)$.
\end{proof}

%% the reason we have C_p and just just Q_p bar is that we eventually want to talk about
%% convergence of power series... which requires completeness

%% Therefore, in Lean this can be changed. But the later will have to be in C_p
%% For generality this will be algebraic closure of an ultrametric normed ring (field?)
%% We then want its completion later on (see how C_p is defined by Maria Ines)

\begin{proposition}
  Let $f(x) = 1 + a_1 x + a_2 x^2 + \cdots a_n x^n$, and assume that the first break of the newton
  polygon occurs at the points $(i,mi)$. Then $f(x)$ has no roots with absolute value less than $p^m$
  and has exactly $i$ roots (counting multiplicity, in $\mathbb{C}_p$) with absolute value $p^m$.
\end{proposition}

\begin{proof}
  If the first break is at $(i,mi)$ then
  \begin{itemize}
    \item if $c < p^m$, $f(x)$ has no roots in the closed ball of radius $c$; and
    \item $f(x)$ factors as the product of a pure polynomial $g(x)$ of slope of $m$ and a polynomial
      which has no roots in the closed ball of radius $p^m$.
  \end{itemize}
  Any root of $g$ is a root of $f$, and therefore has no roots of absolute value less than $p^m$. On
  the other hand, if $\alpha_1, \cdots, \alpha_i$ are the roots of $g$, then we must have
  \[
    g(x) = (1 - \alpha_1^{-1}x) (1 - \alpha_2^{-1}x) \cdots (1 - \alpha_i^{-1}x),
  \]
  and hence the $i$-th coefficient is equal to $(\prod_i \alpha_i)^{-1}$. Since $g$ is pure, this
  must have valuation $mi$, and as $\lvert (a_j) \rvert \geq p^m$, it follows that all $a_j$ have
  absolute value $p^m$.
\end{proof}

\begin{theorem}
  Let $f(x) = 1 + a_1 x + \cdots + a_n x^n$ be a polynomial, let $m_1, \cdots, m_r$ be the slopes
  of its Newton polygon (in increasing order), and let $i_1, \cdots, i_r$ be their corresponding
  projected lengths. Then, for each $k$, $1 \leq k \leq r$, $f$ has exactly $i_k$ roots
  (in $\mathbb{C}_p$, counting multiplicity) of absolute value $p^{m_k}$.
\end{theorem}

\begin{proof}
  We have already seen this for the first segment. So let us move onto the second. Note this suffices
  to complete the proof as the same argument will yield the result for consecutive segments.

  Therefore, let us assume there are breaks at $(i_1,mi)$ and $(i_1 + i_2, m_1 i + m_2 i_2)$. We
  will set $k := i_1 + i_2$. Similar to the first slope we can conclude the following:
  \begin{itemize}
    \item $\lvert a_k \rvert = p^{-m_2 i_2 - m_1 i_1} = p^{-m_2 k} p^{(m_2 - m_1)i}$;
    \item for $i_1 \leq j \leq k$, $\lvert a_j \rvert \leq p^{-m_2 (j - i_1) - m_1 i_1}$; and
    \item for $j < i_1$ and $ j > k$, we have $\lvert a_j \rvert < p^{-m_2 (j - i_1)- m_1 i_1}.$
  \end{itemize}
  In other words, we get for $c = p^{m_2}$ that $\lvert f \rvert_c = p^{(m_2 - m_1)i_1}$, i.e. $f$
  is $k$-distinguished. Therefore, we can go through a completely analogous proof to before to get
  $f$ has exactly $k$ roots in the closed ball of radius $p^{m_2}$, $i_1$ of which have absolute
  value $p^{m_1}$ and $k-i_1 = i_2$ of which have absolute value $p^{m_2}$,
\end{proof}

With this we can now move onto considering power series and give the key theorem.

\begin{theorem} \label{RadofCon}
    Let $m$ be the supremum of the slopes appearing in the Newton polygon of a series
    $f(x) = 1 + a_1 x + a_2 x^2 + \cdots$. Then the radius of convergence of $f$ is $p^m.$
\end{theorem}

\begin{proposition}\label{NP}
    Let $f(x) = 1 + a_1 x + a_2 x^2 + \cdots$ be a power series. Let $m_1, m_2, \dots, m_k$ be the
    first $k$ slopes of the Newton polygon of $f(x)$, and assume that $f(x)$ converges on the
    closed ball of radius $c = p^{m_k}$ (from previous section, this means it is a restricted power
    series for parameter $c$ if $K$ is complete). Let $N$ be the $x$-coordinate of the right
    endpoint of the $k$-th segment of the Newton polygon. Then, there exists a polynomial
    $g(x)$ of degree $N$ and a power series $h(x)$ such that
    \begin{enumerate}
        \item $f(x) = g(x) h(x)$,
        \item $\| f(x) - g(x) \|_c < 1$,
        \item $h(x)$ converges on the closed ball of radius $c$,
        \item $\| h(x) - 1\|_c <1$, and
        \item the Newton polygon of $g(x)$ is the same as the portion of the Newton polygon of
        $f(x)$ contained in the region $0 \leq x \leq N$.
    \end{enumerate}
\end{proposition}

This then allows us to give our connection to zeroes of the power series.

\begin{corollary}
    Let $f(x) = 1 + a_1 x + a_2 x^2 + \cdots$ be a power series which converges on the closed ball
    of radius $c = p^m$. Let $m_1,\dots, m_k$ be the slopes of the Newton polygon of $f(x)$ which
    are less than or equal to $m$, and let $i_1,\dots,i_k$ be their lengths. Then, for each $j$, the
    power series $f(x)$ has $i_j$ zeros with absolute value $p^{m_j}$, and there are no other zeros
    in the closed ball of radius $p^m.$
\end{corollary}

\begin{proof}[Proof of proposition~\ref{NP}]
    We prove this by induction on $k$:\\
    If $k=1$, write the break point as $(N, m_1 N)$. We know all points are on or above the line
    $y = m x$, and the ones where $j > N$ are strictly above it. This translates to
    \begin{itemize}
        \item $| a_j | (p^j)^{m_1} \leq 1$ for all $j$,
        \item $| a_N | (p^N)^{m_1} = 1$, and
        \item $| a_j | (p^{m_1})^j$ < 1 if $j > N.$
    \end{itemize}
    This says that $\|f(x) \|_{p^{m_1}} = 1$ and that the maximum is last realised at the degree
    $N$. Thus, we can apply the Weierstrass preparation theorem to conclude that there is a
    polynomial of degree $N$ and a power series $h(x)$ satisfying
    \[
    \| h(x) - 1 \|_{p^m_1} < 1 \quad \|f(x) - g(x)\|_{p^m_1} < \|f(x) \|_{p^m_1} \quad \text{and} \quad f(x) = g(x) h(x).
    \]
    Combining the middle equation with $\| f(x) \|_{p^{m_1}} < 1$ means
    $\| f(x) - g(x) \|_{p^{m_1}} < 1$.
    Now, using facts about Newton polygons of polynomials, it follows that the zeros of $f(x)$ in
    the closed ball of radius $p^m$ coincide with the those of $g(x)$. That is, we must have the
    Newton polygon of $g(x)$ coincides with the Newton polygon of $f(x)$ on $0 \leq x \leq N$.
    Finally, we have $f(x)$ converges on the closed ball of radius $p^{m_1}$, by Theorem
    \ref{RadofCon}, and thus so does $h(x).$
    \\
    By induction, assume we have $f(x) = g_{k-1}(x)h_{k-1}(x)$ for a polynomial $g_{k-1}(x)$ and
    power series $h_{k-1}(x)$, such that the Newton polygon of $g_{k-1}(x)$ coincides with the
    first $k-1$ segments of the polygon for $f(x)$ and $h_{k-1}(x)$ has no zeroes in the closed ball
    of radius $p^{m_{k-1}}$.
    \\
    Let us now consider the $k$-th segment. By definition, the $k$-th segment ending at $x = N$
    means that for any $i > N$, $(i, \nu_p (a_i))$ lies above the line of slope $m_k$ through the
    point $(N, \nu_p (a_N))$. This means that, writing $c = p^{m_k}$, $\| f(x) \|_{p^{c}} =
    \lvert a_N \rvert c^N$ and $\lvert a_n \rvert c^N > \lvert a_i \rvert c^i $ for any $i > N$.
    Thus, we can apply the $p$-adic Weierstrass preparation theorem to get a polynomial $g(x)$.
    It is then easy to see $g_1 (x) \mid g(x)$, and that its Newton polygon coincides with the
    relevant portion of the Newton polygon of $f(x).$
\end{proof}

\begin{proof}[Proof of corollary]
    Assume we know this statement for polynomials, and the proposition says that the relevant
    part of the Newton polygon of $f(x)$ \emph{is} the Newton polygon of the polynomial $g(x)$.
    Since $g(x)$ is a factor of $f(x)$ and the quotient $h(x)$ clearly has no zeroes in the closed
    ball of radius $p^m$, the conclusion follows.
\end{proof}


%%FILE blueprint/src/chapter/ToRewrite/QMF.tex
\chapter{Quaternionic Modular Forms}

\section{Definition}

We define quaternionic modular forms over definite quaternion algebras explicitely as the following
space.

\begin{definition}
  Let $\kappa : \mathbb{Z}_p^\times \to \mathcal{O}_p^\times$ be a locally analytic character.
  Given $\alpha \in \mathbb{N}$, let
  \[
    \Sigma_\alpha = \left\{ \gamma =  \begin{pmatrix}
      a & b \\
      c & d
    \end{pmatrix}
    \in \text{M}_2(\mathbb{Z}_p) : p^\alpha \mid c, \; p \nmid d, \; \det (\gamma) \neq 0 \right\}.
  \]
  The weight $\kappa$ action of $\gamma =  \begin{pmatrix} a & b \\ c & d \end{pmatrix} \in
  \Sigma_\alpha$ on the Tate algebra $\mathcal{A}_p$ with coefficients in $\mathbb{C}_p$ is given by
  the continuous $\mathbb{C}_p$-linear extension of the map sending
  \[
    z^k \mapsto \frac{\kappa(c z + d)}{(cz + d)^2} \left( \frac{az+b}{cz+d} \right)^k
  \]
  and given $f(z) \in \mathcal{A}_p$ we write $(f||_\kappa \gamma)(z)$ for this action.

  Fix $\alpha$ and $\kappa$ in the above, let $D$ be a definite quaternion algebra and let $U$ be
  an open compact subgroup of $D_f^\times$, where $D_f = D  \otimes_{\mathbb{Q}} \mathbb{A}_f$ is
  $D$ over the finite adeles, of wild level $\geq p^\alpha$; that is the projection
  $U \to D_{p}^\times$ is contained in $\Sigma_\alpha.$
  Now let $A$ be any right $\Sigma_\alpha$-module.

  Then the level $U$, weight $\kappa$ space of automorphic forms is the space
  \[
    \mathcal{L}(U,A) = \{ \varphi : D_f^\times \to A  \;  | \; \varphi (dgu) = \varphi (g)
    ||_\kappa u_p, \; \forall d \in D^\times, \; g \in D^\times_f, \; u \in U \}.
  \]
\end{definition}

\section{Finite Dimensional}

\begin{lemma}
  Using the above notation and supposing $D_f^\times = \coprod_{i \in I} D^\times c_i U$
  is the decomposition into finite double cosets, then
  \[
    \mathcal{L}(U,A) \cong \bigoplus_{i \in I} A^{\Gamma_i}
  \]
  as a $\mathbb{C}_p$ vector space, where $\Gamma_i = c_i^{-1} D^\times c_i \cap U$.
\end{lemma}

\section{Hecke Operators}

We can define Hecke operators on these automorphic forms. We will keep the notation from the above
definitions.
Given $\varphi \in \mathcal{L}(U,A)$ we define a right action of $U$ on $\mathcal{L}(U,A)$:
\[
  (\varphi |_\kappa u) (g) := \varphi (gu^{-1}) ||_\kappa u_p.
\]

With this we have:
\begin{definition}
  Let $\nu \in D^\times_f$ such that $\nu_p \in \Sigma_\alpha$. Then the double coset $U\nu U$ may
  be written as a disjoint union
  \[
    U \nu U = \coprod_{t \in T} U \nu_t
  \]
  where $T$ is a finite set, and $\nu_t \in D_f^\times$. The Hecke operator is the map
  $[U \nu U] : \mathcal{L}(U,A) \to \mathcal{L}(U,A)$ given by
  \[
    [U \nu U] \varphi := \sum_{t \in T} \varphi |_\kappa \nu_t.
  \]
\end{definition}

\begin{definition}
  We set $\varpi_l$ to be the element of $\mathbb{A}_f$ which is $l$ at  $l$ and 1 at all other
  places, and set $\eta_l = \begin{pmatrix} \varpi_l & 0 \\ 0 &1 \end{pmatrix}$.\\
  The standard Hecke operators are then given by
  \[
    T_l := [U \eta_l U] \quad \text{and} \quad S_l:=[U \varpi U]
  \]
  and if $l = p$, we write $T_p = U_p$.
\end{definition}


%%FILE blueprint/src/chapter/ToRewrite/RPS.tex
\chapter{Restricted Power Series}

\section{Definition}

Recall formal multivariate power series over a set $K$ in $n$ variables, written
$K [ \! [ x_1, \dots, x_n ] \! ]$ are infinite sums
\[
  f(x) = \sum a_{(i_1, \cdots, i_n)} x_1^{i_1} \cdots x_n^{i_n},
\]
where the sum is taken over all $n$-tuples $(i_1, \cdots i_n)$ with $i_j \in \mathbb{N}$ and
$a_{(i_1, \cdots, i_n)} \in K$.

When $K$ is a ring, formal multivariate power series in $n$ variables form a ring, and we will refer
to them as the ring of power series over $K$ in $n$ variables. We again denote such
rings using $K [ \! [ x_1, \dots, x_n ] \! ]$.

We define a subset of $K [ \! [ x_1, \dots, x_n ] \! ]$ that consists of power series whose
coefficients approach zero in the cofinite filter.

\begin{definition}
  Let $f$ be a power series over a normed ring $R$ in $n$ variables, write
  $f(x) = \sum a_{(i_1, \cdots, i_n)} x_1^{i_1} \cdots x_n^{i_n}$, and let
  $c = (c_1, \cdots, c_n)$ be an $n$-tuple of real numbers.
  We say $f$ is a restricted power series of parameter $c$ over $R$ in $n$ variables if
  \[
    \lim_{n \to \infty} \lvert a_{(i_1, \cdots, i_n)} \rvert \prod_{j=1}^{n} c_j^{i_j} = 0
  \]
  in the cofinite topology.

  We denote the set of restricted power series of parameter $c$ over $R$ in $n$ variables by
  $R_c \langle \! [ x_1, \dots, x_n ] \! \rangle$.
\end{definition}

\begin{remark}
  Note
  \[
    \lim \lvert a_{(i_1, \cdots, i_n)} \rvert \prod_{j=1}^{n} c_j^{i_j} = 0
  \]
  in the cofinite topology, if for all $0 < \epsilon$ there are finitely many $n$-tuples that have
  $\left \lvert \lvert a_{(i_1, \cdots, i_n)} \rvert \prod_{j=1}^{n} c_j^{i_j} \right \rvert$
  greater than $\varepsilon$.
\end{remark}

Moreover, it actually does not matter about the sign of the values $c_i$, thus we may as well consider
$c_i > 0$ for all $i$.

\begin{lemma}
  A power series $f$ is restricted for parameter $c$ if and only if it is restricted for parameter
  $\lvert c \rvert$.
\end{lemma}

\section{Additive Group}

It turns out that restricted power series of parameter $c$ over a normed ring $R$ in $n$ variables
are not just a set, in fact they form an additive group.

\begin{lemma}
  For any normed ring $R$, natural number $n$, and tuple $c$,
  $R_c \langle \! [ x_1, \dots, x_n ] \! \rangle$ is an additive group.
\end{lemma}

\section{Ring}

Moreover, if we strengthen our condition of normed ring $R$, to a non-archimedean normed ring, we
in fact can show it is a ring.

\begin{theorem}
  For any non-archimedean normed ring $R$, $n \in \mathbb{N}$ and $n$-tuple $c \in \mathbb{R}^n$,
  $R_c \langle \! [ x_1, \dots, x_n ] \! \rangle$ is a ring.
\end{theorem}

The only extra thing we need to show is closure of multiplication. That is the product of two restricted
power series are also restricted.

To prove this in Lean, we instead give a more general statement that will give what we want as an
immediate consequence.

\begin{lemma}
  Let $M$ be an additive monoid with an antidiagonal, and $S$ a non-archimedean normed ring. Let
  $\phi : M \to \mathbb{R}$ be a function such that for all $a, b \in M, \phi(a + b) = \phi(a)
  \phi(b)$. If two functions $f, g : M \to S$ have the property that
  $\lim_i \lvert f i \rvert \phi (i) = 0$ in the cofinite filter and
  $\lim_i \lvert g i \rvert \phi (i) = 0$ then
  \[
  \lim_i \lvert \sum_{\mathclap{\substack{(a,b) \in \\ \textrm{antidiagonal } i}}} f(a) g(b) \rvert
  \phi(i) = 0.
  \]
\end{lemma}

\section{Tate algebra}

With this new definition we can now define a necessary component of our project.

\begin{definition}
  When $c \equiv 1$ and $R$ is a complete non-archimedean field, we call the ring of restricted power
  series over $R$ in $n$ variables, the \emph{Tate algebra} over $K$ in $n$ variables, we denote this
  by $T_n$.
\end{definition}

We can consider Tate algebras not just with variable $c \equiv 1$, that is more general polydiscs,
these turn out to be examples of affiniod algebras. See the Affinoid chapter for more information.


%% Note the following really needs to go below gauss norm as we need the ring of restricted power series
%% to be a normed ring (with the gauss norm!)

\section{Isomorphism}

When considering multivariate power series, it is convenient to consider the following "equality"
\[
  K [ \! [ x_1, \dots, x_n ] \! ] = (K [ \! [ x_1, \dots, x_{n-1} ] \! ])[\![ x_n ]\!]
\]
this is seen by writing $f \in K [ \! [ x_1, \dots, x_n ] \! ]$ in the form
\[
  f = \sum_{k = 0}^\infty g_k (x_1, \dots, x_n) x_n^k.
\]
Of course, by "equality" we really mean the spaces are isomorphic. We would like to have a similar
isomorphism for multivariate restricted power series.

\begin{lemma}
  $R_c \langle x_1, \dots, x_{n + 1} \rangle \equiv
  (R_c \langle x_1, \dots, x_n \rangle) \langle x_{n+1} \rangle.$
\end{lemma}

\begin{proof}
Let $\phi$ be the map given by the isomorphism of multivariate power series detailed above. It
suffices to show that the image by $phi, \phi^{-1}$ of restricted power series are restricted.

Assume for all $1 \leq m \leq n+1$, $c_m \neq 0$, else every power series is restricted trivially
and there is nothing to prove.

Towards this let $f = \sum a_{i_1, \dots i_{n+1}} x_1^{i_1} \cdots x_{n+1}^{i+1} \in
R_c \langle x_1, \dots, x_{n + 1} \rangle$. It being restricted means
\[
  \lim \lvert a_{(i_0, \cdots, i_{n+1})} \rvert \lvert c \rvert^{(i_0, \cdots, i_{n+1})} = 0
\]
under the cofinite filter. Write
\[
  \phi (f) = \sum_{s = 0}^\infty g_s x_n^s
\]
where $g_s \in R [ \! [ x_1, \dots, x_n ] \! ]$.
It suffices to show two things:
\begin{enumerate}
  \item for all $s$, $g_s \in R_c \langle x_1, \dots, x_n \rangle$; and
  \item $\phi (f)$ is restricted in $(R \langle x_1, \dots, x_{n-1} \rangle)[ \! [x_n] \! ]$.
\end{enumerate}

For 1. note for a fixed $s$, $g_s = \sum a_{((i_0, \cdots, i_n),s)} x_1^{i_1} \cdots x_n^{i_n}$;
thus we want to show
\[
  \lim \lvert a_{((i_0, \cdots, i_n), s)} \rvert \prod_{j=1}^{n} \lvert c_j \rvert ^{i_j} = 0
\]
under the cofinite filter. Suppose it does not, that is there exists an $\varepsilon$ such that there
are infinitely many $(i_0, \cdots, i_n)$ with
\[
  \lvert a_{((i_0, \cdots, i_n), s)} \rvert \prod_{j=1}^{n} \lvert c_j \rvert ^{i_j} > \varepsilon.
\]
But then for this fixed $s$ there are infinitely many $(i_0, \cdots, i_n)$ with
\[
  \left( \lvert a_{((i_0, \cdots, i_n), s)} \rvert \prod_{j=1}^{n} \lvert c_j \rvert ^{i_j} \right)
  \lvert c_n \rvert^s > \varepsilon \lvert c_n \rvert^s
\]
which would contradicted the restricted-ness of $f$ - because we have infinitely many $(i_1, \cdots, i_{n+1})$,
namely $(i_1, \cdots, i_n, s)$ with
\[
  \lvert a_{(i_0, \cdots, i_{n+1})} \rvert \lvert c \rvert^{(i_0, \cdots, i_{n+1})} > \varepsilon' =
  \varepsilon \lvert c_n \rvert^s.
\]
Since this is independent of the choice of $s$, we have shown as wanted.

For 2. we want to show
\[
  \lim \lvert g_s \rvert \lvert c_n \rvert^s = 0
\]
in the cofinite filter. Recall by the definition of the gauss norm,
\[
  \lvert g_s \rvert = \max \lvert a_{(i_0, \cdots, i_n, s)} \rvert \prod_{j=1}^{n} \lvert c_j \rvert ^{i_j}
\]
and moreover, since $g_s$ is a restricted power series the gauss norm is attained. That is
there exists $(k_{s_1}, \cdots k_{s_n})$ such that
\[
  \lvert g_s \rvert = \lvert a_{(k_{s_1}, \cdots k_{s_n}, s)} \rvert \prod_{j=1}^{n} \lvert c_j \rvert ^{i_j}.
\]
Now suppose $\phi (f)$ is not restricted, that is there exists an $\varepsilon$ such that there are
infinitely many $s$ with
\[
  \left( \lvert a_{((k_{s_1}, \cdots k_{s_n}), s)} \rvert \prod_{j=1}^{n} \lvert c_j \rvert ^{i_j} \right)
  \lvert c_n \rvert^s > \varepsilon.
\]
But once again this implies there are infinitely many $(i_1, \cdots, i_{n+1})$ such that
\[
  \lvert a_{(i_0, \cdots, i_{n + 1})} \rvert \prod_{j=1}^{n} \lvert c_j \rvert ^{i_j} > \varepsilon,
\]
contradicting $f$ being restricted.

For the other direction consider $p = \sum_{s = 0}^\infty q_s x_n^s \in
(R_c \langle x_1, \dots, x_n \rangle) \langle x_{n+1} \rangle$, with
$q_s = \sum a_{(i_{s_1}, \cdots, i_{s_n})} x_1^{i_{s_1}} \cdots x_n^{i_{s_n}}.$ Then
\[
  \phi^{-1} (p) = \sum b_{(i_1, \cdots, i_{n+1})} x_1^{i_1} \cdots x_n^{i_{n+1}},
\]
where $b_{(i_1, \cdots, i_n, s)} = a_{(i_{s_1}, \cdots, i_{s_n})}$.
Suppose $\varphi(p)$ is not restricted. Then there exists an $\varepsilon$ such that there are
infinitely many tuples $(i_1, \cdots, i_{n+1})$ with
\[
  \lvert b_{(i_1, \cdots, i_{n+1})} \rvert \prod_{j=1}^{n+1} \lvert c_j \rvert ^{i_j} > \varepsilon.
\]

Suppose in this infinite set that $i_{n+1}$ only takes finitely many values. Then for some $s$,
there are infinitely many $(i_{s_1}, \cdots, i_{s_n})$ such that $(i_{s_1}, \cdots, i_{s_n}, s)$
belongs in the set. However, this means there are infinitely many $(i_{s_1}, \cdots, i_{s_n})$ such that
\[
  \left( \lvert  b_{(i_{s_1}, \cdots, i_{s_n}, s)} \rvert \prod_{j=1}^{n} \lvert c_j \rvert ^{i_{s_j}} \right) \lvert c_n \rvert^s =
  \left( \lvert  a_{(i_{s_1}, \cdots, i_{s_n})} \rvert \prod_{j=1}^{n} \lvert c_j \rvert ^{i_{s_j}} \right) \lvert c_n \rvert^s > \varepsilon.
\]
Thus, setting $\varepsilon' = \varepsilon / \lvert c_n \rvert^s$ we have infinitely many
$(i_1, \cdots, i_n)$ such that the following inequality holds
\[
  \lvert  a_{(i_{s_1}, \cdots, i_{s_n})} \rvert \prod_{j=1}^{n} \lvert c_j \rvert ^{i_j} > \varepsilon'
\]
which contradicts the restrictedness of $q_s$.

On the other hand suppose $i_{n+1}$ takes infinitely many values, that is there are infinitely many
$s$ such that
\[
  \left( \lvert  b_{(i_{s_1}, \cdots, i_{s_n}, s)} \rvert \prod_{j=1}^{n} \lvert c_j \rvert ^{i_{s_j}} \right) \lvert c_n \rvert^s
  > \varepsilon.
\]
But after relabeling
\[
  \left( \lvert  b_{(i_{s_1}, \cdots, i_{s_n}, s)} \rvert \prod_{j=1}^{n} \lvert c_j \rvert ^{i_{s_j}} \right) \lvert c_n \rvert^s  =
  \lvert  a_{(k_1, \cdots, k_n)} \rvert \prod_{j=1}^{n+1} \lvert c_j \rvert ^{k_j}.
\]
Thus, for the infinitely many $s$ in the set
\begin{align*}
  \lvert g_s \rvert \lvert c_n \rvert^s & = \max \left\{ \lvert  a_{(i_{s_1}, \cdots, i_{s_n})} \rvert
  \prod_{j=1}^{n+1} \lvert c_j \rvert ^{k_j} \right\} \lvert c_n \rvert ^s \\
  & \geq \lvert  a_{(k_1, \cdots, k_n)} \rvert \prod_{j=1}^{n+1} \lvert c_j \rvert ^{k_j} \\
  & > \varepsilon
\end{align*}
contradicting the restrictedness of $p$.
\end{proof}


%%FILE blueprint/src/chapter/ToRewrite/TopGroup.tex
\chapter{Topological nilpotence and Power bounded elements}

\section{Topological nilpotence}


%%FILE blueprint/src/chapter/ToRewrite/WP.tex
\chapter{Weierstrass preparation}

%% We are going to reformat some of this file by generalising the topological nilpotence and
%% Power bounded defs to as in Wedhorn (see PowerBounded file)
%% The stuff on units etc can then be proven in generality (I think)

When computing zeroes of a power series it would be convenient if we could instead reduce it to
calculating zeroes of a polynomial. This turns out to be possible for restricted power series using
the $p$-adic Weierstrass preparation. This section is dedicated to proving this classical result. We
will follow \ref{BGR}.

\newcommand{\pb}[1]{\overset{\circ}{#1}}
\newcommand{\tn}[1]{\check{#1}}
\newcommand{\res}[1]{\widetilde{#1}}
\newcommand{\quot}[2]{{\raisebox{.2em}{$#1$}\left/\raisebox{-.2em}{$#2$}\right.}}

\section{Notation}

Before we state the theorem we first need to set some notation. Let $R$ be a semi-normed ring.

\begin{definition}
  We say an element $r$ is \emph{topologically nilpotent} if
    \[
     \lim r^n = 0.
    \]
    Denote the set of all topologically nilpotent elements of $R$ as $\tn{R}$.
\end{definition}

Unfortunately, the set $\tn{R}$ is not, in general, an ideal of $R$. However, we can define a
subring of $R$, that has it an ideal.

\begin{definition}
  We say an element $r$ is \emph{power-bounded} if the set
  \[
  \left\{ \lvert r^n \rvert : n \in \N \right\}
  \]
  is bounded.
  Denote the set of all power-bounded elements of $R$ as $\pb{R}$.
\end{definition}

\begin{lemma}
  $\pb{R}$ is a clopen subring of $R$, and $\tn{R}$ is an ideal in $\pb{R}$.
\end{lemma}

\begin{proof}
  Let $a, b \in \pb{R}$. Let $M$ be an upper bound for $\lvert a^n \rvert$ and $\lvert b^n \rvert$
  for all $n$. Then
  \[
  \lvert (a b )^n \rvert \leq \lvert a^n \rvert \lvert b^n \rvert \leq M^2,
  \]
  similarly we can use triangle inequality for closure of addition.
  For negation since, $R$ is a seminormed ring
  \[
  \lvert (- a)^n \rvert = \lvert a^n \rvert \leq M
  \]
  and thus since $1 \in \pb{R}$, we see it is a subring.

  If $a \in \nt{R}$ and $b \in \pb{R}$, then
  \[
  \lvert (a b)^n \rvert \leq \lvert a^n \rvert \lvert b^n \rvert \leq \lvert a^n \rvert M \to 0
  \]
  that is $ a b \in \tn{R}$, and we have shown $\tn{R}$ is an ideal.

  Finally, since $\pb{R}$ is a subgroup and contains $\tn{R}$, which is open in $R$, we have $\pb{R}$
  is clopen in $R$.
\end{proof}

Given these we can then define their quotient, $\res{R} = \pb{R} / \tn{R}$.

\section{Weierstrass division}

Weierstrass preparation will follow as an application of Weierstrass division, so we first prove that.
To do this we need to introduce the following terminology.


\begin{definition}
  Let $s$ be a natural number, we say a strictly convergent power series
  $g \in R \langle x_1 \dots, x_n \rangle >$ is $x_n$-distinguished
  of degree $s$ if, when writing $g = \sum_{v = 0}^\infty g_v(x_1, \dots, x_{n-1}) x_n^v$:
  \begin{itemize}
    \item $g_s$ is a unit; and
    \item $ \lvert g_s \rvert = \lvert g \rvert$ and $\lvert g_s \rvert > \lvert g_v \rvert \; \forall
      v > s$.
  \end{itemize}
\end{definition}

That is, $g_s$ achieves the gauss norm of $g$ and is the last index to do so.

\begin{theorem}
  Let $s$ be a natural number and $g \in T_n$ be $x_n$-distinguished of degree $s$. Then for each
  $f \in T_n$, there exists a unique $ q \in T_n$ and a unique $r \in T_{n-1}[x_n]$ with degree
  $r < s$ such that
  \[
    f = q g + r.
  \]
  Moreover, we have the following estimate
  \[
    \lvert f \rvert = \max \left\{ \lvert q \rvert \lvert g \rvert, \lvert r \rvert \right\},
  \]
  that is
  \[
    \lvert q \rvert \leq \lvert g \rvert^{-1} \quad \text{and} \quad \lvert r \rvert \leq \lvert f \rvert.
  \]
\end{theorem}

We break the proof up into the following steps.

\begin{definition}
  We say a subgroup $H$ of $G$ is $\epsilon$-dense if there is a real number $\epsilon < 1$ such that
  for each $g \in G$ there exists an $h \in H$ with
  \[
    \lvert g + h \rvert \leq \epsilon \lvert g \rvert.
  \]
\end{definition}

\begin{lemma}
  Let $(G, \lvert \cdot \rvert)$ be a normed group and $H$ a subgroup of $G$ $\epsilon$-dense in $G$.
  Then $H$ is dense in $G$.
\end{lemma}

To prove this we will use the following definition.

\begin{definition}
  Let $G$ be a seminormed group and let $H$ be any subgroup of $G$. For each $g \in G$, the nonnegative
  real number
  \[
  \lvert g, H \rvert := \inf_{h \in H} \lvert g + h \rvert
  \]
  is called the distance form $g$ to $H$.
\end{definition}

Using this definition it is clear to see the following:
\begin{itemize}
  \item $\lvert a + h, H \rvert = \lvert a, H \rvert$ for all $h \in H$;
  \item $\lvert g, H \rvert = \inf_{h \in H} d(a,h)$, where $d$ is the normal distance in a seminormed
    group; and
  \item $\overline{H} = \left\{ g \in G : \lvert g, H \rvert = 0 \right\}$.
\end{itemize}

\begin{proof}
  By the final bullet point above it suffices to show that $\lvert g, H \rvert = 0$ for all $g \in G$
  - as then $\overline{H} = G$, and we see $H$ is dense.

  Therefore, assume there is $g \in G$ such that $\lvert g, H \rvert > 0.$ Choose $h_1 \in H$ such that
  \[
    \lvert g + h_1 \rvert < \epsilon^{-1} \lvert g, H \rvert,
  \]
  this is possible since $\epsilon^{-1} > 1.$
  By $\epsilon$-dense, there exists $h_2$ such that
  \[
  \lvert (g + h_1) + h_2 \rvert \leq \epsilon \lvert g + h_1 \rvert.
  \]
  However, $h_1 + h_2 \in H$, thus the inequality obtained by combining the previous two,
  \[
    \lvert g + (h_1 + h_2) \rvert < \lvert g, H \rvert,
  \]
  is impossible. Therefore, $\lvert g, H \rvert = 0$ for all $g \in G$.
\end{proof}

With this we can then start proving parts of Weierstrass division.

\begin{lemma}\label{WD_bounds}
  Assume we are in the case of Weierstrass division and furthermore $\lvert g \rvert = 1$. Then, if
  there exists a representation
  \[
    f = q g + r
  \]
  we have the estimates
  \[
    \lvert q \rvert \leq \lvert f \rvert \quad \text{and} \quad \lvert r \rvert \leq \lvert f \rvert.
  \]
\end{lemma}

\begin{proof}
  Multiplying $f = q g + r$ by a scalar, we can assume
  \[
    \max \left\{ \lvert q \rvert, \lvert r \rvert \right\} = 1
  \]
  which implies $\lvert f \rvert \leq 1$ by the nonarchimedean property and $f = q g + r$.
  Thus, it suffices to show that $\lvert f = 1 \rvert.$ Assume the contrary, that is $\lvert f \rvert < 1$
  and so reduction into the residue field gives
  \[
    0 = \tilde{f} = \tilde{q} \tilde{g} + \tilde{r}.
  \]
  Since
  \[
    \deg \tilde{g} = s > \deg r > \deg \tilde{r}
  \]
  this would imply $\tilde{q} = \tilde{r} = 0$, which is a contradiction to the assumption that
  $\max {\lvert q \rvert, \lvert r \rvert} = 1$.
\end{proof}

\begin{lemma}\label{WD_uniqueness}
  Suppose we have a valid representation for Weierstrass division, then it is unique.
\end{lemma}

\begin{proof}
  Consider two representations
  \[
    f = q g + r \quad \text{and} \quad f = q' g + r',
  \]
  then we have
  \[
    0 = (q - q') g + (r - r').
  \]
  Applying the bounds to this new equation gives:
  \[
    \lvert q - q' \rvert \leq \lvert 0 \rvert = 0
  \]
  and similar for $r - r'$. Thus, $q - q' = 0 = r - r'$ and we see the two representations are equal.
\end{proof}


\begin{lemma}\label{WD_existence}
  There exists a representation
  \[
    f = q g + r
  \]
  in Weierstrass division.
\end{lemma}


%% I think to make sense of the epimorphism I will need prop 2 page 193
%% that proves \res{T}_n = \res{R}[n]... and specifically
%% \pb{T}_n = T_n^{\circ} ( = valuation ring )
%% \np{T}_n = T_n^\vee
%% \res{T}_n = T_n^\tilde

\begin{proof}
  Assume $g = \lvert 1 \rvert$ and consider the set
  \[
    B = \left\{ qg + r : r \in T_{n-1}[x_n], \; \deg r > s, \; q \in T_n \right\},
  \]
  this is a closed subgroup of $T_n$ and it suffices to show $B = T_n$.
  Writing $g = \sum g_s (x_1, \cdots, x_{n-1}) x_n^s$, define
  \[
    \varepsilon = \max_{s < t} \lvert g_t \rvert,
  \]
  since $g$ is $x_n$-distinguished of degree $s$ we must have $\varepsilon < 1$.

  Consider the natural ring epimorphism
  \[
    \tau_\varepsilon : \pb{T}_n \to \res{R}_\varepsilon [x_1, \cdots, x_n],
  \]
  where $R_\epsilon = \left\{ r \in R : \lvert r \rvert \leq \varepsilon \right\}$ and $\res{R}_\varepsilon
  = \pb{R} / R_\varepsilon$, with kernel
  \[
    \ker \tau_\varepsilon = \left\{ f \in T_n : \lvert f \rvert \leq \varepsilon \right\}.
  \]
  $\tau_\varepsilon (g)$ is a unitary polynomial in $x_n$ of degree $s$, and therefore we can apply
  Euclidean division (in $(\res{R}_\varepsilon [x_1, \cdots, x_{n-1}]) [x_n]$) on it.
  %% I need to work out why this is the case...?
  This means for all $f \in \pb{T}_n$ we can find $q \in \pb{T}_n$ and $r \in \pb{T}_{n-1}[x]$ with
  degree $ < s$ such that
  \[
    \tau_\varepsilon (f) = \tau_\varepsilon (q) \tau_\varepsilon (g) + \tau_\varepsilon (r)
  \]
  or equivalently,
  \[
    \lvert f - (qg + r) \rvert \leq \varepsilon
  \]
  since this will be in the kernel of $\tau (\varepsilon)$.
  Hence, for all $f \in T_n$, there exists $b \in B$ such that
  \[
    \lvert f - b \rvert \leq \varepsilon \lvert f \rvert,
  \]
  meaning $B$ is $\varepsilon$ dense in $T_n$ and therefore dense.
  This completes the proof, as $B$ closed and dense implies $B = T_n$.
\end{proof}

\section{Weierstrass preparation}

Earlier we claimed that Weierstrass preparation follows from Weierstrass division. We now show that
this is the case.

\begin{lemma}
  A series $f \in T_n$ with $\lvert f \rvert = 1$ is a unit in $T_n$ if and only if $\lvert f (0) \rvert = 1$
  and $\lvert f - f(0) \rvert < 1$. Thus, $f$ is a unit if and only if $\tilde{f}$ is a unit;
  i.e. a constant since $\res{T}_{n} \isom \tilde{R}[x_1, \cdots, x_{n}]$.
\end{lemma}

\begin{proof}
  Since $\lvert \cdot \rvert$ is a valuation, $f$ is a unit if and only if it is a unit in $\pb{T}$.
  This is equivalent to $a_{(0, \cdots, 0)}$ being a unit in $\pb{R}$ and $a_{(i_1,\cdots, i_n)} \in \tn{R}$
  for all other indices. This is the assertion.
\end{proof}

%% This depends on Proposition 1.4.2 3; trying to work out the proofs for this right now.

\begin{theorem}
  Let $g \in T_n$ be $x_n$-distinguished of degree $s$, for some natural number $s$. Then there is a
  unique monic polynomial $w \in T_{n-1}[x_n]$ of degree $s$ and a unique unit $e \in T_n$ such that
  \[
    g = e w
  \]
  with $\lvert w \rvert = 1$, so that $w$ is $x_n$-distinguished of degree $s$.
\end{theorem}

\begin{proof}
  By Weierstrass division, there exists unique $e' \in T_n$ and $r' \in T_{n-1} [x_n]$ with degree $r' < s$
  such that
  \[
    x_n^s = e' g + r.
  \]
  Then, defining
  \[
    w = x^s_n - r
  \]
  we see $w$ is monic of degree $s$ and $w = e' g$.

  It thus suffices to show that $e'$ is a unit, because then $g = w e^{-1}$ with the wanted properties.

  Using the estimates in Weierstrass division we have
  \[
    \lvert r \rvert \leq \lvert x_n^s \rvert = 1,
  \]
  and so we see that $\lvert w \rvert = 1$ and $x_n$ is distinguished of degree $s$.
  Now assume $\lvert g \rvert = 1$. From $\tilde{w} = \tilde{e'} \tilde{g}$, we conclude that $\tilde{e'}$
  is a unit in $\res{T}_{n-1}$, and since $\res{T}_{n-1} \isom \tilde{R}[x_1, \cdots, x_{n-1}]$ this means
  $\tilde{e'}$ is a constant.
  A fortiori this means it is a unit in $\res{T_n}$ and so invoking the above lemma it is a unit in $T_n$
  and we are done with the existence part.

  For uniqueness, suppose $g = e w$, define $r = x^s_n - w$. Then
  \[
    x_n^s = e^{-1} g + r
  \]
  is a valid representation satisfying the conditions for Weierstrass division and therefore $e$ and
  $w$ are unique, and moreover so is $w$.
\end{proof}


\section{Aside lemmas to fit in}

\begin{lemma}
  Let $A$ be a complete topological ring, then each element of the form $e = 1 - y$, where $y$ is topologically
  nilpotent, is a unit in $A$.
  We have $ e{-1} = \sum_{n=0}^\infty = 1 + z$, where $z$ is topologically nilpotent.
\end{lemma}

\begin{proof}
  The ring being complete we can set
  \[
    v = 1 + z, \quad \text{where} \quad z = \sum_{n=1}^\infty y^n.
  \]
  Since topological nilpotents are closed and $y^n$ is topologically nilpotent for all $n \geq 1$, we
  get $z$ is topologically nilpotent. We have
  \[
    e v = (1 - y) \sum_{n = 0}^{\infty} y^n = 1
  \]
  i.e. $v = e^{-1}$.
\end{proof}


\begin{lemma}
  Let $A$ be a complete topological space. Then a powerbounded element $a$ is a unit if and only if
  its residue class $\tilde{a} \in \tilde{A}$ is a unit.
\end{lemma}

\begin{proof}
  If $a$ is a uni then it is immediate that $\tilde{a}$ is a unit (take its inverse and reduce to
  its residue class, then they are inverses in $\tilde{A}$).

  To show the converse, consider a power bounded $b$ with $\tilde{a} \tilde{b} = \tilde{1}$. This
  means $a b = 1 - x$ for some topologically nilpotent $x$. By the previous lemma we conclude that
  $ab$ is a unit and hence $a$ is a unit.
\end{proof}

\begin{lemma}
  If $R$ is a complete normed space and $c$ is a positive real number. Then $R_c\lvert x \rvert$
  (restricted power series for parameter $c$) is complete with respect to the gauss norm.
\end{lemma}

\begin{proof}
  Let $(f_i) = (s\sum_{n = 0}^\infty a_{n,i} x^n)_i$ be a Cauchy sequence in $R\lvert x \rvert$.
  Since
  \[
  \lvert a_{n,i+1} - a_{n,i} \rvert * c^n \leq \lvert f_{i+1} - f_i \rvert
  \]
  (by writing out the definition of gauss norm of $f_{i+1} - f_i$) for fixed $n$, we see that each
  sequence $(a_{n,i})_i$ is a Cauchy sequence in $R$.

  Let $a_n$ be a limit of this sequence and set
  \[
    f = \sum_{n = 0}^\infty a_n x^n.
  \]
  We want to show that $f$ is a restricted power series for parameter $c$ and that
  \[
    \lim_{i \to \infty} \lvert f - f_i \rvert = 0.
  \]
  We may assume $\lvert f_j - f_i \rvert \leq 1 / i$ for all $j \geq i$.
  From all $\lvert a_{n,j} - a_{n,i} \rvert * c^n \leq 1/i$ we deduce by continuity of the norm on $R$
  that $\lvert a_n - a_{n,i} \rvert * c^n \leq 1 / i$ for all $n \geq 0$ and all $i$. For $n$ big enough
  we have $\lvert a_{n,i} \rvert *c^n < 1 /i$ since $f_i$ is restricted.
  Therefore, $\lvert a_n \rvert * c^n \leq 1/i$ and we get $f$ is restricted.

  Furthermore, we have
  \[
    \lvert f - f_i \rvert = \max \lvert a_n - a_{n,i} \rvert c^n \leq 1/i
  \]
  and, so we see $f$ is the limit of $f_i$.
\end{proof}

\begin{lemma}
  If $R$ is complete so is $\pb{R}$ the ring of power bounded elements.
\end{lemma}

\begin{proof}
  $\pb{R}$ is closed in a complete space.
\end{proof}

\begin{lemma}
  We have $\pb{R} \langle x \rangle = \pb{R \langle x \rangle}$ and
  $\tn{R} \langle x \rangle = \tn{R \langle x \rangle}.$
\end{lemma}

\begin{proof}
  %% This needs to be copied out... concern is that this may need to be an isometry...
  %% which is not hard but just more work
  %% perhaps I can see if the proof really needs it
  %% Otherwise perhaps I just want this as a statement saying that something in the LHS
  %% is powerbounded or topologically nilpotent respectively in the whole space
  %% which seems like it would be better
  %% I can then also have this as an iff if I want
\end{proof}


\begin{lemma}
  Let $R$ be complete. An element $f = \sum_{n = 0}^\infty a_n x^n \in \pb{R} \langle x \rangle$ (that
  is every coefficient is power bounded) is a unit in $\pb{R} \langle x \rangle$ if and only if $a_0$
  is a unit and $a_n$ is topologically nilpotent for all $n > 0$.
\end{lemma}

\begin{proof}
  Since $\pb{R}$ is complete we have $\pb{A} \langle x \rangle$ is also complete. Then by an earlier
  lemma we have that $f$ is a unit in $\pb{R} \langle x \rangle$ if and only if its residue class
  $ \tilde{f} = \sum_{n=0}^\infty \tilde{a}_n x^n$ is a unit.
  Moreover, we have that $a_0$ is a unit in $\pb{R}$ if and only if $\tilde{a_0}$ is a unit.

  Therefore, it is enough to show that any unit in the residue class is a unit in $\tilde{R}$, that
  is its a constant power series.
  %% This statement may have some nuance...
  %% For example, we want the reduction \tilde{R \langle x \rangle } to be \tilde{R} \langle x \rangle
  %% This may be done seperately or see if it can be avoided.
\end{proof}


%%FILE blueprint/src/chapter/ToRewrite/slopes.tex
\chapter{Slopes and eigenvalues}

In this section we combine the work on compact operators and Newton polygons.
Recall:
\begin{itemize}
  \item slopes of the Newton polygon gives information on the valuation of roots of the
    power series; %reference
  \item compact operators can be associated with a Fredholm determinant, that is a characteristic
    power series.
\end{itemize}

Therefore, since roots of the characteristic power series % is this still obvious with this definition
correspond to eigenvalues of the operator, we can construct the Newton polygon of the characteristic
power series to get information on the valuations of the eigenvalues.

This can be summed up in the following statement.

\begin{lemma}
  Let $u$ be a compact operator on a Banach space. Then the slopes of the Newton polygon of $u$,
  that is the slopes of the characteristic power series of $u$, are precisely the $p$-adic valuations
  of the eigenvalues of $u$.
\end{lemma}
%% need to decide if $p$-adic is the correct choice - or if I can work in some more generality


%%FILE blueprint/src/chapter/WeierstrassPrep.tex
\chapter{Weierstrass Division and Preparation}
\label{chap:wp}

Over $\mathbb{C}$, a holomorphic function on a disc with finitely many zeros factors as a
polynomial with those zeros times a nowhere-vanishing function.  The nonarchimedean analogue is
Weierstrass preparation, and it is sharper: the factorisation is algebraic, it is unique, and the
polynomial factor has degree exactly the number of zeros counted with multiplicity.

The classical statement is for the Tate algebra $K \langle T \rangle$ over a complete
nonarchimedean field, and the hypothesis on the divisor is phrased through the residue field:
$g$ is \emph{distinguished of degree $s$} when $\lvert g_s \rvert = \lVert g \rVert$ and
$\lvert g_i \rvert < \lVert g \rVert$ for $i > s$.  Two features of that setting are inconvenient
for us.  First, we shall want division over the ring $R_c \langle T \rangle$ at radii $c \neq 1$;
second, and more seriously, in the multivariate induction the coefficient ring is itself a Tate
algebra, which is a normed \emph{ring} and not a field.

The right generality is due to Martin \cite{Martin}, whose \S1.3 replaces ``unit'' by
``multiplicative unit'' in the hypothesis and thereby removes every restriction on the
coefficient ring and on the radius.  It is his form that we formalise, and the classical
field-theoretic statement is recovered as a specialisation.

\section{Distinguished series}

Throughout $A$ is a commutative ring with a complete ultrametric norm, $c > 0$ is a radius, and
$A_c \langle T \rangle$ denotes the restricted power series of parameter $c$ over $A$, carrying
the Gauss norm $\lVert f \rVert_c = \sup_i \lvert f_i \rvert c^i$.  We call $\lvert f_i \rvert c^i$
the $i$-th \emph{Gauss term} of $f$.

Recall that $a \in A$ is a \emph{multiplicative unit} if it is a unit and $\lvert a x\rvert =
\lvert a \rvert \lvert x \rvert$ for every $x$; over a field every nonzero element is one, and
this is the only place where the field hypothesis was ever used.

\begin{definition}
  \label{def:ismuldistinguished}
  \uses{def:gaussnorm, prop:restricted-ring, def:achievesgaussnorm}
  \lean{PowerSeries.IsMulDistinguished}
  \leanok
  Let $f \in A_c\langle T \rangle$ and $s \in \N$.  We say $f$ is \emph{distinguished of order
  $s$} if
  \begin{enumerate}
    \item $f_s$ is a multiplicative unit of $A$;
    \item the $s$-th Gauss term attains the Gauss norm, $\lvert f_s \rvert c^s = \lVert f
      \rVert_c$; and
    \item it strictly dominates every \emph{later} Gauss term: $\lvert f_i \rvert c^i < \lvert
      f_s \rvert c^s$ for all $i > s$.
  \end{enumerate}
\end{definition}

Condition (iii) says nothing about the coefficients below $s$; this asymmetry is exactly what
makes the definition usable at every radius, since only the tail has to be contracted.  Over a
field, or more generally when every unit of $A$ is a multiplicative unit, this coincides with the
classical notion.

The whole theory rests on a single computation, which says that multiplying by a distinguished
series is norm-multiplicative and tracks where the norm of the product is attained.  For $q \neq
0$ write $k_0(q)$ for the greatest index at which $q$ achieves its Gauss norm; this exists
because the Gauss terms of a restricted series tend to zero.

\begin{lemma}[{\cite[Lemma 1.26]{Martin}}]
  \label{lem:martin126}
  \uses{def:ismuldistinguished}
  \lean{PowerSeries.Restricted.norm_mul_of_isMulDistinguished}
  \leanok
  Let $g$ be distinguished of order $s$ and let $q \in A_c\langle T\rangle$ be nonzero.  Then
  \[
    \lVert g q \rVert_c = \lVert g \rVert_c \lVert q \rVert_c,
  \]
  and the product achieves its Gauss norm at the index $s + k_0(q)$, and at no larger index.
\end{lemma}

\begin{proof}
  \leanok
  \uses{def:ismuldistinguished}
  The coefficient of $T^{s + k_0}$ in $gq$ is $g_s q_{k_0} + \sum_{i + j = s + k_0, \, i \neq s}
  g_i q_j$.  In each remaining term either $i > s$, so that the $i$-th Gauss term of $g$ is
  strictly dominated by the $s$-th, or $i < s$ and hence $j > k_0$, so that the $j$-th Gauss term
  of $q$ is strictly dominated by the $k_0$-th by maximality of $k_0$.  Either way the term is
  strictly smaller than $\lvert g_s q_{k_0}\rvert c^{s + k_0}$, which is $\lVert g \rVert_c
  \lVert q \rVert_c$ because $g_s$ is a multiplicative unit.  The ultrametric equality case gives
  the claim, and the same estimate applied at larger indices gives the maximality.
\end{proof}

\section{Division}

\begin{theorem}[Weierstrass division, {\cite[Proposition 1.27]{Martin}}]
  \label{thm:wdivision}
  \uses{lem:martin126, thm:restricted-complete}
  \lean{PowerSeries.Restricted.weierstrassDivision_exists_of_isMulDistinguished}
  \leanok
  Let $A$ be complete, $c > 0$, and let $g \in A_c\langle T\rangle$ be distinguished of order
  $s$.  Then every $f \in A_c\langle T \rangle$ can be written
  \[
    f = g q + r, \qquad q \in A_c\langle T\rangle, \quad r \in A[T], \quad \deg r < s,
  \]
  and moreover $\lVert f \rVert_c = \max ( \lVert g \rVert_c \lVert q \rVert_c , \lVert r
  \rVert_c )$.
\end{theorem}

\begin{theorem}[Uniqueness]
  \label{thm:wdivision-unique}
  \uses{thm:wdivision}
  \lean{PowerSeries.Restricted.weierstrassDivision_q_unique_of_isMulDistinguished,
        PowerSeries.Restricted.weierstrassDivision_r_unique_of_isMulDistinguished}
  \leanok
  The pair $(q, r)$ of Theorem~\ref{thm:wdivision} is unique.
\end{theorem}

\begin{proof}
  \leanok
  \uses{lem:martin126}
  Uniqueness is the norm identity: if $gq + r = gq' + r'$ then $g(q - q') = r' - r$ has degree
  $< s$, while by Lemma~\ref{lem:martin126} the left-hand side achieves its Gauss norm at an
  index $\geq s$ unless $q = q'$.
\end{proof}

The existence half is where the completeness is spent.  Martin's argument, which we follow, is
a contraction:

\begin{proof}[Proof of Theorem~\ref{thm:wdivision}]
  \leanok
  \uses{lem:martin126}
  Truncate $g$ to a polynomial $g'$ of degree $s$ agreeing with $g$ up to $T^s$; its leading
  coefficient $g_s$ is invertible, so Euclidean division in $A[T]$ divides any polynomial by
  $g'$.  Applying this to a polynomial approximation of $f$ produces a one-step approximate
  division whose error is controlled by the tail factor
  \[
    \theta = \sup_{i > s} \frac{\lvert g_i \rvert c^i}{\lvert g_s \rvert c^s} < 1,
  \]
  which is $< 1$ precisely by condition (iii) of Definition~\ref{def:ismuldistinguished}.
  Iterating contracts the error geometrically; the limit exists by completeness and the bounds of
  Lemma~\ref{lem:martin126} keep it inside the set of admissible pairs, which is closed.
\end{proof}

Note where the hypotheses were used: strict domination is needed \emph{above} the distinguished
degree only, and no hypothesis on $c$ ever enters.  This is the point of Martin's formulation.

\section{Preparation}

\begin{theorem}[Weierstrass preparation, {\cite[Corollary 1.28]{Martin}}]
  \label{thm:wprep}
  \uses{thm:wdivision, thm:wdivision-unique}
  \lean{PowerSeries.Restricted.weierstrassPreparation_exists_of_isMulDistinguished}
  \leanok
  Let $A$ be complete and let $g \in A_c \langle T \rangle$ be distinguished of order $s$.  Then
  \[
    g = e \, \omega
  \]
  with $\omega \in A[T]$ monic of degree $s$ and $e$ a multiplicative unit of $A_c\langle T
  \rangle$.  The factorisation is unique.
\end{theorem}

\begin{proof}
  \leanok
  \uses{thm:wdivision, lem:martin126}
  Divide $T^s$ by $g$: $T^s = g q + r$ with $\deg r < s$, and set $\omega := T^s - r = gq$.  It
  remains to see that $q$ is a multiplicative unit.  By Lemma~\ref{lem:martin126} the product
  $gq = \omega$ achieves its norm at $s + k_0(q)$, and $\omega$ is monic of degree $s$, forcing
  $k_0(q) = 0$.  Comparing coefficients of $T^s$ gives $1 = \sum_i g_i q_{s - i}$, which exhibits
  $g_s q_0$ as $1$ minus a strictly smaller quantity, hence a multiplicative unit; and then $q =
  q_0 (1 + \text{small})$ is a multiplicative unit of $A_c\langle T\rangle$.  Uniqueness follows
  from uniqueness of the division.
\end{proof}

\begin{corollary}
  \label{cor:wprep-field}
  \uses{thm:wprep}
  \lean{PowerSeries.Restricted.weierstrassPreparation_polynomial_of_isMulDistinguished}
  \leanok
  Over a complete nontrivially normed field every nonzero element is a multiplicative unit, so
  Definition~\ref{def:ismuldistinguished} is the classical distinguishedness condition and
  Theorems~\ref{thm:wdivision} and \ref{thm:wprep} are the classical Weierstrass division and
  preparation theorems in $K_c\langle T \rangle$, at every radius $c$.
\end{corollary}

\section{The multivariate theorems}

Martin closes \S1.3 by remarking that the several-variable statements follow by applying the
one-variable theorems over the Tate algebra in the remaining variables.  We make that remark
formal.

Let $c = (c_0, \dots, c_n)$ be a polyradius and split off the first variable: the splitting
isomorphism
\[
  R_c \langle X_0, \dots, X_n \rangle \;\cong\; \bigl( R_{(c_1, \dots, c_n)} \langle X_1, \dots,
  X_n \rangle \bigr)_{c_0} \langle X_0 \rangle
\]
identifies a multivariate restricted series with a one-variable series whose coefficients lie in
the tail Tate algebra.  The correct hypothesis is therefore distinguishedness in $X_0$ over that
ring.

\begin{definition}
  \label{def:ismuldistinguishedx0}
  \uses{def:ismuldistinguished, thm:finsuccequiv}
  \lean{MvPowerSeries.Restricted.IsMulDistinguishedX0}
  \leanok
  $f \in R_c\langle X_0, \dots, X_n\rangle$ is \emph{distinguished in $X_0$ of order $s$} if its
  image under the splitting isomorphism is distinguished of order $s$ in the sense of
  Definition~\ref{def:ismuldistinguished} over the coefficient ring $R_{(c_1, \dots, c_n)}\langle
  X_1, \dots, X_n\rangle$.
\end{definition}

\begin{theorem}
  \label{thm:mvweierstrass}
  \uses{def:ismuldistinguishedx0, thm:wdivision, thm:wprep, thm:finsuccequiv}
  \lean{MvPowerSeries.Restricted.weierstrassDivision_exists_of_isMulDistinguishedX0,
        MvPowerSeries.Restricted.weierstrassPreparation_exists_of_isMulDistinguishedX0}
  \leanok
  Let $R$ be a complete ultrametric normed commutative ring, $c$ any polyradius, and $g$
  distinguished in $X_0$ of order $s$.  Then division by $g$ with remainder of $X_0$-degree $< s$
  exists and is unique, and $g$ factors uniquely as $e \omega$ with $\omega$ monic of degree $s$
  in $X_0$ over the tail algebra and $e$ a multiplicative unit.
\end{theorem}

\begin{proof}
  \leanok
  \uses{thm:wdivision, thm:wprep, thm:finsuccequiv}
  The tail algebra is a complete ultrametric normed commutative ring, so
  Theorems~\ref{thm:wdivision} and \ref{thm:wprep} apply to it verbatim; transport the
  conclusions back along the splitting isomorphism.
\end{proof}

%% The `IsDistinguished`-flavoured field-facing corollaries, and the Gauss-extension route that
%% preceded this development, are in PhD/BirkovichWP/. They are kept as history: the Martin route
%% supersedes them and needs no hypothesis on the radius.


%%FILE blueprint/src/chapter/biblio.tex
\bibliographystyle{plain} % We choose the "plain" reference style
\bibliography{PhD}


%%FILE blueprint/src/content.tex
\input{chapter/New/intro}

\input{chapter/RestrictedPowerSeries}
\input{chapter/WeierstrassPrep}
\input{chapter/NewtonPolygons}
\input{chapter/CompactOperators}
\input{chapter/QMF}
\input{chapter/Slopes}
\input{chapter/Jacobs}
\input{chapter/RieszColeman}
\input{chapter/LWXHalo}
\input{chapter/LWXModel}
\input{chapter/LWXSlopes}

\input{chapter/biblio}

